\documentclass[11pt]{article}

\usepackage[utf8]{inputenc}
\usepackage[T1]{fontenc}
\usepackage{lmodern}
\usepackage{amsmath,amssymb,amsthm,mathtools}
\usepackage{stmaryrd}
\usepackage[margin=1in]{geometry}
\usepackage{enumitem}
\usepackage{xcolor}
\usepackage{aliascnt}
\usepackage{hyperref}
\hypersetup{
  colorlinks=true,
  linkcolor=blue!60!black,
  citecolor=blue!60!black,
  urlcolor=blue!60!black
}

% theorem environments
\newtheorem{theorem}{Theorem}
\newaliascnt{lemma}{theorem}
\newtheorem{lemma}[lemma]{Lemma}
\aliascntresetthe{lemma}
\newaliascnt{proposition}{theorem}
\newtheorem{proposition}[proposition]{Proposition}
\aliascntresetthe{proposition}
\newaliascnt{corollary}{theorem}
\newtheorem{corollary}[corollary]{Corollary}
\aliascntresetthe{corollary}
\newaliascnt{question}{theorem}
\newtheorem{question}[question]{Question}
\aliascntresetthe{question}
\theoremstyle{definition}
\newaliascnt{definition}{theorem}
\newtheorem{definition}[definition]{Definition}
\aliascntresetthe{definition}
\newaliascnt{example}{theorem}
\newtheorem{example}[example]{Example}
\aliascntresetthe{example}
\newaliascnt{remark}{theorem}
\newtheorem{remark}[remark]{Remark}
\aliascntresetthe{remark}

\usepackage[nameinlink,capitalize]{cleveref}
\crefname{theorem}{Theorem}{Theorems}
\crefname{lemma}{Lemma}{Lemmas}
\crefname{proposition}{Proposition}{Propositions}
\crefname{corollary}{Corollary}{Corollaries}
\crefname{question}{Question}{Questions}
\crefname{definition}{Definition}{Definitions}
\crefname{example}{Example}{Examples}
\crefname{remark}{Remark}{Remarks}
\Crefname{theorem}{Theorem}{Theorems}
\Crefname{lemma}{Lemma}{Lemmas}
\Crefname{proposition}{Proposition}{Propositions}
\Crefname{corollary}{Corollary}{Corollaries}
\Crefname{question}{Question}{Questions}
\Crefname{definition}{Definition}{Definitions}
\Crefname{example}{Example}{Examples}
\Crefname{remark}{Remark}{Remarks}

\sloppy
\setlength{\emergencystretch}{3em}

% macros
\newcommand{\R}{\mathbb{R}}
\newcommand{\N}{\mathbb{N}}
\newcommand{\E}{\mathbb{E}}
\newcommand{\Prob}{\mathbb{P}}
\newcommand{\A}{\mathcal{A}}
\newcommand{\F}{\mathcal{F}}
\newcommand{\K}{\mathcal{K}}
\newcommand{\T}{\mathcal{T}}
\newcommand{\Q}{\mathcal{Q}}
\newcommand{\Sset}{\mathcal{S}}
\newcommand{\U}{\mathcal{U}}
\newcommand{\X}{\mathcal{X}}
\newcommand{\C}{\mathcal{C}}
\newcommand{\Homp}{\operatorname{Hom}^{+}}
\newcommand{\Homalg}{\operatorname{Hom}}
\newcommand{\Emb}{\operatorname{Emb}}
\newcommand{\Ext}{\operatorname{Ext}}
\newcommand{\supp}{\operatorname{supp}}
\newcommand{\Forb}{\operatorname{Forb}}
\newcommand{\brak}[1]{\left\langle #1\right\rangle}
\newcommand{\avg}[2]{\left\llbracket #1 \right\rrbracket_{#2}}
\newcommand{\qmap}{\mathfrak q}
\newcommand{\eps}{\varepsilon}
\newcommand{\one}{\mathbf 1}
\newcommand{\vtype}{\mathsf{1}}
\DeclareMathOperator{\cl}{cl}

\title{Completeness of Forbidden-Subgraph Reasoning in Flag Algebras:\\
       A Root-Planting Criterion}
\author{}
\date{\today}

\begin{document}
\maketitle

\begin{abstract}
Razborov's flag algebras are among the most successful tools in extremal graph theory,
underlying many computer-assisted density bounds. Often one wants such a bound only inside a hereditary
class of graphs---those avoiding some forbidden subgraph, like the triangle-free or
$C_4$-free graphs.  There are two natural ways to make the method respect the restriction.
One can compute in the \emph{quotient} flag algebra of the class---the algebra the method
actually uses, in which the forbidden configurations have been set to zero.  Or one can
stay among all graphs and ask that the inequality hold for almost every local view of a
limit of graphs from the class: this is what we call \emph{ensemble} semantics.  The
quotient method is always sound, so a quotient proof is a genuine bound for the class; we
ask when it is also \emph{complete}, capturing every true restricted inequality.

Our main result (a support-closure criterion) is that the two semantics agree, for every
flag-algebra element, exactly when a geometric condition we call \emph{root-plantability}
holds: every labelled limit of the class must be realisable as a typical local view of an
unlabelled one.  This holds whenever the class is blow-up-closed---recovering the
clique-free classes, triangle-free among them---with substitution-closure as a stronger sufficient
condition.  It can also fail: the $C_4$-free graphs are a counterexample.  The common
cause is \emph{pinning}, a local quantity frozen across the class while the quotient
still sees it vary.

This incompleteness, however, never weakens an actual density bound: such a bound concerns
unlabelled graphs, where the two semantics always coincide, and the gap lives only in the
labelled intermediate steps a proof passes through---closing once one takes the natural
limit.  We illustrate the whole picture on the $C_5$-free graphs, where root-plantability
holds for a single root and for a non-adjacent pair of roots but fails for an edge, for
which a ``book'' graph obstructs it, refuting the natural all-types conjecture.  Finally,
ensemble semantics is strictly more flexible: it can impose non-hereditary constraints,
such as fixing a limit's edge density, and thereby read off the local structure of
extremal configurations.  For instance, the verified $K_4$-free $P_4$ certificate, read
through this lens, forces every extremiser to be the balanced complete tripartite graph,
with a quantitative stability version.
\end{abstract}

\section{Introduction}

A recurring task in extremal combinatorics is to prove a density inequality that holds
only under a forbidden-subgraph constraint---for instance, ``in every triangle-free graph,
this combination of subgraph densities is non-negative.''  Razborov's flag algebras
\cite{Razborov2007} are among the most successful tools for such problems, and they
underlie many computer-assisted density bounds.  This paper concerns a foundational question that the
constrained setting raises: \emph{how should a flag-algebra proof enforce the constraint,
and does the choice cost anything?}

There are two natural answers.  Let $\T_1$ be the class of graphs we care about---say the
triangle-free graphs, or more generally any \emph{hereditary} class---one closed under
\emph{induced} subgraphs: whenever a graph is in the class, so is every graph obtained by
deleting vertices and keeping all edges among the vertices that remain.  We always use this
induced-subgraph relation, never the weaker one that may also delete edges; the distinction
matters here.  Equivalently, $\T_1$ is cut out by forbidding, as induced subgraphs, the
members of some (possibly infinite) family---and a class named ``$H$-free'' is of this
kind, since having no $H$ as an ordinary subgraph is the same as omitting,
as induced subgraphs, the finitely many graphs on $|H|$ vertices that contain $H$.
One option is to pass to the \emph{quotient} flag algebra of $\T_1$ and ask for
non-negativity there; this is the algebra a flag-algebra computation actually works in,
the one in which the forbidden configurations have simply been set to zero.  The other
option is to stay in the ambient world $\T_0$ of \emph{all} graphs and test the inequality
along the constrained \emph{limits}: take a convergent sequence of graphs from $\T_1$,
single out (``label'') a few of its vertices at random, and require the inequality to hold
almost surely for the resulting local view.  We call these the \emph{quotient} and
\emph{ensemble} semantics.  The first is what one computes with; the second is the
property one really wants throughout the class.

One direction is automatic: quotient non-negativity always implies ensemble
non-negativity (\cref{thm:support-criterion}).  So forbidden-subgraph flag-algebra
reasoning is unconditionally \emph{sound}---a certificate proved in the quotient is a
genuine density bound for the constrained class, and its validity is never in question.
The substantive question is the converse, \emph{completeness}: can \emph{every} true
constrained inequality be certified in the quotient, or are there inequalities that hold
throughout the class yet are invisible to the algebra one computes in?

\paragraph{The criterion.}
Our main result identifies exactly when completeness holds.  Fix a type $\sigma$:
a small graph whose distinguished vertices, with their prescribed adjacencies, are
called the \emph{roots}.  The two semantics then live in one and the same space of
\emph{labelled limits}.  Such a limit is a \emph{view} from an embedding of $\sigma$:
choose an embedding $\sigma\to G$ in a large graph of the class, treat the images of the
roots as the roots, record the density of every \emph{flag} (a finite graph with a
specified embedding of $\sigma$ inside it), and pass to the limit---the local picture the
roots see.
Each semantics singles
out a subset of these views.  The set $Q_\sigma$ ($Q$ for \emph{quotient}) collects the
views the quotient can exhibit: those obtained by placing an embedding of $\sigma$
\emph{anywhere} in graphs of the class---even at \emph{exceptional} vertices, ones that a
uniformly random choice almost never lands on.  The set $S_\sigma$
($S$ for \emph{support}) collects the views ensemble semantics actually sees: those that
arise with positive probability when the embedding of $\sigma$ is chosen at a
\emph{random}---hence typical---spot in a large graph of the class.  Soundness is the
inclusion $S_\sigma\subseteq Q_\sigma$, which always holds.  We prove
(\cref{thm:support-criterion}) that the two semantics coincide for \emph{every}
flag-algebra element if and only if the reverse inclusion holds too, that is,
\[
  S_\sigma=Q_\sigma .
\]
We call this condition \emph{root-plantability}: every view the quotient can exhibit must
be \emph{typical}---a limit of views that occur with positive probability when the roots
are dropped at random into a large graph of the class.  Informally, any view the
quotient can display at some special, hand-picked vertex must also be realisable at a
random, hence typical, vertex of a large constrained graph; the roots can be
\emph{planted} at typical vertices.  When this fails, some inequality holds almost surely
under random rooting yet cannot be seen in the quotient, and the quotient calculus is
\emph{incomplete}.  (\Cref{def:root-planting} makes the set $S_\sigma$ precise.)

\paragraph{When root-planting works.}
Root-plantability holds whenever the constrained class is closed under
\emph{independent blow-ups}
(replacing each vertex by an independent set of clones): blowing one vertex up into a
\emph{clone class} of positive density---a positive fraction of all the vertices, so that a
random vertex lands in it with probability bounded away from zero---lets one realise at
such a typical vertex what the quotient could show only at an exceptional one.
This already covers the clique-free classes ($K_r$-free graphs, $r\ge3$), the
triangle-free graphs among them
(\cref{thm:clone-root-plantable,cor:clique-free}).  The same argument has a
``true-twin'' version using cliques in place of independent sets
(\cref{thm:true-clone-root-plantable}, which yields e.g.\ the cluster graphs).  Both are
instances of a single criterion: root-plantability holds for every \emph{blow-up-closed}
class---one in which a single vertex may always be blown up to an arbitrarily large graph,
of the class's own choosing, without leaving the class (\cref{thm:blowup-root-plantable}).
This is what the proof actually needs, because the blown-up interior is invisible to the
sampling argument; closure under \emph{substitution}---replacing each vertex by an entire
graph drawn from the class---is one strictly stronger sufficient condition
(\cref{thm:substitution-root-plantable}).  We also isolate a finite, checkable form of the
argument---the \emph{finite planting} criterion (\cref{thm:finite-local-planting})---and show
it can be met by a purely \emph{local} repair, changing only a sparse set of edges, even when
the blown-up graph does not automatically lie in the class and no global closure property
applies (\cref{thm:sparse-repair-planting}); this is what the $C_5$-free constructions below
use.

\paragraph{When it fails.}
Heredity alone is not enough.  Consider the $C_4$-free graphs.  Their large members are
sparse---every $n$-vertex $C_4$-free graph has only $o(n^2)$ edges---so almost every vertex
has vanishingly few neighbours.  The ensemble drops the single root at a \emph{random}
vertex and measures its \emph{one-root edge density}---the fraction of the other vertices
that root is joined to.  In a graph this sparse a random vertex has hardly any neighbours,
so for almost every choice of the root this fraction is essentially $0$; in the limit it
is \emph{almost surely} $0$ (that is, $0$ for all but a negligible set of root choices).
The quotient, by contrast, may put the root wherever it likes.  The class contains
arbitrarily large \emph{stars}---a single central vertex joined to many leaves and to
nothing else---and rooting at such a centre gives one-root edge density $1$, because the
centre is joined to \emph{every} other vertex.  That centre is an exceptional vertex: the
quotient can sit on it, while the ensemble, only ever landing on a typical vertex, never
does.  So the inequality ``the one-root edge density is $0$'' holds throughout the ensemble
yet fails in the quotient, and the class is not root-plantable
(\cref{thm:degenerate-obstruction,cor:c4-counterexample}).  The same failure befalls the
$K_{s,t}$-free, $C_{2k}$-free, and planar graphs (\cref{cor:degenerate-family}).  One
might guess this is a defect of \emph{sparsity}, but root-plantability is preserved by
complementation (\cref{lem:complementation}), so every sparse obstruction has a
\emph{dense} mirror image (\cref{cor:codegenerate}); neither sparsity nor density is the
dividing line.  The uniform cause is what we call \emph{pinning}: the density of
a fixed flag or algebra element is frozen to a single value across all constrained
limits while the quotient still attains another (\cref{thm:pinning}).  For
edge-deletion-closed hereditary classes, any such pinning is to a boundary value $0$ or $1$
(\cref{thm:no-interior}).

\paragraph{The gap does not affect density bounds.}
This incompleteness may sound alarming, but it never weakens an actual application.  A
density bound---a Mantel- or Tur\'an-type inequality---is a statement about
\emph{unlabelled} graphs, the \emph{empty type}, and there the two semantics provably
coincide (\cref{prop:empty-type}): no density bound is ensemble-true but quotient-false.
The incompleteness lives entirely in the intermediate, \emph{labelled} steps a proof
passes through.  Even there it is benign: the witnesses at non-empty types that expose the gap
unlabel to zero---they average away to nothing once the labels are dropped
(\cref{prop:ideal-zero})---and every contribution the ensemble relaxation can make is,
after unlabelling, a limit of ordinary \emph{sum-of-squares contributions}: the standard
non-negative building blocks of a quotient certificate, each obtained by squaring a
flag-algebra element and averaging it down to the unlabelled level
(\cref{thm:no-closed-certificate-gap}).  What could in principle remain is
only a difference at the level of exact, finite certificates---never a new asymptotic
bound.

\paragraph{The \texorpdfstring{$C_5$}{C5}-free graphs.}
These are the first dense, non-blow-up-closed class for which the question is genuinely
nontrivial, and they show that root-plantability must be asked \emph{one type at a time}.
At the one-vertex type, and at the two-vertex non-edge type (two non-adjacent roots), a local
repair of the blow-up construction works, so root-plantability holds
(\cref{thm:c5-one-root,thm:c5-nonedge-root}).  At the two-vertex \emph{edge} type (two
adjacent roots) it
fails: a typical edge has almost no common neighbours (triangles are too rare in
$C_5$-free graphs), while a sparse ``book'' graph has an edge whose endpoints share every
other vertex (\cref{thm:c5-edge-not-root-plantable}).  So the natural all-types conjecture
is false---though, consistent with the previous paragraph, even this genuine obstruction
leaves every $C_5$-free density bound untouched (\cref{cor:c5-edge-closed-inert}).

\paragraph{A bonus: constraints the quotient cannot express.}
Finally, ensemble semantics is strictly more flexible than the quotient, because it can
impose \emph{non-hereditary} limit constraints---conditions such as ``edge density exactly
$\tfrac12$'' that do not define a hereditary class and so have no quotient algebra at all.
Fixing such a constraint and reading off the local structure it forces turns a
flag-algebra certificate into information about extremal configurations.  On the Mantel
equality slice---the triangle-free limits that actually attain edge density
$\tfrac12$---this recovers the local statistics of the extremal complete bipartite
limit, for instance that a typical vertex has degree $\tfrac12$ (\cref{thm:relative-mantel}).
Applied to the verified $K_4$-free $P_4$ certificate, the same mechanism forces every
extremiser to be the balanced complete tripartite graph outright---the two mined
labelled equations are rigid by themselves, with no appeal to the equality case of
Tur\'an's theorem (\cref{thm:k4free-p4-tripartite})---with a quantitative stability
version (\cref{thm:k4free-p4-quant-stability}) and an
$r$-partite generalisation (\cref{thm:parametric-p4-equality-slice}).

\section{Flag-algebra background}\label{sec:background}

We recall the flag-algebra apparatus \cite{Razborov2007}, at the level of detail a reader
new to it needs.  Everything is phrased for finite simple graphs.  Throughout, $\T$ stands
for a hereditary class of finite simple graphs---either the class
$\T_0=\T_{\mathrm{Graph}}$ of all finite simple graphs, or a forbidden-subgraph subclass
$\T_1$---and ``graph in $\T$'' means a finite graph belonging to that class.

\paragraph{Types and flags.}
A \emph{type} $\sigma$ is a fixed graph in $\T$ on $k$ \emph{roots}---$k$
distinguished vertices together with the edges among them.  An \emph{embedding}
of $\sigma$ into a graph $H$ is an injective map from the roots of $\sigma$ to
$V(H)$ that preserves both edges
and non-edges: two vertices are adjacent in $\sigma$ if and only if their images
are adjacent in $H$.  Equivalently, it identifies $\sigma$ with an induced
subgraph of $H$, with the roots kept in their prescribed order.
A \emph{$\sigma$-flag} is a graph in $\T$ with a distinguished embedding of
$\sigma$ singled out inside it---informally, a finite graph in which $k$ roots
form the pattern $\sigma$.  The \emph{empty type}
($k=0$) marks nothing, so a flag of the empty type is just a graph in $\T$.
Flags are taken up to root-preserving isomorphism, and $\F^\sigma[\T]$
(abbreviated $\F^\sigma$) is the countable set of all such isomorphism classes
(so two flags related by a root-preserving isomorphism are the same element of
$\F^\sigma$).  For $\sigma$-flags $F$ and $G$ with $|F|\le|G|$, the
\emph{density} $p(F,G)\in[0,1]$ is the probability that a uniformly random set of
$|F|-k$ unlabelled vertices of $G$, taken together with the $k$ roots, induces a flag
isomorphic to $F$; for the empty type this is the ordinary induced-subgraph density.

\paragraph{The flag algebra.}
$\A^\sigma[\T]$ is the real vector space spanned by $\F^\sigma$ (formal real combinations
of flags), modulo the linear relations
\[
  F=\sum_{|H|=|F|+1}p(F,H)\,H
\]
relating each flag to the flags one vertex larger.  This quotient is compatible with the
induced-density evaluation $p(\,\cdot\,,G)$: for every larger $G$ the identity
$p(F,G)=\sum_{|H|=|F|+1}p(F,H)\,p(H,G)$ holds, so $p(\,\cdot\,,G)$ respects the defining
relations and descends to a well-defined linear functional on $\A^\sigma[\T]$.  We keep the
same notation for this extension: if $f\in\A^\sigma[\T]$ is represented by
$\sum_i a_iF_i$, then
\[
  p(f,G):=\sum_i a_i\,p(F_i,G),
\]
independently of the chosen representative.  The space
carries a commutative, associative \emph{product} that records how densities combine: the
product $F_1\cdot F_2$ of $\sigma$-flags on $n_1$ and $n_2$ vertices is the formal
combination of the $\sigma$-flags $H$ on $n_1+n_2-k$ vertices, each weighted by the
probability that splitting the unlabelled vertices of $H$ uniformly into parts of sizes
$n_1-k$ and $n_2-k$ exhibits $F_1$ on one part (with the roots) and $F_2$ on the other.
The product is rigged so that densities multiply in the limit (made precise next); its
unit is the flag $\one_\sigma$ consisting of the type alone.  We write $\A^0[\T]$ for the
empty-type algebra.

\paragraph{Positive homomorphisms as limits.}
A \emph{positive homomorphism} is a linear, unit-preserving, multiplicative map
$\phi:\A^\sigma[\T]\to\R$---so $\phi(\one_\sigma)=1$ and
$\phi(F_1\cdot F_2)=\phi(F_1)\,\phi(F_2)$---that is non-negative on flags, i.e.\
$\phi(F)\ge0$ for every $F\in\F^\sigma$.  The set of all of them is denoted
$\Homp(\A^\sigma[\T],\R)$.  By Razborov's theory these are exactly the limits of
convergent flag sequences: a sequence of finite $\sigma$-flags of growing size
\emph{converges} when every density $p(F,\cdot)$ converges, its limit is the homomorphism
$\phi$ with $\phi(F)=\lim p(F,\cdot)$, and conversely every positive homomorphism arises
this way \cite[Theorem 3.3]{Razborov2007}.  Terminologically, from now on a
\emph{limit} means such a positive homomorphism unless we explicitly mention an
approximating sequence.  Thus an \emph{unlabelled limit} is a point of
$\Homp(\A^0[\T],\R)$, while a \emph{$\sigma$-labelled limit} is a point of
$\Homp(\A^\sigma[\T],\R)$.  We use \emph{roots} for the distinguished vertices
of a type or flag, and reserve \emph{labelled} for the associated flag-algebra
objects, such as labelled flags, labelled algebras, and labelled limits.  We use
``limit'' and ``positive homomorphism'' interchangeably; $\phi(F)$ is the density of the flag $F$ in
the limit $\phi$.  On a single flag
this value lies in $[0,1]$: for each size the flags of that
size sum to the unit,
$\sum_{|H|=|F|}H=\one_\sigma$ (a vertex set induces exactly one flag), so
$0\le\phi(F)\le\sum_{|H|=|F|}\phi(H)=\phi(\one_\sigma)=1$.  Hence evaluation on all flags,
$\phi\mapsto(\phi(F))_{F\in\F^\sigma}$, embeds $\Homp(\A^\sigma[\T],\R)$ into the compact
metrizable cube $[0,1]^{\F^\sigma}$; the image is closed---cut out by the linear and
multiplicative homomorphism conditions---so $\Homp(\A^\sigma[\T],\R)$ is itself compact
metrizable.

The averaging (unlabelling) operator
\[
  \avg{\cdot}{\sigma}:\A^\sigma[\T]\to\A^0[\T]
\]
forgets the labels of a $\sigma$-flag and averages over the choice of labelling.  The
only properties we use are that it is linear and that on a single flag $F$ with
underlying unlabelled graph $M$ it returns a \emph{non-negative scalar multiple}
of $M$.  Recall that $\one_\sigma\in\A^\sigma[\T]$, the flag consisting of the type alone, is the
multiplicative unit of $\A^\sigma[\T]$.  Its unlabelling
\[
  \brak{\sigma}:=\avg{\one_\sigma}{\sigma}\in\A^0[\T]
\]
is the unlabelled density of embeddings of the type $\sigma$: if the unlabelled limit
$\phi_0$ is represented by finite graphs $G_n$, then $\phi_0(\brak\sigma)$ is the
limiting probability that $k$ vertices of $G_n$, chosen uniformly at random and placed in a
random order, induce the labelled graph $\sigma$.

\paragraph{Non-degenerate types.}
A type $\sigma$ is \emph{non-degenerate} for the class $\T$ if
$\brak\sigma\neq0$ in $\A^0[\T]$.  By the unlabelling property above,
$\brak\sigma=\avg{\one_\sigma}{\sigma}$ is a non-negative multiple of the unlabelled graph
underlying $\sigma$, so it is a \emph{non-negative element}: $\phi_0(\brak\sigma)\ge0$ for
every unlabelled limit $\phi_0$.  Since an element of $\A^0[\T]$ is the zero element exactly
when every unlabelled limit assigns it the value $0$, non-degeneracy is equivalently the
existence of an unlabelled limit $\phi_0$ with $\phi_0(\brak\sigma)>0$.  Intuitively, the
pattern $\sigma$ is actually realised, occurring with positive density in some unlabelled
limit of the class, so that an embedding of $\sigma$ can be sampled and labelled.  A
degenerate type occurs with density zero in every unlabelled limit, so the random rooting
below acts on nothing, and the root-planting story of this paper is vacuous for it.

If
$\phi_0\in\Homp(\A^0[\T],\R)$ and $\phi_0(\brak\sigma)>0$, Razborov's extension
theorem gives a unique probability measure on
$\Homp(\A^\sigma[\T],\R)$, denoted here by
\[
  \Ext_\sigma(\phi_0),
\]
characterised by
\begin{equation}\label{eq:extension-expectation}
  \E_{\phi^\sigma\sim\Ext_\sigma(\phi_0)}[\phi^\sigma(g)]
  =
  \frac{\phi_0(\avg{g}{\sigma})}{\phi_0(\brak\sigma)}
  \qquad(g\in\A^\sigma[\T]).
\end{equation}
The positivity assumption on $\phi_0(\brak\sigma)$ is exactly what permits this
conditioning: the denominator is the limiting probability of seeing the root type.
Intuitively, $\Ext_\sigma(\phi_0)$ is the law of the labelled limit seen from a uniformly
random embedding of $\sigma$ in finite graphs converging to $\phi_0$;
\eqref{eq:extension-expectation} records that the expected density of $g$ is the
unlabelled density of $g$, normalised by that of $\sigma$.  Moreover, if a finite sequence
converges to $\phi_0$, then the
corresponding finite random-rooting distributions converge weakly to $\Ext_\sigma(\phi_0)$
\cite[Theorems 3.5 and 3.12]{Razborov2007}.

Recall that the \emph{support} of a Borel probability measure $P$ on a compact metrizable
space $X$, written $\supp(P)$, is the smallest closed set of full $P$-measure; equivalently,
\[
  \supp(P)=\{x\in X:P(U)>0\text{ for every open }U\ni x\}.
\]
We will repeatedly use the following elementary consequence.

\begin{lemma}[Almost-sure non-negativity and support]\label{lem:support-as}
Let $P$ be a Borel probability measure on a compact metrizable space $X$, and let
$h:X\to\R$ be continuous.  Then
\[
  P[h(x)\ge0]=1
\]
if and only if $h(x)\ge0$ for every $x\in\supp(P)$.
\end{lemma}

\begin{proof}
($\Rightarrow$) Suppose $P[h(x)\ge0]=1$, equivalently $P[h(x)<0]=0$, and let
$x\in\supp(P)$.  If $h(x)<0$, then $\{y:h(y)<0\}$ is an open set (as $h$ is continuous)
containing $x$, so by the defining property of the support it has positive $P$-measure,
contradicting $P[h(x)<0]=0$.  Hence $h(x)\ge0$.

($\Leftarrow$) Suppose $h\ge0$ on $\supp(P)$.  Then
$\{y:h(y)\ge0\}\supseteq\supp(P)$, and the support has full measure, so
$P[h(x)\ge0]\ge P(\supp(P))=1$; that is, $P[h(x)\ge0]=1$.
\end{proof}

In applications $X$ will be a homomorphism space and $h(\chi)=\chi(f)$; the lemma lets us
replace almost-sure non-negativity under a random extension by pointwise non-negativity on
its support.

\section{Constrained classes and supported homomorphisms}

Let $\T_1$ be a hereditary subclass of $\T_0$.  For a non-degenerate type $\sigma$ of
$\T_1$, define
\[
  X_\sigma:=\Homp(\A^\sigma[\T_0],\R).
\]
Call a $\sigma$-flag \emph{forbidden} when its underlying unlabelled graph belongs to
$\T_0$ but not to $\T_1$.  Being forbidden depends only on
that underlying graph, not on how the flag is labelled.  The quotient map
\[
  \qmap_\sigma:\A^\sigma[\T_0]\twoheadrightarrow\A^\sigma[\T_1]
\]
is precisely the map that sends each forbidden flag to $0$ and every other flag to itself.
It pulls each $\psi\in\Homp(\A^\sigma[\T_1],\R)$ back to a point of $X_\sigma$, the composite
${\psi} \circ {\qmap_\sigma}$: this is an algebra homomorphism, and it is non-negative on every
flag because $\qmap_\sigma$ sends flags to flags or to $0$ and $\psi$ is non-negative on
flags.  We write
\[
  Q_\sigma
  :=
  \{\psi\circ\qmap_\sigma:
      \psi\in\Homp(\A^\sigma[\T_1],\R)\}
  \subseteq X_\sigma.
\]
Thus $Q_\sigma$ is the constrained, or $\T_1$-supported, part of the ambient homomorphism
space.  We call $Q_\sigma$ the space of constrained
$\sigma$-labelled limits, and $Q_0:=Q_\emptyset$ the space of constrained unlabelled
limits.  The set $Q_\sigma$ has a simple intrinsic description, one that never mentions the
quotient: a homomorphism $\chi\in X_\sigma$ lies in $Q_\sigma$ if and only if it vanishes
on every forbidden flag.  One direction is immediate.  Since $\qmap_\sigma$ sends each forbidden
flag to $0$, any $\chi=\psi\circ\qmap_\sigma$ satisfies
$\chi(F)=\psi(\qmap_\sigma(F))=\psi(0)=0$ on every forbidden flag $F$, so every member of
$Q_\sigma$ vanishes on the forbidden flags.

The converse---that a homomorphism vanishing on all forbidden flags is itself of the form
$\psi\circ\qmap_\sigma$---is the substantive direction, and it is where heredity does the
work.  The strategy is to \emph{factor} $\chi$ through the quotient map.  By the universal
property of a quotient, a homomorphism on $\A^\sigma[\T_0]$ has the form
$\psi\circ\qmap_\sigma$ precisely when it vanishes on the kernel of $\qmap_\sigma$, so the
first task is to identify that kernel.

Forming $\A^\sigma[\T_1]$ amounts to declaring the forbidden flags equal to $0$.  To do
this \emph{inside an algebra}, and be left with an algebra, one must quotient not by an
arbitrary subspace but by an \emph{ideal}: a subspace closed under multiplication by every
algebra element.  A priori the ideal generated by the forbidden flags could be larger than
their linear span, since it must also contain every product of a forbidden flag with an
arbitrary element, and such a product need not, on its face, be a combination of forbidden
flags.  Heredity is exactly what rules this out: the product of a forbidden flag with any
flag is again a combination of forbidden flags, because each flag occurring in the product
contains the forbidden underlying graph as an induced subgraph, and any graph that contains
a non-$\T_1$ graph as an induced subgraph is itself outside $\T_1$.  This last step uses
only that $\T_1$ is hereditary---equivalently, that it is cut out by forbidding a family of
induced subgraphs, here the entire complement $\T_0\setminus\T_1$: closure under induced
subgraphs says exactly that a graph with a forbidden induced subgraph cannot lie in $\T_1$.
So the forbidden flags already span an ideal, not merely a subspace, and the ideal they
generate is exactly their span.  Hence the elements that $\qmap_\sigma$ collapses to $0$ are exactly the linear
combinations of forbidden flags---nothing more.

Now let $\chi$ vanish on every forbidden flag.  By linearity it vanishes on their whole
span, which is the kernel just identified; equivalently, $\chi$ assigns the same value to
any two elements that $\qmap_\sigma$ identifies, so it cannot distinguish $\A^\sigma[\T_0]$
from its quotient.  By the universal property it therefore factors through $\qmap_\sigma$:
the assignment $\psi(\qmap_\sigma a):=\chi(a)$ is well defined, and it is the unique
homomorphism $\psi$ on
$\A^\sigma[\T_1]$ with $\chi=\psi\circ\qmap_\sigma$ (unique because $\qmap_\sigma$ is
surjective).  It inherits linearity, unit-preservation, and multiplicativity from $\chi$
because $\qmap_\sigma$ is a surjective algebra map, and it is non-negative on flags because
every flag of $\A^\sigma[\T_1]$ is $\qmap_\sigma(G)$ for a non-forbidden flag $G$, whence
$\psi(\qmap_\sigma G)=\chi(G)\ge0$.  Thus $\psi\in\Homp(\A^\sigma[\T_1],\R)$ and
$\chi=\psi\circ\qmap_\sigma\in Q_\sigma$.

This description also makes the closedness of $Q_\sigma$ transparent.  For a single
forbidden flag $F$, the condition $\chi(F)=0$ is a \emph{closed} one: the evaluation
$\chi\mapsto\chi(F)$ is continuous---recall that $X_\sigma$ carries the topology of
pointwise convergence of flag densities (\cref{sec:background})---so the homomorphisms
satisfying it form the preimage $\{\chi:\chi(F)=0\}$ of the closed set $\{0\}$.  The space
$Q_\sigma$ consists of the homomorphisms that satisfy all of these conditions at once, one
per forbidden flag, so it is an intersection of closed sets and is therefore closed in
$X_\sigma$.

Informally, the next lemma says that randomly choosing roots in a constrained unlabelled limit
almost surely yields a constrained labelled one.

\begin{lemma}[Support passes to random extensions]\label{lem:support-passes-general}
Let $\phi_0\in Q_0$ and suppose $\phi_0(\brak\sigma)>0$.  If
$\phi^\sigma\sim\Ext_\sigma(\phi_0)$ is formed in the ambient class $\T_0$,
then
\[
  \phi^\sigma\in Q_\sigma
  \qquad\text{almost surely.}
\]
\end{lemma}

\begin{proof}
Let $F$ be a forbidden $\sigma$-flag, and write $M\notin\T_1$ for its underlying unlabelled
graph.  The downward average $\avg{F}{\sigma}$ is a non-negative scalar multiple of $M$ as
an element of $\A^0[\T_0]$.  Since $\phi_0\in Q_0$ vanishes on every graph outside $\T_1$,
we have $\phi_0(M)=0$, and hence $\phi_0(\avg{F}{\sigma})=0$.  By \eqref{eq:extension-expectation},
\[
  \E[\phi^\sigma(F)]=0.
\]
But $\phi^\sigma(F)\ge0$ pointwise.  Hence $\phi^\sigma(F)=0$ almost surely.
Intersecting the resulting probability-one events over the countable set of
forbidden $\sigma$-flags shows that $\phi^\sigma\in Q_\sigma$ almost surely.
\end{proof}

\section{The support-closure criterion}

The following definition is the main conceptual point of this paper.

\begin{definition}[Root-planting set and root-plantability]\label{def:root-planting}
Fix a non-degenerate type $\sigma$ of $\T_1$; equivalently,
$\phi_0(\brak\sigma)>0$ for some constrained unlabelled limit $\phi_0\in Q_0$.  Define
\[
  S_\sigma
  :=
  \overline{
  \bigcup_{\substack{\phi_0\in Q_0\\ \phi_0(\brak\sigma)>0}}
  \supp(\Ext_\sigma(\phi_0))}
  \subseteq X_\sigma,
\]
where the closure is taken in the compact ambient space
$X_\sigma=\Homp(\A^\sigma[\T_0],\R)$.
We say that the triple $(\T_0,\T_1,\sigma)$ is \emph{root-plantable} if
\[
  S_\sigma=Q_\sigma.
\]
\end{definition}

By \cref{lem:support-passes-general}, one always has $S_\sigma\subseteq Q_\sigma$.
The point of root-plantability is the reverse inclusion: every constrained
labelled limit should be approximable by support points of random extensions from
constrained unlabelled limits.

\begin{theorem}[Support-closure criterion]\label{thm:support-criterion}
Let $\T_1$ be a hereditary subclass of $\T_0$, and let $\sigma$ be a non-degenerate type of
$\T_1$ (so $\phi_0(\brak\sigma)>0$ for some $\phi_0\in Q_0$).  For
$f\in\A^\sigma[\T_0]$, consider the following two
conditions.
\begin{enumerate}[label=\textup{(\alph*)}, leftmargin=2.8em]
  \item\label{cond:quotient}
  Quotient semantics:
  \[
    \chi(f)\ge0\qquad\text{for every }\chi\in Q_\sigma.
  \]
  Equivalently, $\psi(\qmap_\sigma f)\ge0$ for every
  $\psi\in\Homp(\A^\sigma[\T_1],\R)$.
  \item\label{cond:ensemble}
  Ensemble semantics:
  for every $\phi_0\in Q_0$ with $\phi_0(\brak\sigma)>0$,
  \[
    \Prob_{\phi^\sigma\sim\Ext_\sigma(\phi_0)}[\phi^\sigma(f)\ge0]=1.
  \]
\end{enumerate}
Then \ref{cond:quotient} always implies \ref{cond:ensemble}.  Moreover,
\ref{cond:quotient} and \ref{cond:ensemble} are equivalent for every
$f\in\A^\sigma[\T_0]$ if and only if $(\T_0,\T_1,\sigma)$ is root-plantable.
\end{theorem}

In words: a certificate proved in the quotient is always a valid constrained density
bound, and it captures \emph{every} such bound precisely when $\sigma$ is
root-plantable; when root-plantability fails, some inequality that holds almost surely
under random rooting cannot be seen in the quotient.

\begin{proof}
We use the standard evaluation pairing between the algebra and its positive
homomorphisms.  Thus an element $f\in\A^\sigma[\T_0]$ also denotes the continuous
function
\[
  X_\sigma\longrightarrow\R,\qquad \chi\longmapsto \chi(f).
\]
In this sense, ``$f\ge0$ on $S_\sigma$'' means $\chi(f)\ge0$ for every
$\chi\in S_\sigma$, and similarly for $Q_\sigma$.

By \cref{lem:support-passes-general}, every support appearing in the definition
of $S_\sigma$ is contained in $Q_\sigma$.  Hence non-negativity on $Q_\sigma$
implies non-negativity on all those supports, and by \cref{lem:support-as} this
is precisely the implication from quotient semantics to ensemble semantics.

For the converse direction, note first that ensemble semantics for a fixed $f$
is equivalent, again by \cref{lem:support-as}, to non-negativity of $f$ on the
union of the supports in \cref{def:root-planting}.  Since $f$ is continuous on
$X_\sigma$, this is equivalent to non-negativity on $S_\sigma$.
Thus the comparison is simply
\[
  f\ge0\text{ on }S_\sigma
  \qquad\text{versus}\qquad
  f\ge0\text{ on }Q_\sigma.
\]
If $S_\sigma=Q_\sigma$, the two conditions are equivalent for every $f$.

Conversely suppose $S_\sigma\ne Q_\sigma$.  Choose
$\psi\in Q_\sigma\setminus S_\sigma$.  The compact space $X_\sigma$ is a closed subspace of the product
$[0,1]^{\F^\sigma}$, hence is metrizable; so its subspace $Q_\sigma$ is metrizable as
well, and therefore normal in the topological sense: disjoint closed subsets can be
separated by a continuous function.  The sets
$S_\sigma$ and $\{\psi\}$ are disjoint closed subsets of $Q_\sigma$.
By Urysohn's lemma, followed by an affine rescaling, there is a continuous function
$h\in C(Q_\sigma)$ such that $h=1$ on $S_\sigma$ and $h(\psi)=-1$.
Also, $Q_\sigma$ is closed in $X_\sigma$ by the closedness observation in the previous
section; more explicitly, it is cut out by the equations $\chi(F)=0$ for forbidden flags
$F$.  Hence Tietze's extension theorem applies, and
$h$ extends to a continuous function $H\in C(X_\sigma)$.  Finally, the evaluation
functions coming from $\A^\sigma[\T_0]$ form a unital real subalgebra of
$C(X_\sigma)$ and separate points of $X_\sigma$, since distinct positive homomorphisms
disagree on some flag-algebra element.  Therefore Stone--Weierstrass, in the form used
by Razborov \cite[Proposition 3.7]{Razborov2007}, says that every continuous function on
$X_\sigma$ can be approximated arbitrarily well in sup norm by elements of
$\A^\sigma[\T_0]$.  Choose
$f\in\A^\sigma[\T_0]$ with $\|f-H\|_\infty<1/3$.  Then $f>0$ on
$S_\sigma$, so ensemble semantics holds for $f$, while $f(\psi)<0$, so quotient
semantics fails.  Thus equivalence for all $f$ forces $S_\sigma=Q_\sigma$.
\end{proof}

\section{Clone-closed graph classes are root-plantable}

We now give a checkable sufficient condition for root-plantability of graph
classes, built on an independent blow-up construction.  The idea is that when a class is
closed under independent blow-ups, a root can be duplicated into a clone class
of positive density, so a configuration the quotient exhibits only at an exceptional vertex
can instead be realised at a random---hence typical---vertex, which is just what
root-plantability requires.  Recall that $\T_0=\T_{\mathrm{Graph}}$ denotes the class of
all finite simple graphs.  If $\K$ is a hereditary class of finite graphs, we write
$\T_\K$ for the same class in the constrained-theory notation: equivalently, $\T_\K$
is the class of graphs avoiding every induced subgraph not in $\K$.  Thus in this
section $\T_1=\T_\K$ is just $\K$ viewed as the constrained class inside
$\T_{\mathrm{Graph}}$.

\begin{definition}[Independent blow-up]\label{def:independent-blow-up}
Let $G$ be a finite graph and let $m_v\ge1$ be an integer for each
$v\in V(G)$.  The independent blow-up $G^{\mathbf m}$ is obtained by replacing
each vertex $v$ by an independent clone class $C_v$ of size $m_v$, and by putting
all edges between $C_v$ and $C_w$ whenever $vw\in E(G)$.
\end{definition}

\begin{definition}[Clone-closed hereditary class]
A hereditary graph class $\K$ is \emph{clone-closed} if every independent
blow-up of every graph in $\K$ again belongs to $\K$.
\end{definition}

\begin{example}
For every $r\ge3$, the class of $K_r$-free graphs is clone-closed.  A clique in
an independent blow-up uses at most one vertex from each clone class, and hence
projects to a clique of the same size in the original graph.  In particular, the
triangle-free class is clone-closed.
\end{example}

\subsection{The planted blow-up estimate}

We first fix the notation used in the estimates.  Let $(G,\theta)$ be a
$\sigma$-flag with $|\sigma|=k$ and $|G|=n>k$, and fix a parameter
$0<\lambda\le 1/2$.  The blow-up will be used to make the planted roots
visible with positive probability, while keeping the densities
$p(F,\cdot)$ of every fixed $\sigma$-flag $F$ close to their values in
$(G,\theta)$.  For that purpose we first prescribe the desired
fractions of the future blow-up lying over each vertex of $G$: the roots
together receive total mass $\lambda$, and the unlabelled vertices
share the remaining mass uniformly.  These desired fractions are the
\emph{target weights}.  If $k>0$, set
\[
  w_{\theta(i)}=\frac{\lambda}{k}\quad(1\le i\le k),
  \qquad
  w_v=\frac{1-\lambda}{n-k}\quad(v\notin\operatorname{im}\theta).
\]
If $k=0$, set $w_v=1/n$ for all $v\in V(G)$.

The target weights are ideal limiting fractions, not a finite graph by
themselves.  To build an actual $N$-vertex blow-up, one must choose integer
clone-class sizes whose normalised sizes approximate these fractions.  Thus, for
an integer $N\ge n$, let
\[
  \mathbf m=(m_v)_{v\in V(G)}
\]
be any vector of positive integers with $\sum_v m_v=N$, and let
$B_N=G^{\mathbf m}$ be the corresponding independent blow-up.  We write $C_v$
for the clone class of $v$, so $|C_v|=m_v$.  The graph $B_N$ depends only on
$G$ and the chosen vector $\mathbf m$; the parameters $\theta$ and $\lambda$
enter through the target proportions $w_v$ and hence through the error term
\[
  \operatorname{err}_N
  :=
  \sum_{v\in V(G)}\left|\frac{|C_v|}{N}-w_v\right|.
\]
For fixed $G,\theta,\lambda$, the vectors $\mathbf m$ can be chosen, for
instance by rounding the numbers $w_vN$, so that
$\operatorname{err}_N\to0$ as $N\to\infty$.

\begin{lemma}[Planted blow-up approximation]\label{lem:planted-estimate}
Let $m\ge k$ be fixed.  There is a constant $C_m$, depending only on $m$ and
on no other data, such that the following holds for every
$0<\lambda\le 1/2$.  Let $F_0$ be any $\sigma$-flag with at most $m$ vertices,
let $(G,\theta)$ be a $\sigma$-flag on $n$ vertices with
$n\ge\max\{m,k+1\}$, and choose any positive clone-class sizes
$\mathbf m=(m_v)_{v\in V(G)}$ with total size $N\ge n$.  Let
$B_N=G^{\mathbf m}$ be the resulting independent blow-up, define
$\operatorname{err}_N$ from this $\mathbf m$ and from the target weights
defined above using $(G,\theta)$ and $\lambda$, and write $C_v$ for the clone
class of $v$.
If
$\widehat\theta:\sigma\to B_N$ satisfies
$\widehat\theta(i)\in C_{\theta(i)}$ for all $i$, then
\begin{equation}\label{eq:planted-estimate}
  \left|p(F_0,(B_N,\widehat\theta))-p(F_0,(G,\theta))\right|
  \le
  C_m\left(\lambda+\frac1{n-k}+\operatorname{err}_N\right).
\end{equation}
Consequently, if $f\in\A^\sigma[\T_{\mathrm{Graph}}]$ is represented by a finite
linear combination of $\sigma$-flags with at most $m$ vertices, then there is a
constant $C_f$ depending only on this chosen finite linear combination
(in particular, independent of $n$, $G$, $\theta$, $\lambda$, $N$,
$\mathbf m$, and $\widehat\theta$) such that
\begin{equation}\label{eq:planted-linear-estimate}
  \left|p(f,(B_N,\widehat\theta))-p(f,(G,\theta))\right|
  \le
  C_f\left(\lambda+\frac1{n-k}+\operatorname{err}_N\right).
\end{equation}
\end{lemma}

\begin{proof}
It suffices to prove \eqref{eq:planted-estimate}; the estimate for $f$ follows
by linearity because $f$ is a finite linear combination of flags.  Indeed, if
$f=\sum_r a_r F_r$ with each $|F_r|\le m$, one may take
$C_f=C_m\sum_r |a_r|$.

Before turning to the sampling argument, note that the size assumptions have two
roles.  First, $|F_0|\le m\le n$, so both
$p(F_0,(G,\theta))$ and $p(F_0,(B_N,\widehat\theta))$ are defined.  Second, the
inequality $N\ge n$ is used in
the finite denominator comparisons below.

Let $\ell=|F_0|\le m$ and put $q:=\ell-k$.  If $q=0$, then $F_0$ is just the
type $\sigma$ itself; no additional vertices are sampled and the two
probabilities are equal to $1$.  Hence assume $q>0$.  To compute
$p(F_0,(B_N,\widehat\theta))$, choose a uniformly random $q$-element subset of
$B_N\setminus\operatorname{im}\widehat\theta$.  For the estimates only, expose
this unordered subset as an ordered sample
\[
  X_1,\ldots,X_q
\]
drawn without replacement, and then forget the order.  This produces the same
uniform $q$-subset, and all events below depend only on that subset.  Let
\[
  U:=V(G)\setminus\operatorname{im}\theta,
  \qquad
  R:=C_{\theta(1)}\cup\cdots\cup C_{\theta(k)} .
\]
Thus $R$ is the union of the clone classes attached to the roots
(and $R=\emptyset$ if $k=0$).

The proof has two parts.  First we show that, with probability
$O_m(\lambda+1/(n-k)+\operatorname{err}_N)$, the sample is bad: it either meets
one of the root clone classes in $R$ or chooses two vertices from the same
unlabelled clone class.  Second, on the complementary good event, projection
from the blow-up to $G$ preserves the induced $\sigma$-flag.  The only remaining
question on the good event is whether the projected $q$-set is close to a
uniform $q$-subset of $U$; this will contribute only
$O_m(\operatorname{err}_N)$ in total variation.  Combining these two errors gives
\eqref{eq:planted-estimate}.

By the definition of $\operatorname{err}_N$,
\[
  \frac{|R|}{N}\le \lambda+\operatorname{err}_N .
\]
Here $R$ includes the $k$ planted roots themselves, although those
vertices are not in the sampling space; this only makes the following upper
bound larger.
Since each exposed position $X_j$ is marginally uniform on the
$N-k$ vertices of $B_N\setminus\operatorname{im}\widehat\theta$, we have
\[
  \Prob[X_j\in R]
  \le
  \frac{|R|}{N-k}
  =
  \frac{N}{N-k}\frac{|R|}{N}
  \le
  (m+1)(\lambda+\operatorname{err}_N).
\]
Here the last inequality uses only $N\ge n\ge k+1$ and $k\le m$, which imply
$N/(N-k)\le k+1\le m+1$.  A union bound over the $q\le m$ exposed positions then gives
\[
  \Prob[\text{the sample meets }R]
  =
  \Prob\!\left[\bigcup_{j=1}^q\{X_j\in R\}\right]
  \le
  K_m^{(1)}(\lambda+\operatorname{err}_N),
\]
for a constant $K_m^{(1)}$ depending only on $m$.

Now condition on the sample avoiding $R$.  Under this conditioning, the exposed
sequence is a uniform ordered sample without replacement from the union of the
unlabelled clone classes.  For each $v\in U$, write $s_v:=|C_v|$ and
$S:=\sum_{u\in U}s_u$.  Each exposed position then falls in the clone class
$C_v$ with marginal probability $s_v/S$.

Here is the normalisation estimate explicitly.  Put
\[
  \beta:=\sum_{v\in U}w_v
  =
  \begin{cases}
    1-\lambda,& k>0,\\
    1,& k=0,
  \end{cases}
  \qquad
  a_v:=\frac{s_v}{N},\qquad
  a:=\frac{\beta}{n-k},\qquad
  A:=\frac SN=\sum_{v\in U}a_v .
\]
The definition of $\operatorname{err}_N$ gives
\[
  \sum_{v\in U}
  |a_v-a|
  \le \operatorname{err}_N,
  \qquad
  |A-\beta|
  =
  \left|\sum_{v\in U}(a_v-a)\right|
  \le \operatorname{err}_N .
\]
Since $\beta\ge1/2$ and $a/\beta=1/(n-k)$, we have
\[
\begin{aligned}
  \sum_{v\in U}\left|\frac{s_v}{S}-\frac1{n-k}\right|
  &=
  \sum_{v\in U}\left|\frac{a_v}{A}-\frac{a}{\beta}\right|  \\
  &\le
  \sum_{v\in U}\left|\frac{a_v}{A}-\frac{a_v}{\beta}\right|
  +
  \sum_{v\in U}\left|\frac{a_v-a}{\beta}\right|  \\
  &=
  \frac{|A-\beta|}{\beta}
  +
  \frac1{\beta}\sum_{v\in U}|a_v-a|  \\
  &\le 4\,\operatorname{err}_N .
\end{aligned}
\]
Thus the probability vector
\[
  \left(\frac{s_v}{S}\right)_{v\in U}
\]
has $\ell^1$-distance at most $4\operatorname{err}_N$ from the uniform vector on
the $n-k$ vertices of $U$.  In particular, each coordinate satisfies
$s_v/S\le 1/(n-k)+K_m^{(0)}\operatorname{err}_N$ for a constant
$K_m^{(0)}$ depending only on $m$.

Now fix two exposed positions $a<b$.  We bound the probability that $X_a$ and
$X_b$ land in the same unlabelled clone class, still under the conditioning
that the whole sample avoids $R$.  If $X_a\in C_v$, then $X_b$ is uniform among
the remaining $S-1$ vertices in the unlabelled clone classes.  Therefore
\[
  \Prob[X_b\in C_v\mid X_a\in C_v,\ \text{the sample avoids }R]
  =
  \frac{s_v-1}{S-1}
  \le
  \frac{s_v}{S}.
\]
Averaging over the possible clone class $C_v$ of $X_a$ gives the explicit bound
\[
  \begin{aligned}
  &\Prob[X_a,X_b\text{ lie in the same unlabelled clone class}
        \mid \text{the sample avoids }R]  \\
  &\qquad =
  \sum_{v\in U}
    \frac{s_v}{S}\cdot\frac{s_v-1}{S-1}
  \le
  \sum_{v\in U}\left(\frac{s_v}{S}\right)^2
  \le
  \max_{v\in U}\frac{s_v}{S}
  \le
  \frac1{n-k}+K_m^{(0)}\operatorname{err}_N.
  \end{aligned}
\]
Let $E_{\mathrm{coll}}$ be the event that two sampled vertices
lie in the same unlabelled clone class.  Equivalently, for the ordered
exposure, $E_{\mathrm{coll}}$ is the union over pairs $a<b$ of the events that
$X_a$ and $X_b$ land in the same clone class.  A second union bound over the at
most $\binom q2\le\binom m2$ pairs gives
\[
  \Prob[E_{\mathrm{coll}}\mid \text{the sample avoids }R]
  \le
  K_m^{(2)}\left(\frac1{n-k}+\operatorname{err}_N\right),
\]
for a constant $K_m^{(2)}$ depending only on $m$.

Call the sample good if it avoids $R$ and no two sampled vertices lie in the same
unlabelled clone class.  Combining the two bounds above,
\[
  \Prob[\text{bad}]
  \le
  K_m^{(3)}\left(\lambda+\frac1{n-k}+\operatorname{err}_N\right),
\]
where $K_m^{(3)}$ again depends only on $m$.

It remains to compare what happens on the good event with the corresponding
$\sigma$-flag density in the base flag $(G,\theta)$.  Let
\[
  \pi:B_N\to G
\]
be the projection sending each clone class $C_v$ to its base vertex $v$.  On the
good event, the vertices $\pi(X_1),\ldots,\pi(X_q)$ are distinct elements of
$U$.

First compare their distribution with the uniform distribution on $q$ distinct
vertices of $U$.  Let $\mu$ be the probability vector on $U$ given by
$\mu(v)=s_v/S$, and let $\nu$ be the uniform probability vector on $U$.  The
normalisation estimate above says
\[
  \|\mu-\nu\|_1\le 4\operatorname{err}_N .
\]
Here is the same comparison for ordered distinct tuples.  Let $r:=|U|=n-k$, and
let $D$ be the set of ordered $q$-tuples of distinct elements of $U$.  On the
good event, the projected ordered tuple $(\pi(X_1),\ldots,\pi(X_q))$ has the
probability
\[
  P_\mu(v_1,\ldots,v_q)
  =
  \frac{\prod_{j=1}^q \mu(v_j)}
       {Z_\mu},
  \qquad
  Z_\mu:=\sum_{(u_1,\ldots,u_q)\in D}\prod_{j=1}^q \mu(u_j),
\]
for $(v_1,\ldots,v_q)\in D$.  Indeed, the numerator is proportional to the
number $s_{v_1}\cdots s_{v_q}$ of ways to choose one vertex from each prescribed
clone class, and $Z_\mu$ is the normalising constant after restricting to
distinct base vertices.  The uniform ordered distinct tuple has the analogous law
\[
  P_\nu(v_1,\ldots,v_q)
  =
  \frac{\prod_{j=1}^q \nu(v_j)}
       {Z_\nu},
  \qquad
  Z_\nu:=\sum_{(u_1,\ldots,u_q)\in D}\prod_{j=1}^q \nu(u_j).
\]
Since $\nu$ is uniform,
\[
  Z_\nu=\frac{r(r-1)\cdots(r-q+1)}{r^q}.
\]
Because $r\ge q$ and $q\le m$, this number is bounded below by a positive
constant $\alpha_m$ depending only on $m$.

The target for the next few estimates is the following total-variation bound,
where $d_{\mathrm{TV}}$ denotes total variation distance: for some constant
$K_m^{(4)}$ depending only on $m$,
\begin{equation}\label{eq:ordered-good-tv-target}
  d_{\mathrm{TV}}(P_\mu,P_\nu)\le K_m^{(4)}\operatorname{err}_N .
\end{equation}
After forgetting the order, the same bound will hold for the projected $q$-set
on the good event.

For $(v_1,\ldots,v_q)\in D$, write
\[
  w_\mu(v_1,\ldots,v_q):=\prod_{j=1}^q\mu(v_j),
  \qquad
  w_\nu(v_1,\ldots,v_q):=\prod_{j=1}^q\nu(v_j).
\]
The unnormalised weights are close.  For each tuple, a telescoping expansion gives
\[
  \left|
    \prod_{j=1}^q\mu(v_j)-\prod_{j=1}^q\nu(v_j)
  \right|
  \le
  \sum_{i=1}^q
  |\mu(v_i)-\nu(v_i)|
  \prod_{j<i}\mu(v_j)\prod_{j>i}\nu(v_j).
\]
Now sum this inequality over $D$.  This gives
\[
\begin{aligned}
&\sum_{(v_1,\ldots,v_q)\in D}
  \left|
    \prod_{j=1}^q\mu(v_j)-\prod_{j=1}^q\nu(v_j)
  \right|  \\
&\qquad\le
  \sum_{i=1}^q
  \sum_{(v_1,\ldots,v_q)\in D}
  |\mu(v_i)-\nu(v_i)|
  \prod_{j<i}\mu(v_j)\prod_{j>i}\nu(v_j).
\end{aligned}
\]
It remains to bound the inner sum for each fixed $i$.  All terms in it are
non-negative, and $D\subseteq U^q$, so we may enlarge the inner sum from $D$ to
all of $U^q$:
\[
\begin{aligned}
&\sum_{(v_1,\ldots,v_q)\in D}
  |\mu(v_i)-\nu(v_i)|
  \prod_{j<i}\mu(v_j)\prod_{j>i}\nu(v_j)  \\
&\qquad\le
  \sum_{(v_1,\ldots,v_q)\in U^q}
  |\mu(v_i)-\nu(v_i)|
  \prod_{j<i}\mu(v_j)\prod_{j>i}\nu(v_j)  \\
&\qquad=
  \left(\sum_{v_i\in U}|\mu(v_i)-\nu(v_i)|\right)
  \prod_{j<i}\left(\sum_{v_j\in U}\mu(v_j)\right)
  \prod_{j>i}\left(\sum_{v_j\in U}\nu(v_j)\right)  \\
&\qquad=
  \|\mu-\nu\|_1 .
\end{aligned}
\]
Summing this bound over the $q$ possible values of $i$ gives
\begin{equation}\label{eq:good-unnormalized-weight-bound}
  \sum_{(v_1,\ldots,v_q)\in D}
  |w_\mu(v_1,\ldots,v_q)-w_\nu(v_1,\ldots,v_q)|
  \le
  q\|\mu-\nu\|_1
  \le
  m\|\mu-\nu\|_1
  \le
  K_m''\operatorname{err}_N,
\end{equation}
where the last inequality uses $\|\mu-\nu\|_1\le4\operatorname{err}_N$ and
$K_m''$ depends only on $m$.  Since
$Z_\mu-Z_\nu=\sum_D(w_\mu-w_\nu)$, \eqref{eq:good-unnormalized-weight-bound} also gives
\begin{equation}\label{eq:good-normalizer-bound}
  |Z_\mu-Z_\nu|\le K_m''\operatorname{err}_N .
\end{equation}
If
$K_m''\operatorname{err}_N\ge\alpha_m/2$, then the target
\eqref{eq:ordered-good-tv-target} is automatic after increasing $K_m^{(4)}$:
indeed, total variation distance is always at most $1$, while
$1\le (2K_m''/\alpha_m)\operatorname{err}_N$.  Otherwise
$K_m''\operatorname{err}_N<\alpha_m/2$.  Since $Z_\nu\ge\alpha_m$ and
\eqref{eq:good-normalizer-bound} holds, we have
\[
  Z_\mu
  \ge
  Z_\nu-|Z_\mu-Z_\nu|
  \ge
  \alpha_m-\frac{\alpha_m}{2}
  =
  \frac{\alpha_m}{2}.
\]
In this case the normalisation can be estimated directly:
the two bounds being reused are
\eqref{eq:good-unnormalized-weight-bound} for $\sum_D|w_\mu-w_\nu|$ and
\eqref{eq:good-normalizer-bound} for $|Z_\mu-Z_\nu|$.  Thus
\[
\begin{aligned}
  \sum_D\left|\frac{w_\mu}{Z_\mu}-\frac{w_\nu}{Z_\nu}\right|
  &\le
  \frac1{Z_\mu}\sum_D|w_\mu-w_\nu|
  +
  \left|\frac1{Z_\mu}-\frac1{Z_\nu}\right|\sum_D w_\nu  \\
  &=
  \frac1{Z_\mu}\sum_D|w_\mu-w_\nu|
  +
  \frac{|Z_\mu-Z_\nu|}{Z_\mu}
  \le
  \frac{4K_m''}{\alpha_m}\operatorname{err}_N .
\end{aligned}
\]
Thus, after enlarging $K_m^{(4)}$ if necessary, the target
\eqref{eq:ordered-good-tv-target} holds in this case as well.
Forgetting the order cannot increase total variation distance, so the projected
$q$-set on the good event is within $K_m^{(4)}\operatorname{err}_N$ of a
uniformly chosen $q$-subset of $U$.

Finally, on the good event the induced $\sigma$-flag is preserved by projection.
The planted roots $\widehat\theta(i)\in C_{\theta(i)}$ have exactly the same
adjacencies to a sampled vertex in $C_v$ as $\theta(i)$ has to $v$ in $G$, and
two sampled unlabelled vertices lie in distinct clone classes, so their adjacency
is exactly the adjacency of their two base vertices.  Therefore the $\sigma$-flag
seen in $(B_N,\widehat\theta)$ is the same as the $\sigma$-flag seen in
$(G,\theta)$ on the projected $q$-set.

Let $h$ be the indicator that a $q$-set of base vertices induces $F_0$ together
with the labels.  Let $Q_{\mathrm{good}}$ be the distribution of the projected
$q$-set conditioned on the good event, and let $U_q$ be the uniform distribution
on $q$-subsets of $U$.  If $\gamma:=\Prob[\text{good}]$, then the contribution
from bad samples is some number $(1-\gamma)b$ with $0\le b\le1$, and
\[
  p(F_0,(B_N,\widehat\theta))
  =
  \gamma\,\E_{Q_{\mathrm{good}}}h+(1-\gamma)b .
\]
On the other hand,
\[
  p(F_0,(G,\theta))
  =
  \E_{U_q}h
  =
  \gamma\,\E_{U_q}h+(1-\gamma)\E_{U_q}h .
\]
Subtracting these two expressions gives
\[
\begin{aligned}
&\left|p(F_0,(B_N,\widehat\theta))-p(F_0,(G,\theta))\right| \\
&\qquad\le
  \gamma\left|\E_{Q_{\mathrm{good}}}h-\E_{U_q}h\right|
  +(1-\gamma)\left|b-\E_{U_q}h\right| \\
&\qquad\le
  d_{\mathrm{TV}}(Q_{\mathrm{good}},U_q)+\Prob[\text{bad}]  \\
&\qquad\le
  K_m^{(4)}\operatorname{err}_N
  +
  K_m^{(3)}\left(\lambda+\frac1{n-k}+\operatorname{err}_N\right).
\end{aligned}
\]
Taking the lemma's constant $C_m$ larger than $K_m^{(3)}+K_m^{(4)}$ proves
\eqref{eq:planted-estimate}.
\end{proof}

\begin{lemma}[Positive probability of the planted root]\label{lem:planted-mass}
Keep the blow-up notation introduced above: $B_N=G^{\mathbf m}$, with clone
classes $C_v$, target weights defined from $(G,\theta)$ and $\lambda$, and error
$\operatorname{err}_N$; write $k:=|\sigma|$.  Consider the following finite
random-rooting experiment on the unlabelled graph $B_N$: choose uniformly from
all embeddings $\widehat\theta:\sigma\to B_N$, where the roots of
$\sigma$ are kept in their prescribed order.  If either $k=0$, or $k>0$ and
$\operatorname{err}_N\le\lambda/(2k)$, then this experiment gives
probability at least
\[
  c_{k,\lambda}:=
  \begin{cases}
    1, & k=0,\\[1mm]
    \left(\dfrac{\lambda}{2k}\right)^k, & k>0,
  \end{cases}
\]
to the event that the chosen embedding satisfies
$\widehat\theta(i)\in C_{\theta(i)}$ for all $i=1,\ldots,k$.
\end{lemma}

\begin{proof}
For $k=0$ there is only the empty embedding, so the probability is $1$.  Assume
$k>0$.  The random-rooting experiment above is not a uniform map from $\sigma$ to $B_N$;
it is uniform over the finite set of embeddings of $\sigma$ into
$B_N$.  Every choice of vertices
$\widehat\theta(i)\in C_{\theta(i)}$ is an embedding of $\sigma$,
because $\theta$ is an embedding into $G$ and the blow-up preserves the
adjacencies between distinct base vertices.  Since
$\operatorname{err}_N\le\lambda/(2k)$, each planted clone class satisfies
\[
  \frac{|C_{\theta(i)}|}{N}
  \ge
  w_{\theta(i)}-\operatorname{err}_N
  =
  \frac{\lambda}{k}-\operatorname{err}_N
  \ge
  \frac{\lambda}{2k}.
\]
The clone classes $C_{\theta(1)},\ldots,C_{\theta(k)}$ are distinct, since
$\theta$ is injective.  Thus the choices of the $k$ roots may be made
independently: there are $|C_{\theta(i)}|$ choices for the image of the
$i$th root, and every resulting choice is one of the planted embeddings counted above.
Hence the number of planted embeddings is at least
\[
  \prod_{i=1}^k |C_{\theta(i)}|
  \ge
  \left(\frac{\lambda}{2k}\right)^k N^k .
\]
The last inequality is just the product of the individual lower bounds
$|C_{\theta(i)}|\ge(\lambda/(2k))N$.
The denominator in the random experiment is the total number of embeddings of
$\sigma$ into $B_N$.  This denominator is at most the number of all maps from the
$k$ roots to $V(B_N)$, namely $N^k$.  Therefore
the probability of choosing a planted embedding is at least
\[
  \frac{(\lambda/(2k))^kN^k}{N^k}
  =
  \left(\frac{\lambda}{2k}\right)^k .
\]
\end{proof}

\subsection{Root-plantability theorem}

\begin{theorem}[Clone-closed classes are root-plantable]\label{thm:clone-root-plantable}
Let $\K$ be a clone-closed hereditary graph class, let
$\T_1=\T_\K$, and let $\sigma$ be a non-degenerate type of $\T_\K$.  Then
\[
  S_\sigma=Q_\sigma.
\]
Consequently, quotient semantics and ensemble semantics are equivalent for every
$f\in\A^\sigma[\T_{\mathrm{Graph}}]$.
\end{theorem}

\begin{proof}
The inclusion $S_\sigma\subseteq Q_\sigma$ is \cref{lem:support-passes-general}.
We prove the reverse inclusion.  Let $\psi\in Q_\sigma$.  We must show that every
neighbourhood of $\psi$ in $X_\sigma$ intersects the support of some
$\Ext_\sigma(\phi_0)$ with $\phi_0\in Q_0$ and $\phi_0(\brak\sigma)>0$.  Write
$k:=|\sigma|$.

The topology on $X_\sigma$ is the subspace topology inherited from the product
\[
  [0,1]^{\F^\sigma[\T_{\mathrm{Graph}}]},
\]
where $\F^\sigma[\T_{\mathrm{Graph}}]$ is the set of $\sigma$-flags in the
ambient class of all finite graphs, and the coordinate indexed by a flag $F$
records the value $\chi(F)$.
Therefore finite-coordinate cylinders form a neighbourhood basis.  Concretely,
every neighbourhood of $\psi$ contains one specified by finitely many
$\sigma$-flags $F_1,\ldots,F_s$ and positive tolerances.  Shrinking the tolerances
to a common positive number $\eps$, it is enough to work with a neighbourhood of the form
\[
  U=
  \{\chi\in X_\sigma:
    |\chi(F_i)-\psi(F_i)|<\eps\text{ for all }i\}.
\]
By the representation theorem in the constrained class $\T_\K$, choose a
convergent sequence of $\sigma$-flags
$(G_t,\theta_t)$ with $G_t\in\K$, $|G_t|\to\infty$, and
$p(\cdot,(G_t,\theta_t))\to\psi$.  Let $m\ge k$ be an integer bounding the sizes
of all the flags $F_i$; if the list is empty, take $m=k$.
Choose $0<\lambda\le 1/2$ so small that $C_m\lambda<\eps/10$, where $C_m$ is the
constant from \cref{lem:planted-estimate}.  Then fix one $t$ large enough that
\[
  |G_t|\ge\max\{m,k+1\},
  \qquad
  |p(F_i,(G_t,\theta_t))-\psi(F_i)|<\eps/10
  \quad\text{for all }i,
  \qquad
  \frac{C_m}{|G_t|-|\sigma|}<\eps/10.
\]
Finally take a sequence of clone-size vectors
$\mathbf m^{(j)}=(m^{(j)}_v)_{v\in V(G_t)}$ with total sizes $N_j\to\infty$
and with $\operatorname{err}_{N_j}\to0$, where the error is computed from
$\mathbf m^{(j)}$ and the target weights attached to $(G_t,\theta_t)$ and
$\lambda$; for all large $j$, require
\[
  C_m\operatorname{err}_{N_j}<\eps/10
  \qquad
  \text{and}
  \qquad 
  \text{if }k>0,\text{ then } \operatorname{err}_{N_j}\le\lambda/(2k);
\]
discarding finitely many initial terms, assume these requirements hold for every
$j$, and put
\[
  B_j:=G_t^{\mathbf m^{(j)}}.
\]
Since $G_t\in\K$ and $\K$ is clone-closed, every $B_j$ lies in $\K$.  Passing to
a subsequence if necessary, the unlabelled graphs $B_j$ converge to some
$\phi_0\in Q_0$.  The counting argument in \cref{lem:planted-mass} gives a positive
lower bound, independent of $j$, on the density of embeddings of $\sigma$ in $B_j$;
hence $\phi_0(\brak\sigma)>0$.

Let $Y_\sigma:=[0,1]^{\F^\sigma[\T_{\mathrm{Graph}}]}$, so that
$X_\sigma\subseteq Y_\sigma$ is a closed subspace.  For each fixed $j$ and each
embedding $\widehat\theta\in\Emb_\sigma(B_j)$, form the point
$x_{j,\widehat\theta}\in Y_\sigma$ whose $F$-coordinate is
$p(F,(B_j,\widehat\theta))$ whenever $|F|\le |B_j|$; for the remaining
coordinates with $|F|>|B_j|$, fill in any fixed value, say $0$.  This convention
for too-large flags is irrelevant for limits, since $|B_j|\to\infty$ and every
fixed coordinate is eventually a genuine finite density.  Let $P_j$ be the
pushforward of the uniform probability measure on the finite set
$\Emb_\sigma(B_j)$ under the map
$\widehat\theta\mapsto x_{j,\widehat\theta}$.

Thus $P_j$ is a well-defined finite probability measure on $Y_\sigma$ for each
$j$.  We are not choosing embeddings for different $j$ on a common probability
space, nor claiming that a randomly chosen labelled sequence converges almost
surely.  The finite-graph form of the random-extension theorem is precisely a
weak-convergence statement for these measures: since $B_j\to\phi_0$ and
$\phi_0(\brak\sigma)>0$, the measures $P_j$ converge weakly to
$\Ext_\sigma(\phi_0)$, viewed as a probability measure on $Y_\sigma$ concentrated
on $X_\sigma$.  Equivalently for our purposes, every continuous function of
finitely many flag-density coordinates has the limiting distribution prescribed
by $\Ext_\sigma(\phi_0)$.

Consider the closed cylinder in the ambient product space
\[
  \widetilde C=
  \{x\in Y_\sigma:
    |x(F_i)-\psi(F_i)|\le\eps/2\text{ for all }i\}.
\]
Put $C:=\widetilde C\cap X_\sigma$.  Then $C\subseteq U$.  For all sufficiently
large $j$, write $C_v^{(j)}$ for the clone class of $v\in V(G_t)$ in $B_j$, and
call an embedding $\widehat\theta:\sigma\to B_j$ \emph{planted} if
$\widehat\theta(a)\in C_{\theta_t(a)}^{(j)}$ for each root $a$ of
$\sigma$.  We now explain what \cref{lem:planted-estimate} gives for a planted
embedding.  Since each $F_i$ has at most $m$ vertices, the lemma applies with
base flag $(G_t,\theta_t)$ and with the clone-size vector
$\mathbf m^{(j)}$.  Hence
\[
  |p(F_i,(B_j,\widehat\theta))-p(F_i,(G_t,\theta_t))|
  \le
  C_m\left(\lambda+\frac1{|G_t|-|\sigma|}+\operatorname{err}_{N_j}\right)
  <\frac{3\eps}{10}
\]
for every $i$, by the choices of $\lambda$, $t$, and $j$.  Also
$|p(F_i,(G_t,\theta_t))-\psi(F_i)|<\eps/10$, so the triangle inequality gives
\[
  |p(F_i,(B_j,\widehat\theta))-\psi(F_i)|<\frac{4\eps}{10}<\frac{\eps}{2}
  \qquad(1\le i\le s).
\]
In other words, the point $x_{j,\widehat\theta}\in Y_\sigma$ associated to every
planted embedding lies in the closed cylinder $\widetilde C$.

Now use the positive mass of planted embeddings.  By \cref{lem:planted-mass},
when an embedding of $\sigma$ into $B_j$ is chosen uniformly, the
probability that it is planted is at least $c_{|\sigma|,\lambda}$.  Since every
planted embedding contributes a point of $\widetilde C$, the pushforward measure
$P_j$ satisfies
\[
  P_j(\widetilde C)\ge c_{|\sigma|,\lambda}>0
\]
for all sufficiently large $j$.

Finally pass this positive finite mass to the limiting random-extension measure.
The set $\widetilde C$ is closed in $Y_\sigma$, and $P_j$ converges weakly to
$\Ext_\sigma(\phi_0)$.  Therefore the closed-set form of the Portmanteau theorem
gives
\[
  \Ext_\sigma(\phi_0)(\widetilde C)
  \ge
  \limsup_{j\to\infty} P_j(\widetilde C)
  \ge c_{|\sigma|,\lambda}>0.
\]
Since $\Ext_\sigma(\phi_0)$ is concentrated on $X_\sigma$, this says
$\Ext_\sigma(\phi_0)(C)>0$.  A closed set of positive measure must meet the
support of that measure, so the support of $\Ext_\sigma(\phi_0)$ intersects
$C\subseteq U$.  Thus the support intersects this basic neighbourhood of $\psi$.
Because every neighbourhood of $\psi$ contains such a basic neighbourhood, every
neighbourhood of $\psi$ meets the support of some admissible random extension.
Hence $\psi\in S_\sigma$.
\end{proof}

\begin{corollary}[Clique-free and triangle-free classes]\label{cor:clique-free}
Fix $r\ge3$, and let $\K$ be the class of $K_r$-free graphs.  Let $\sigma$ be a
non-degenerate type in $\T_\K$, and let $f\in\A^\sigma[\T_{\mathrm{Graph}}]$.
Then the following two conditions, formed with respect to the constrained class
$\T_\K$, are equivalent:
\[
  \chi(f)\ge0\quad\text{for every }\chi\in Q_\sigma
\]
and
\[
  \Prob_{\phi^\sigma\sim\Ext_\sigma(\phi_0)}
  [\phi^\sigma(f)\ge0]=1
  \quad\text{for every }\phi_0\in Q_0
  \text{ with }\phi_0(\brak\sigma)>0.
\]
In particular, when $r=3$, the same equivalence holds for every non-degenerate
triangle-free type $\sigma$ and every $f\in\A^\sigma[\T_{\mathrm{Graph}}]$.
\end{corollary}

\begin{proof}
The class of $K_r$-free graphs is hereditary, and it is clone-closed by the
example above.  Therefore \cref{thm:clone-root-plantable} applies to $\T_\K$ and
gives root-plantability for every non-degenerate type $\sigma$ in this class.
The stated equivalence is exactly the quotient/ensemble equivalence in the final
sentence of that theorem.  The triangle-free case is the special case $r=3$.
\end{proof}

\section{Complete blow-ups and true twins}

The previous theorem uses the false-twin mechanism.  Let us spell out the
terminology.  In a finite graph, two distinct vertices $x$ and $y$ are
\emph{false twins} if they are not adjacent and have the same neighbours outside
the pair $\{x,y\}$: for every other vertex $z$, we have $zx\in E$ if and only
if $zy\in E$.  Duplicating a vertex as a false twin therefore creates a new
vertex with the same external neighbourhood as the old one, but with no edge
between the old and new copies.  Repeating this operation replaces each vertex
by an independent clone class, exactly as in the independent blow-ups of
\cref{def:independent-blow-up}.  Conversely, each one-vertex false-twin
duplication is a special independent blow-up.  Thus false-twin closure is
clone-closure in another language.

The true-twin analogue is different.  Two vertices are \emph{true twins} if
they are adjacent and have the same neighbours outside the pair.  Thus
duplicating a vertex as a true twin creates a new copy joined to the old one,
and repeated true-twin duplications replace vertices by cliques rather than by
independent sets.

\begin{definition}[Complete blow-up]\label{def:complete-blow-up}
Let $G$ be a finite graph and let $m_v\ge1$ be an integer for each
$v\in V(G)$.  The complete blow-up $G^{\mathbf m,+}$ is obtained by replacing
each vertex $v$ by a clique $C_v$ of size $m_v$, and by putting all edges
between $C_v$ and $C_w$ whenever $vw\in E(G)$.
\end{definition}

\begin{definition}[True-clone-closed hereditary class]
A hereditary graph class $\K$ is \emph{true-clone-closed} if every complete
blow-up of every graph in $\K$ again belongs to $\K$.
\end{definition}

\begin{lemma}[Complete blow-up approximation]\label{lem:true-planted-estimate}
The estimates of \cref{lem:planted-estimate,lem:planted-mass} remain valid,
with the same constants and the same hypotheses on $N$ and $\operatorname{err}_N$,
if the independent blow-up $B_N=G^{\mathbf m}$ is replaced by the complete
blow-up $B_N^+=G^{\mathbf m,+}$ with the same clone class sizes.
\end{lemma}

\begin{proof}
Work in the notation of
\cref{lem:planted-estimate,lem:planted-mass}; in particular, write
$k:=|\sigma|$ and $n:=|G|$.
We compare the complete blow-up $B_N^+=G^{\mathbf m,+}$ with the independent
blow-up $B_N=G^{\mathbf m}$ using the same clone classes $C_v$ and the same
projection map $\pi:C_v\mapsto v$.  The only difference between the two finite
graphs is internal to the clone classes: in $B_N$ each $C_v$ is independent,
whereas in $B_N^+$ each $C_v$ is a clique.  All edges between distinct clone
classes are governed by the same base graph $G$.

First consider the approximation estimate for a $\sigma$-flag $F_0$ with
$|F_0|\le m$.  Fix a planted embedding $\widehat\theta:\sigma\to B_N^+$ with
$\widehat\theta(i)\in C_{\theta(i)}$, and compute
$p(F_0,(B_N^+,\widehat\theta))$ by sampling the $q=|F_0|-|\sigma|$ additional
vertices from $B_N^+\setminus\operatorname{im}\widehat\theta$.  Use exactly the
same bad event as in \cref{lem:planted-estimate}: the sample is bad if it meets
one of the root clone classes or if two sampled vertices lie in the same
unlabelled clone class.  The probability of this bad event depends only on the
sizes of the clone classes and on sampling without replacement.  The two
union-bound estimates in the proof of \cref{lem:planted-estimate} therefore
apply without change and give
\[
  \Prob[\text{bad}]
  \le
  K_m^{(3)}\left(\lambda+\frac1{n-k}+\operatorname{err}_N\right),
\]
with the same constant $K_m^{(3)}$ used there.

On the complementary good event, the extra edges inside the complete clone
classes are never observed.  Indeed, no sampled vertex lies in a root clone
class, and no two sampled vertices lie in the same unlabelled clone class.
Thus every adjacency involving a sampled vertex is an adjacency between two
distinct clone classes, hence is exactly the adjacency of the corresponding base
vertices in $G$.  The planted roots also lie in distinct clone
classes, because $\theta$ is injective; since $\theta$ is an embedding of
the type $\sigma$ into $G$, their mutual adjacencies are the same as in the base
type.
Consequently, on the good event, the induced $\sigma$-flag seen in
$(B_N^+,\widehat\theta)$ is exactly the induced $\sigma$-flag seen in
$(G,\theta)$ after projecting the sampled vertices to $G$.

It remains only to compare the projected good sample with a uniform sample in
$V(G)\setminus\operatorname{im}\theta$.  This comparison uses only the numbers
$|C_v|$ and not the edge relation inside any clone class.  The normalisation
estimate for the vector $(|C_v|/S)_{v\notin\operatorname{im}\theta}$, the
collision estimate, and the total-variation comparison of the projected good
sample with a uniform sample are therefore exactly the estimates already proved
in \cref{lem:planted-estimate}.  Combining these clone-size estimates with the
good-event preservation just described gives
\[
  \left|p(F_0,(B_N^+,\widehat\theta))-p(F_0,(G,\theta))\right|
  \le
  C_m\left(\lambda+\frac1{n-k}+\operatorname{err}_N\right),
\]
with the same constant $C_m$.  The extension from a single flag $F_0$ to a
finite linear combination is by linearity: if $f=\sum_r a_rF_r$ with
$|F_r|\le m$, then the same argument gives the bound with
$C_f=C_m\sum_r|a_r|$.

Now consider the planted-mass estimate.  Choose a uniformly random embedding of
$\sigma$ into $B_N^+$.  Every choice of one vertex from each root clone class
$C_{\theta(i)}$ gives an embedding of $\sigma$: the root clone classes are
distinct, and all adjacencies between them are inherited from the base flag
$(G,\theta)$.  Since the clone classes are distinct, these choices are independent:
for the $i$th root there are $|C_{\theta(i)}|$ possible images, and the product counts
these choices of embeddings.  The
hypothesis $\operatorname{err}_N\le\lambda/(2k)$ gives, as in
\cref{lem:planted-mass},
$|C_{\theta(i)}|\ge(\lambda/(2k))N$ for every $i$.  Multiplying these $k$
lower bounds gives the same lower bound on the number of planted embeddings,
\[
  \prod_{i=1}^k |C_{\theta(i)}|
  \ge
  \left(\frac{\lambda}{2k}\right)^k N^k
  \qquad(k>0),
\]
under the same hypothesis $\operatorname{err}_N\le\lambda/(2k)$.  The total
number of embeddings is still at most the number $N^k$ of all maps from the
$k$ roots to $V(B_N^+)$.  Dividing gives the
same lower bound $c_{k,\lambda}$, and the case $k=0$ is again trivial.
\end{proof}

\begin{theorem}[True-clone-closed classes are root-plantable]\label{thm:true-clone-root-plantable}
Let $\K$ be a true-clone-closed hereditary graph class, let
$\T_1=\T_\K$, and let $\sigma$ be a non-degenerate type of $\T_\K$.  Then
\[
  S_\sigma=Q_\sigma.
\]
Consequently, quotient semantics and ensemble semantics are equivalent for every
$f\in\A^\sigma[\T_{\mathrm{Graph}}]$.
\end{theorem}

\begin{proof}
The inclusion $S_\sigma\subseteq Q_\sigma$ is the general support-passing
statement, \cref{lem:support-passes-general}.  We prove the reverse inclusion.
Fix $\psi\in Q_\sigma$.  It suffices to show that every neighbourhood of $\psi$
in $X_\sigma$ meets the support of some admissible random extension.

As before, $X_\sigma$ has the subspace topology inherited from
\[
  [0,1]^{\F^\sigma[\T_{\mathrm{Graph}}]},
\]
so finite-coordinate cylinders form a neighbourhood basis.  Thus every
neighbourhood of $\psi$ contains one of the form
\[
  U=
  \{\chi\in X_\sigma:
    |\chi(F_i)-\psi(F_i)|<\eps\text{ for all }i=1,\ldots,s\}
\]
for finitely many $\sigma$-flags $F_1,\ldots,F_s$ and some $\eps>0$, after
shrinking the coordinate tolerances to a common value.  We will find a
constrained unlabelled limit $\phi_0\in Q_0$ with $\phi_0(\brak\sigma)>0$ whose
random $\sigma$-extension has support meeting this $U$.

Since $\psi\in Q_\sigma$, the representation theorem in the constrained class
$\T_\K$ gives a convergent sequence of $\sigma$-flags
$(G_t,\theta_t)$ with $G_t\in\K$, $|G_t|\to\infty$, and
$p(\cdot,(G_t,\theta_t))\to\psi$.  Let $m\ge|\sigma|$ bound the sizes of the
finitely many flags $F_i$; if the list is empty, take $m=|\sigma|$.  Choose
$0<\lambda\le1/2$ so small that $C_m\lambda<\eps/10$, where $C_m$ is the
constant from \cref{lem:true-planted-estimate}.  Then choose $t$ large enough
that
\[
  |G_t|\ge\max\{m,|\sigma|+1\},\qquad
  |p(F_i,(G_t,\theta_t))-\psi(F_i)|<\eps/10\quad(1\le i\le s),
\]
and
\[
  \frac{C_m}{|G_t|-|\sigma|}<\eps/10 .
\]

For this fixed base flag $(G_t,\theta_t)$, choose clone-size vectors
$\mathbf m^{(j)}=(m_v^{(j)})_{v\in V(G_t)}$ with total sizes $N_j\to\infty$ and
with errors $\operatorname{err}_{N_j}\to0$, computed using the same target
weights determined by $(G_t,\theta_t)$ and $\lambda$.  Discarding finitely many
initial terms, assume that for every $j$,
\[
  C_m\operatorname{err}_{N_j}<\eps/10
  \qquad\text{and}\qquad
  \text{if }|\sigma|>0,\text{ then }
  \operatorname{err}_{N_j}\le\lambda/(2|\sigma|).
\]
Put
\[
  B_j^+:=G_t^{\mathbf m^{(j)},+}.
\]
Write $C_v^{(j)}$ for the clone class of $v\in V(G_t)$ in $B_j^+$.  An embedding
$\widehat\theta:\sigma\to B_j^+$ will be called \emph{planted} if
$\widehat\theta(a)\in C_{\theta_t(a)}^{(j)}$ for every root $a$ of
$\sigma$.
Because $G_t\in\K$ and $\K$ is true-clone-closed, every $B_j^+$ belongs to
$\K$.  Passing to a subsequence if necessary, the unlabelled graphs $B_j^+$
converge to some $\phi_0\in Q_0$.

The planted-mass part of \cref{lem:true-planted-estimate} gives a constant
$c_{|\sigma|,\lambda}>0$, independent of $j$, such that a uniformly random
embedding of $\sigma$ into $B_j^+$ is planted with probability at least
$c_{|\sigma|,\lambda}$.  In particular, the density of embeddings of
$\sigma$ in $B_j^+$ stays bounded below by a positive constant, so the limiting
unlabelled homomorphism satisfies $\phi_0(\brak\sigma)>0$.

Let
\[
  Y_\sigma:=[0,1]^{\F^\sigma[\T_{\mathrm{Graph}}]}.
\]
For each $j$, define a probability measure $P_j$ on $Y_\sigma$ as follows:
choose uniformly an embedding
$\widehat\theta\in\Emb_\sigma(B_j^+)$ and record the labelled flag-density vector
$x_{j,\widehat\theta}$ with coordinates
$x_{j,\widehat\theta}(F)=p(F,(B_j^+,\widehat\theta))$ whenever this finite
density is defined; for coordinates indexed by flags larger than $B_j^+$, fill
in any fixed value, say $0$.  This convention is irrelevant for weak limits,
because each fixed flag eventually has at most $|B_j^+|$ vertices.
Equivalently, $P_j$ is the pushforward of the uniform measure on
$\Emb_\sigma(B_j^+)$ under
$\widehat\theta\mapsto x_{j,\widehat\theta}$.  The finite-graph random-extension
theorem gives
\[
  P_j \Rightarrow \Ext_\sigma(\phi_0)
\]
weakly as probability measures on $Y_\sigma$, with the limit measure
concentrated on $X_\sigma$.

Now define the closed cylinder
\[
  \widetilde C=
  \{x\in Y_\sigma:
    |x(F_i)-\psi(F_i)|\le\eps/2\text{ for all }i=1,\ldots,s\}
\]
and put $C:=\widetilde C\cap X_\sigma$.  Then $C\subseteq U$.  If
$\widehat\theta$ is a planted embedding into $B_j^+$, the approximation part of
\cref{lem:true-planted-estimate} gives, for every $i$,
\[
  |p(F_i,(B_j^+,\widehat\theta))-p(F_i,(G_t,\theta_t))|
  \le
  C_m\left(\lambda+\frac1{|G_t|-|\sigma|}
      +\operatorname{err}_{N_j}\right)
  < \frac{3\eps}{10}.
\]
Together with
$|p(F_i,(G_t,\theta_t))-\psi(F_i)|<\eps/10$, this implies
\[
  |p(F_i,(B_j^+,\widehat\theta))-\psi(F_i)|<\eps/2 .
\]
Thus every planted embedding contributes a point of $\widetilde C$.  Since
planted embeddings have probability at least $c_{|\sigma|,\lambda}$,
\[
  P_j(\widetilde C)\ge c_{|\sigma|,\lambda}
\]
for every $j$.

Because $\widetilde C$ is closed in $Y_\sigma$, Portmanteau gives
\[
  \Ext_\sigma(\phi_0)(\widetilde C)
  \ge
  \limsup_j P_j(\widetilde C)
  \ge
  c_{|\sigma|,\lambda}>0 .
\]
The measure $\Ext_\sigma(\phi_0)$ is concentrated on $X_\sigma$, so
$\Ext_\sigma(\phi_0)(C)>0$.  Hence the support of
$\Ext_\sigma(\phi_0)$ meets $C\subseteq U$.  Since $U$ was an arbitrary basic
neighbourhood of $\psi$, every neighbourhood of $\psi$ meets the support of some
admissible random extension, and therefore $\psi\in S_\sigma$.  This proves
$Q_\sigma\subseteq S_\sigma$, hence $S_\sigma=Q_\sigma$.

The final assertion follows from the support-closure criterion,
\cref{thm:support-criterion}: once $S_\sigma=Q_\sigma$, quotient semantics and
ensemble semantics are equivalent for every
$f\in\A^\sigma[\T_{\mathrm{Graph}}]$.
\end{proof}

\begin{corollary}[Cluster graphs]\label{cor:cluster-graphs}
The class of cluster graphs, equivalently finite disjoint unions of cliques, is
root-plantable at every non-degenerate type.  It is not clone-closed.
\end{corollary}

\begin{proof}
Let $\K$ be the class of cluster graphs, i.e. finite disjoint unions of cliques.
This class is hereditary: deleting vertices from a disjoint union of cliques
again gives a disjoint union of cliques.

We first check true-clone-closure.  Let $G\in\K$, and write its connected
components as cliques $D_1,\ldots,D_r$.  In the complete blow-up
$G^{\mathbf m,+}$, the union $\bigcup_{v\in D_a} C_v$ is a clique: each
individual clone class is a clique, and all edges between clone classes over
adjacent vertices of $D_a$ are present.  No edges are added between clone
classes lying over different components.  Thus $G^{\mathbf m,+}$ is again a
disjoint union of cliques.  Hence $\K$ is true-clone-closed, so
\cref{thm:true-clone-root-plantable} gives root-plantability at every
non-degenerate type.

It remains to show that $\K$ is not clone-closed.  The graph $K_2$ belongs to
$\K$, but its independent blow-up with both clone classes of size $2$ is
$K_{2,2}$, the complete bipartite graph with two vertices on each side.  This
graph contains an induced $P_3$, the path on three vertices, obtained by taking
two vertices from one side and one from the other.  Since $P_3$ is not a
disjoint union of cliques and $\K$ is hereditary, $K_{2,2}\notin\K$.  Thus an
independent blow-up of a graph in $\K$ can leave $\K$, so $\K$ is not
clone-closed.
\end{proof}

\section{A common generalisation: blow-up-closed classes}\label{sec:blowup-closed}

The proofs of \cref{thm:clone-root-plantable,thm:true-clone-root-plantable} differ only
in what is put inside each clone class.  Inspecting the estimates shows that this
interior is never observed.  On the good sampling event, the sampled vertices lie in
distinct clone classes, so every adjacency that matters is a between-class adjacency
and is governed by the base graph $G$ alone
(\cref{lem:planted-estimate,lem:true-planted-estimate}).  Independent blow-ups and
complete blow-ups therefore differ only in data the sampling argument ignores.  The same
proof works for any graph placed inside each blown-up vertex.  This section isolates the
exact closure hypothesis needed for that argument; clone-closure, true-clone-closure,
and substitution-closure are all special cases.

\begin{definition}[Vertex blow-up]\label{def:vertex-blowup}
Let $G$ be a finite graph, $v\in V(G)$, and $H$ a finite graph whose vertex set is
disjoint from $G$.  The
\emph{blow-up of $v$ to $H$}, written $G[v\to H]$, is obtained by deleting $v$ and adding
a disjoint copy of $H$, joining every vertex of $H$ to every $G$-neighbour of $v$ and
keeping the internal edges of $H$.  (Equivalently $G[v\to H]=G[H_w:w\in V(G)]$ with
$H_v=H$ and $H_w$ a single vertex for $w\ne v$.)
\end{definition}

Taking $H$ edgeless recovers a one-vertex independent blow-up
(\cref{def:independent-blow-up}); taking $H$ complete recovers a one-vertex complete
blow-up (\cref{def:complete-blow-up}).

\begin{definition}[Blow-up-closed hereditary class]\label{def:blow-up-closed}
A hereditary graph class $\K$ is \emph{blow-up-closed} if for every $G\in\K$, every
$v\in V(G)$, and every $N\in\mathbb N$ there is a finite graph $H$ with $|H|\ge N$ such
that $G[v\to H]\in\K$.
\end{definition}

In words: in a blow-up-closed class one may always blow up a single vertex to an
arbitrarily large graph without leaving the class.  The class chooses the interior $H$;
the graph $H$ is not required to belong to $\K$ by itself.  This is the key difference
from substitution-closure, which demands that \emph{every} in-class interior be allowed
(see \cref{rem:strictness} below).

\begin{lemma}[From single vertices to a full uniform blow-up]\label{lem:blowup-iterate}
Let $\K$ be blow-up-closed, $G\in\K$, and $m\ge1$.  Then there are graphs
$(H_v)_{v\in V(G)}$, each with $|H_v|=m$, such that the simultaneous blow-up
$G[H_v:v\in V(G)]$ belongs to $\K$.
\end{lemma}

\begin{proof}
List the vertices of the original graph as $v_1,\dots,v_n$ and process them in that
order.  Put $G_0=G$.  At stage $i$, the vertices $v_i,\dots,v_n$ have not yet been
replaced, while $v_1,\dots,v_{i-1}$ have already been replaced by their new fibres.  Since
$G_{i-1}\in\K$, apply \cref{def:blow-up-closed} inside the current graph $G_{i-1}$ to the
still-present vertex $v_i$ with $N=m$.  This gives a graph $H'_i$ with $|H'_i|\ge m$ and
$G_{i-1}[v_i\to H'_i]\in\K$.

Choose an induced subgraph $H_i\subseteq H'_i$ on exactly $m$ vertices.  By heredity,
$G_{i-1}[v_i\to H_i]$ is again in $\K$, because it is an induced subgraph of
$G_{i-1}[v_i\to H'_i]$.  Set $G_i=G_{i-1}[v_i\to H_i]$.
Replacing $v_i$ does not change any fibre already created over
$v_1,\dots,v_{i-1}$, and every new vertex in $H_i$ inherits exactly the old adjacencies of
$v_i$ to the vertices that remain or have already been blown up.  Thus, after all
vertices have been processed, the resulting graph is precisely the simultaneous blow-up
$G[H_v:v\in V(G)]$ with $|H_v|=m$ for every original vertex $v$, and it lies in $\K$.
\end{proof}

The remaining ingredient is that the planting estimates
\cref{lem:planted-estimate,lem:true-planted-estimate} hold for \emph{arbitrary} interiors,
not just independent and complete ones.  This is the formal content of the opening remark:
the interior is never observed.

\begin{lemma}[Planting is blind to the interior]\label{lem:general-planting-estimate}
Let $(G,\theta)$ be a $\sigma$-flag with $|\sigma|=k$ and $|G|=n>k$.  Let
$B=G[H_v:v\in V(G)]$ be any simultaneous blow-up of $G$, with
$s_v:=|H_v|\ge1$ and total size $N=\sum_v s_v$.  Fix $0<\lambda\le1/2$, define the target
weights $w_v$ from $(G,\theta)$ and $\lambda$ as in \cref{lem:planted-estimate}, and put
\[
  \operatorname{err}_N:=\sum_{v\in V(G)}\left|\frac{s_v}{N}-w_v\right|.
\]
Then the conclusions of \cref{lem:planted-estimate,lem:planted-mass} hold for $B$ in
place of the independent blow-up with the same constants.  In particular, for every
integer $m\ge k$, every planted embedding
$\widehat\theta$ (meaning $\widehat\theta(i)\in H_{\theta(i)}$ for each labelled
vertex $i$), and every test $\sigma$-flag $F_0$ with at most $m$ vertices,
\[
  \left|p(F_0,(B,\widehat\theta))-p(F_0,(G,\theta))\right|
  \le
  C_m\left(\lambda+\frac1{n-k}+\operatorname{err}_N\right),
\]
and, if $k=0$ or $\operatorname{err}_N\le\lambda/(2k)$, the planted embeddings have
probability at least $c_{k,\lambda}$ under a uniformly random embedding of
$\sigma$ into $B$.  Here $C_m$ and $c_{k,\lambda}$ are the constants from
\cref{lem:planted-estimate,lem:planted-mass}.
\end{lemma}

\begin{proof}
This is the proof of \cref{lem:true-planted-estimate} with ``complete'' replaced by
``arbitrary''.  For the approximation estimate, fix a test flag $F_0$, sample the
$q=|F_0|-k$ extra vertices, and condition on the good event that no sampled vertex lies in
a root clone class and that no two sampled vertices lie in the same clone class.  On
this event every sampled vertex lies in a distinct clone class, so every adjacency
involving a sampled vertex is an adjacency between two \emph{distinct} clone classes and
therefore equals the corresponding base adjacency in $G$, whatever the interiors $H_v$
are.  Hence the induced $\sigma$-flag seen in $B$ is exactly the one seen in $(G,\theta)$
after projecting to $G$, and an interior edge is never observed.  The probability of the
bad event and the total-variation comparison with a uniform sample depend only on the
sizes $s_v$ and on sampling without replacement, exactly as in
\cref{lem:planted-estimate}.

For the planted-mass bound, every choice of one vertex from each root clone class
$H_{\theta(i)}$ gives an embedding of $\sigma$: the root clone classes are distinct
and their mutual adjacencies are inherited from $(G,\theta)$.  So the planted count
$\prod_i s_{\theta(i)}$ and its lower bound are unchanged.  By the same token an
embedding of $\sigma$ already arises from any transversal of $k$ root clone classes
whose projected vertices form an embedding of $\sigma$ in $G$, so the positive lower
bound on the $\sigma$-type density of $B$ also carries over.
\end{proof}

\begin{theorem}[Blow-up-closed classes are root-plantable]\label{thm:blowup-root-plantable}
Let $\K$ be a blow-up-closed hereditary graph class, let $\T_1=\T_\K$, and let $\sigma$
be a non-degenerate type of $\T_\K$.  Then $S_\sigma=Q_\sigma$, and consequently
quotient semantics and ensemble semantics are equivalent for every
$f\in\A^\sigma[\T_{\mathrm{Graph}}]$.
\end{theorem}

\begin{proof}
The inclusion $S_\sigma\subseteq Q_\sigma$ is the general support-passing statement,
\cref{lem:support-passes-general}.  We prove the reverse inclusion.  Fix
$\psi\in Q_\sigma$, and let $U$ be a basic neighbourhood of $\psi$ in $X_\sigma$, say
\[
  U=\{\chi\in X_\sigma:
      |\chi(F_i)-\psi(F_i)|<\eps\text{ for }i=1,\ldots,s\},
\]
where the $F_i$ are finitely many $\sigma$-flags.  Write $k=|\sigma|$, and let
$m\ge k$ bound the sizes of the $F_i$.

Represent $\psi$ by in-class flags $(G_t,\theta_t)$ with $G_t\in\K$,
$|G_t|\to\infty$, and $p(\cdot,(G_t,\theta_t))\to\psi$.  Choose one large base flag
$(G_t,\theta_t)$, and write $n=|G_t|$, so that
\[
  |p(F_i,(G_t,\theta_t))-\psi(F_i)|<\eps/10
  \quad(1\le i\le s),
  \qquad
  \frac{C_m}{n-k}<\eps/10.
\]
If $k>0$, also choose $t$ so large that $k/n\le1/2$ and $C_m k/n<\eps/10$, and set
$\lambda:=k/n$.  If $k=0$, choose any $0<\lambda\le1/2$ with
$C_m\lambda<\eps/10$.  This is the small point hidden by the phrase ``uniform blow-up'':
for $k>0$, uniform clone classes have normalised weight $1/n$ at every vertex, exactly
matching the target weights for $\lambda=k/n$.

For each $j\ge1$, \cref{lem:blowup-iterate} gives a simultaneous blow-up
\[
  B_j=G_t[H_v^{(j)}:v\in V(G_t)]\in\K,
  \qquad |H_v^{(j)}|=j\text{ for every }v.
\]
Thus the normalised clone-class sizes are all $1/n$.  With the choice of $\lambda$ above,
the error term in \cref{lem:general-planting-estimate} is $0$ when $k>0$; when $k=0$ it is
also $0$, since the target weights are uniform.  Passing to a subsequence, the unlabelled
graphs $B_j$ converge to some constrained limit $\phi_0\in Q_0$.  Moreover, there are
many embeddings of the type in $B_j$: for $k>0$, choosing one vertex from each labelled
clone class gives $j^k=(1/n)^k|V(B_j)|^k$ embeddings of $\sigma$, and for $k=0$ the
empty type has density $1$.  Thus the $\sigma$-density of $B_j$ is
bounded below by a positive constant independent of $j$, so
$\phi_0(\brak\sigma)>0$.  The planted-mass estimate also gives a positive lower bound,
independent of $j$, for the probability that a uniformly random embedding of
$\sigma$ into $B_j$ is one of these planted embeddings.

Let $Y_\sigma=[0,1]^{\F^\sigma[\T_{\mathrm{Graph}}]}$, and let $P_j$ be the pushforward
to $Y_\sigma$ of the uniform measure on embeddings
$\widehat\theta:\sigma\to B_j$, recording all finite flag densities.  The finite-graph
random-extension theorem gives
\[
  P_j\Rightarrow \Ext_\sigma(\phi_0),
\]
with the limit measure concentrated on $X_\sigma$.  Define the closed cylinder
\[
  \widetilde C=
  \{x\in Y_\sigma:
    |x(F_i)-\psi(F_i)|\le\eps/2\text{ for all }i\}.
\]
For a planted embedding $\widehat\theta$, \cref{lem:general-planting-estimate} gives
\[
  |p(F_i,(B_j,\widehat\theta))-p(F_i,(G_t,\theta_t))|
  \le C_m\left(\lambda+\frac1{n-k}\right)<\frac{2\eps}{10}
\]
for every $i$.  Together with the choice of $(G_t,\theta_t)$, this puts every planted
embedding in $\widetilde C$.  Since planted embeddings have probability bounded below by a
constant $c>0$, we have $P_j(\widetilde C)\ge c$ for all $j$.

By Portmanteau, the closedness of $\widetilde C$ gives
\[
  \Ext_\sigma(\phi_0)(\widetilde C)\ge\limsup_j P_j(\widetilde C)\ge c>0.
\]
Because $\Ext_\sigma(\phi_0)$ is concentrated on $X_\sigma$, its support meets
$\widetilde C\cap X_\sigma\subseteq U$.  Thus every basic neighbourhood of $\psi$ meets
the support of some admissible random extension, so $\psi\in S_\sigma$.
\end{proof}

\Cref{thm:blowup-root-plantable} subsumes the two previous root-plantability theorems and
also yields the substitution theorem below, because each closure hypothesis in the next
corollary implies blow-up-closure.

\begin{corollary}\label{cor:closures-imply-blowup}
Each of the following hereditary classes is blow-up-closed, and hence root-plantable at
every non-degenerate type:
\begin{enumerate}
\item every clone-closed class (\cref{thm:clone-root-plantable}): blow up $v$ to an
  edgeless graph $\overline{K_N}$;
\item every true-clone-closed class (\cref{thm:true-clone-root-plantable}): blow up $v$
  to a clique $K_N$;
\item every infinite substitution-closed class (\cref{def:substitution-closed} below):
  blow up $v$ to any in-class graph $H$ with $|H|\ge N$, which exists by infinitude.
\end{enumerate}
\end{corollary}

\begin{proof}
In each case $G[v\to H]$ is the corresponding one-vertex blow-up.  For (1),
$G[v\to\overline{K_N}]$ is an independent blow-up of $G$, so it lies in $\K$ by
clone-closure; for (2), $G[v\to K_N]$ is a complete blow-up, in $\K$ by
true-clone-closure.  For (3), an infinite hereditary class contains graphs of arbitrarily
large order; otherwise, if all orders were bounded, there would be only finitely many
graphs in the class.  Choose such an in-class graph $H$ with $|H|\ge N$.  Then
$G[v\to H]=G[H_w:w\in V(G)]$ with the in-class fibre $H$ at $v$ and single vertices
elsewhere, so it lies in $\K$ by substitution-closure.  Apply
\cref{thm:blowup-root-plantable}.
\end{proof}

For completeness we record the substitution-closure hypothesis named in
\cref{cor:closures-imply-blowup}(3), the strongest of the three.

\begin{definition}[Graph substitution and substitution-closure]\label{def:substitution-closed}
With $G$ and the $H_v$ as above, the substitution $G[H_v:v\in V(G)]$ is the disjoint union
of the $H_v$ together with all edges between $H_v$ and $H_w$ for every $vw\in E(G)$.  A
hereditary class $\K$ is \emph{substitution-closed} if $G[H_v:v\in V(G)]\in\K$ whenever
$G\in\K$ and $H_v\in\K$ for all $v$.
\end{definition}

\begin{theorem}[Substitution-closed classes are root-plantable]\label{thm:substitution-root-plantable}
Every infinite substitution-closed hereditary graph class is blow-up-closed, hence
root-plantable at every non-degenerate type.
\end{theorem}

\begin{proof}
Immediate from \cref{cor:closures-imply-blowup}(3) and
\cref{thm:blowup-root-plantable}.
\end{proof}

\begin{remark}[Substitution-closure is strictly stronger]\label{rem:strictness}
Substitution-closure demands $G[H_v:v\in V(G)]\in\K$ for \emph{all} in-class fibres $H_v$,
whereas the blow-up criterion only needs \emph{one} large blow-up of each base graph to stay in
$\K$.  The gap is real and excludes \cref{thm:clone-root-plantable,thm:true-clone-root-plantable}:
\begin{itemize}
\item Triangle-free graphs are clone-closed but not substitution-closed: $K_2\in\K$, yet
  $K_2[K_2,K_2]=K_4\supseteq K_3$.
\item Cluster graphs (disjoint unions of cliques, equivalently $P_3$-free graphs) are
  true-clone-closed but not substitution-closed: $K_2[\overline{K_2},\overline{K_2}]=C_4$
  has an induced $P_3$.  (Cluster graphs are not clone-closed either, for the same $C_4$;
  cf.\ \cref{cor:cluster-graphs}.)
\end{itemize}
Both classes are nonetheless blow-up-closed --- the first by edgeless interiors, the second
by clique interiors --- so \cref{thm:blowup-root-plantable} covers them while
substitution-closure does not.  Blow-up-closure is exactly the existential ``some
interior works'' weakening of substitution-closure's universal ``every interior works''.
\end{remark}

\section{Finite planting and root-plantable \texorpdfstring{$C_5$}{C5}-free types}

This section is the first place where the choice of type becomes a real part of the
phenomenon.  The closure hypotheses above are type-uniform: once a class is clone-closed,
true-clone-closed, blow-up-closed, or substitution-closed in the relevant sense, the
root-plantability conclusion holds for every non-degenerate type.  Finite planting is
different.  It is a property \emph{at a fixed type} $\sigma$, and a local repair may work
for one type while failing for another.  The $C_5$-free class below is the basic example:
the one-root and two-root non-edge types are plantable, but the two-root edge type is not.

We now isolate a finite, checkable criterion for root-plantability---\emph{finite
planting}---which is implied by the global closure hypotheses of the previous sections but
can also be met by a purely \emph{local} repair when no such closure is available.  Fix a
constrained
labelled limit $\psi\in Q_\sigma$; concretely, $\psi$ is a positive homomorphism of the
constrained $\sigma$-flag algebra, viewed inside the ambient labelled homomorphism space.
By the representation theorem, $\psi$ is the limit of a sequence of finite
$\sigma$-flags
\[
  (G_t,\theta_t),\qquad G_t\in\K,\qquad |G_t|\to\infty,
\]
meaning that for every fixed $\sigma$-flag $F$ one has
$p(F,(G_t,\theta_t))\to\psi(F)$.  These pairs are the concrete finite witnesses in
this section: they are finite graphs with specified roots whose $\sigma$-flag densities approximate the
limiting view $\psi$.

To \emph{plant} such a finite witness means to build a larger graph $H$ in the class with
many embeddings $\widehat\theta:\sigma\to H$ whose $\sigma$-flag densities are close to
those of $(G,\theta)$.  In the proof of \cref{thm:blowup-root-plantable}, this was done
by a global operation: blow up the vertices of $G$ while staying inside the class.  The
roots then have large clone classes, and choosing one vertex from each such
class gives a positive-density set of embeddings $\sigma\to H$.  For these embeddings, the
densities of flags of any fixed size still look almost like they did in $(G,\theta)$.
In other words, a view that was visible only by hand-picking the roots has been made
visible to random roots.

The next examples, especially the $C_5$-free class, do not have such a global closure
property.  A naive blow-up can create forbidden cycles.  What sometimes happens instead
is more local: only a sparse, root-centred part of the graph is altered, with no closure
property of the class assumed.  One first creates many possible embeddings of the
type---in the $C_5$ constructions below, by replacing each root with a positive-density set
of new vertices, any one of which can serve as the image of that root---and then repairs the
few dangerous edges.  If the repair changes so little of the graph that the $\sigma$-flag densities up to the
chosen size are essentially unchanged, these embeddings still reproduce the local view of
$(G,\theta)$.  The purpose of this section is to isolate
exactly what one has to check at the level of finite graphs: the repaired graph must
remain in the class, contain a positive-density set of embeddings $\sigma\to H$, and
preserve the $\sigma$-flag densities up to the chosen size.  The payoff is
\cref{thm:finite-local-planting}: carrying out this finite-graph verification is already
enough to prove root-plantability.

The definition below formalises this finite check.  Let $k=|\sigma|$ be the number of
roots, fix a resolution $m\ge k$, and fix a tolerance $\eps>0$.  We only ask
to preserve $\sigma$-flag densities $p(F,\cdot)$ for flags $F$ with at most $m$ vertices,
and only up to error $\eps$.  The requirement is that every sufficiently large finite
witness $(G,\theta)$ in the class can be replaced by a larger in-class graph $H$ together
with a set $\Theta\subseteq\Emb_\sigma(H)$ of embeddings $\sigma\to H$.  The
lower bound $|\Theta|\ge\delta |V(H)|^k$ ensures that a uniformly random embedding
of $\sigma$ into $H$ lands in $\Theta$ with probability at least $\delta$, since
$|\Emb_\sigma(H)|\le |V(H)|^k$.  For each $\widehat\theta\in\Theta$, the densities
$p(F,(H,\widehat\theta))$ must be close to $p(F,(G,\theta))$ for every such flag $F$.
Thus the finite approximations to $\psi$ can be made typical under random rooting
without assuming any global blow-up closure of the class.

\begin{definition}[Finite planting]\label{def:finite-local-planting}
Let $\K$ be a hereditary graph class and let $\sigma$ be a type in $\K$ with
$k$ roots.  Write $\Emb_\sigma(H)$ for the set of embeddings
$\sigma\to H$.  We say that $\K$ has the \emph{finite planting property at
$\sigma$} if, for every integer $m\ge k$ and every $\eps>0$, there are an integer
$n_0=n_0(m,\eps)\ge m$ and a constant $\delta=\delta(m,\eps)>0$ with the following
property.
Every finite $\sigma$-flag $(G,\theta)$ with $G\in\K$ and $|G|\ge n_0$ admits a graph
$H\in\K$ and a set $\Theta\subseteq\Emb_\sigma(H)$ such that:
\begin{enumerate}[label=\textup{(\roman*)}, leftmargin=2.5em]
  \item $|V(H)|\ge |G|$;
  \item $|\Theta|\ge \delta |V(H)|^k$;
  \item for every $\widehat\theta\in\Theta$ and every ambient $\sigma$-flag $F$
        with $|F|\le m$,
        \[
          \left|p(F,(H,\widehat\theta))-p(F,(G,\theta))\right|<\eps .
        \]
\end{enumerate}
\end{definition}

\begin{theorem}[Finite planting implies root-plantability]\label{thm:finite-local-planting}
Let $\K$ be a hereditary graph class and let $\sigma$ be a non-degenerate type
in $\K$.  If $\K$ has the finite planting property at $\sigma$, then
\[
  S_\sigma=Q_\sigma.
\]
Consequently, quotient semantics and ensemble semantics are equivalent for every
$f\in\A^\sigma[\T_{\mathrm{Graph}}]$.
\end{theorem}

\begin{proof}
The inclusion $S_\sigma\subseteq Q_\sigma$ is the general support-passing
statement.  We prove the reverse inclusion.

Let $\psi\in Q_\sigma$, and let $U$ be a basic neighbourhood of $\psi$ in
$X_\sigma$, specified by ambient $\sigma$-flags $F_1,\ldots,F_s$ and a tolerance
$\eps>0$.  Then
\[
  U=
  \{\chi\in X_\sigma:
    |\chi(F_i)-\psi(F_i)|<\eps\text{ for all }i\}.
\]
Write $k=|\sigma|$ for the number of roots, and pick an integer $m\ge k$
bounding the sizes of the $F_i$.
Choose $n_0=n_0(m,\eps/3)$ and $\delta=\delta(m,\eps/3)$ supplied by the finite
planting property.  Also, choose a convergent sequence of $\sigma$-flags
$(G_t,\theta_t)$ with $G_t\in\K$, $(G_t,\theta_t)\to\psi$, and $|G_t|\to\infty$.
For all sufficiently large $t$, we have $|G_t|\ge n_0$ and
\[
  |p(F_i,(G_t,\theta_t))-\psi(F_i)|<\eps/3
  \qquad(1\le i\le s).
\]

Now apply the finite planting property to each such $(G_t,\theta_t)$, with the integer
$m$ and tolerance $\eps/3$.  Thus, after discarding finitely many terms, we have
graphs $H_t\in\K$ and sets $\Theta_t\subseteq\Emb_\sigma(H_t)$ such that
$|H_t|\ge |G_t|\to\infty$, $|\Theta_t|\ge\delta |V(H_t)|^k$, and every
$\widehat\theta\in\Theta_t$ satisfies
\[
  |p(F_i,(H_t,\widehat\theta))-p(F_i,(G_t,\theta_t))|<\eps/3
  \qquad(1\le i\le s).
\]
Passing to a subsequence, assume the unlabelled graphs $H_t$ converge to some
$\phi_0\in Q_0$.  Since $|\Emb_\sigma(H_t)|\ge |\Theta_t|\ge\delta |V(H_t)|^k$,
the type density of $\sigma$ stays bounded away from zero, so
$\phi_0(\brak\sigma)>0$.

Let $P_t$ be the finite random-rooting distribution obtained by choosing a uniformly
random element of $\Emb_\sigma(H_t)$ and recording all $\sigma$-flag densities in
\[
  Y_\sigma:=[0,1]^{\F^\sigma[\T_{\mathrm{Graph}}]}.
\]
By Razborov's random-extension theorem, $P_t$ converges weakly to
$\Ext_\sigma(\phi_0)$ as a probability measure on the ambient product space.  Let
\[
  \widetilde C=
  \{x\in Y_\sigma:
    |x(F_i)-\psi(F_i)|\le 2\eps/3\text{ for all }i\}.
\]
Then $\widetilde C\cap X_\sigma\subseteq U$.  Every embedding in $\Theta_t$ gives a
point of $\widetilde C$, because its $F_i$-coordinates differ from those of $\psi$ by
at most the two $\eps/3$ errors above.  The total number of embeddings of $\sigma$ into
$H_t$ is at most $|V(H_t)|^k$, so
\[
  P_t(\widetilde C)\ge\delta
\]
for all large $t$.  Since $\widetilde C$ is closed, the closed-set form of
Portmanteau gives
\[
  \Ext_\sigma(\phi_0)(\widetilde C)\ge\delta>0.
\]
The measure $\Ext_\sigma(\phi_0)$ is concentrated on $X_\sigma$, so its support
intersects $\widetilde C\cap X_\sigma\subseteq U$.  Since $U$ was arbitrary,
$\psi\in S_\sigma$.
\end{proof}

\begin{remark}
The clone-based root-plantability theorems above
(\cref{thm:clone-root-plantable,thm:true-clone-root-plantable}) verify the finite
planting property by taking $\Theta$ to be the embeddings obtained by choosing one vertex
from each root clone class.  The one-root and two-root non-edge $C_5$-free constructions
below prove it differently: they first blow up the roots and then delete a sparse set of
dangerous neighbourhood edges.  This is why the criterion is useful beyond closure
properties.
\end{remark}

\subsection{Sparse root-blow-up repairs}

Here is a concrete, checkable way to verify finite planting.  Blow up only
the roots into small positive-density sets from which the embeddings of
$\sigma$ will be chosen, leave all unlabelled vertices as singletons, and repair only a
sparse set of old--old edges.

\begin{definition}[Sparse root-blow-up repair]\label{def:sparse-root-repair}
Let $\K$ be a hereditary graph class and let $\sigma$ be a type in $\K$ with
$k\ge1$ roots.  We say that $\K$ has \emph{sparse root-blow-up
repairs at $\sigma$} if the following holds.
For every $0<\lambda\le1$ and every $\rho>0$, there is an integer
$n_0=n_0(\lambda,\rho)$ such that for every $\sigma$-flag $(G,\theta)$ with $G\in\K$
and $n:=|G|\ge n_0$, there are an integer $L$ with
\[
  \frac{\lambda n}{2}\le L\le \lambda n
\]
and a graph $H\in\K$ with vertex set
\[
  V(H)=U\uplus R_1\uplus\cdots\uplus R_k,
  \qquad
  U:=V(G)\setminus\operatorname{im}\theta,
  \qquad
  |R_i|=L,
\]
such that:
\begin{enumerate}[label=\textup{(\roman*)}, leftmargin=2.5em]
  \item for $i\ne j$, the pair $R_i,R_j$ is complete or empty according as
        $\theta(i)\theta(j)$ is an edge or a non-edge of $G$;
  \item for every $u\in U$ and every $x\in R_i$, the edge $xu$ is present in
        $H$ if and only if $\theta(i)u$ is present in $G$;
  \item the old part $H[U]$ differs from $G[U]$ on at most $\rho n^2$ unordered
        pairs.
\end{enumerate}
\end{definition}

\begin{theorem}[Sparse root repairs imply finite planting]\label{thm:sparse-repair-planting}
Let $\K$ be a hereditary graph class and let $\sigma$ be a non-degenerate type
in $\K$ with $k\ge1$ roots.  If $\K$ has sparse root-blow-up repairs
at $\sigma$, then $\K$ has the finite planting property at $\sigma$.
Consequently, $\K$ is root-plantable at $\sigma$.
\end{theorem}

\begin{proof}
We show that $\K$ has the finite planting property at $\sigma$;
root-plantability then follows from \cref{thm:finite-local-planting}.  Fix
$m\ge k$ and $\eps>0$.  By \cref{def:finite-local-planting} we must, from $m$ and
$\eps$ alone, produce an integer $n_0\ge m$ and a constant $\delta>0$ so that every
$\sigma$-flag $(G,\theta)$ with $G\in\K$ and $|G|\ge n_0$ admits a graph $H\in\K$
and a set $\Theta\subseteq\Emb_\sigma(H)$ meeting conditions (i)--(iii) there.  In
every case $H$ will be a sparse root blow-up repair of $(G,\theta)$ and
$\Theta=R_1\times\cdots\times R_k$ the embeddings choosing one vertex from each root
set; conditions (i) and (ii) are then immediate vertex counts and (iii) a two-event
sampling estimate.  We may therefore fix all constants in advance and only afterwards
introduce $(G,\theta)$.

Fix the constants first.  Set $C_m:=4m^2$.  Choose
$0<\lambda\le1$ and $\rho>0$ so small that
\[
  C_m(\lambda+\rho)<\eps ,
\]
and set
\[
  \delta:=\left(\frac{\lambda}{2(1+k\lambda)}\right)^k>0 .
\]
Let $n_0$ be the largest of $m$, $2k$, $\lceil2/\lambda\rceil$, and the threshold
provided by \cref{def:sparse-root-repair} for this $\lambda$ and $\rho$.  Since
$\lambda$ and $\rho$ are chosen from $m$ and $\eps$ alone (with $k=|\sigma|$ fixed),
both $\delta$ and $n_0$ depend only on $m$ and $\eps$, as
\cref{def:finite-local-planting} requires.

Now let $(G,\theta)$ be a $\sigma$-flag with $G\in\K$ and $n:=|G|\ge n_0$, and let
$H$ be the sparse root blow-up repair supplied by \cref{def:sparse-root-repair} for
this $\lambda$ and $\rho$, so that $V(H)=U\uplus R_1\uplus\cdots\uplus R_k$ with
$U=V(G)\setminus\operatorname{im}\theta$ and each root set of common size
$L\in[\lambda n/2,\lambda n]$.  Put
\[
  \Theta:=R_1\times\cdots\times R_k\subseteq\Emb_\sigma(H),
\]
the inclusion holding because the adjacencies prescribed among the $R_i$, and between
the $R_i$ and $U$, make every choice of one vertex per root set an embedding of
$\sigma$.  We verify conditions (i)--(iii) in turn.

For (i), since $L\ge\lambda n/2\ge1$ (using $n\ge\lceil2/\lambda\rceil$), the vertex
count $|V(H)|=n-k+kL\ge n=|G|$.

For (ii), the same count gives $|V(H)|=n-k+kL\le(1+k\lambda)n$, so, using
$L\ge\lambda n/2$,
\[
  |\Theta|=L^k
  \ge\left(\frac{\lambda n}{2}\right)^k
  \ge\left(\frac{\lambda}{2(1+k\lambda)}\right)^k|V(H)|^k
  =\delta\,|V(H)|^k .
\]

It remains to verify (iii).  Fix an embedding $\widehat\theta\in\Theta$ and a
$\sigma$-flag $F$ with $\ell:=|F|\le m$.  To compute $p(F,(H,\widehat\theta))$,
sample $q:=\ell-k$ additional vertices from
$H\setminus\operatorname{im}\widehat\theta$.  Two events are bad.  First, the sample may meet the union of the
root sets, $R:=R_1\cup\cdots\cup R_k$, of size $|R|=kL$.  The $q$ sampled vertices are
drawn from $W:=V(H)\setminus\operatorname{im}\widehat\theta$, which has
\[
  |W|=|V(H)|-k=n-2k+kL\ge n-k\ge n/2,
\]
the inequalities using $L\ge1$ and $n\ge n_0\ge2k$.  Since $L\le\lambda n$, at most
$kL\le k\lambda n$ of the vertices of $W$ lie in $R$.  Expose the sample as an ordered
draw without replacement $X_1,\dots,X_q$; each $X_j$ is marginally uniform on $W$, so
$\Prob[X_j\in R]\le kL/|W|$.  A union bound over the $q=\ell-k\le m$ positions then gives
\[
  \Prob[\text{the sample meets }R]
  \le q\cdot\frac{kL}{|W|}
  \le m\cdot\frac{k\lambda n}{n/2}
  =2mk\lambda.
\]

Second, the sample may contain both endpoints of one of the altered old pairs.  Let $D$ be
the set of unordered pairs $\{u,u'\}\subseteq U$ on which $H[U]$ and $G[U]$ differ; by
condition (iii) of \cref{def:sparse-root-repair}, $|D|\le\rho n^2$.  If $q\le1$ the sample
has at most one vertex, so it contains no pair, and
\[
  \Prob[\text{the sample contains an unordered pair from }D] = 0 .
\]
Now assume that $q\ge2$.  A fixed pair $\{u,u'\}\subseteq U$ then lies entirely in the sample with probability
\[
  \frac{\binom{|W|-2}{q-2}}{\binom{|W|}{q}}
  =\frac{q(q-1)}{|W|\,(|W|-1)}
  \le\Bigl(\frac{q}{|W|}\Bigr)^{2}
  \le\frac{4q^2}{n^2},
\]
the first inequality holding because $2\le q\le|W|$ (so $|W|-1\ge1$ and
$\tfrac{q-1}{|W|-1}\le\tfrac{q}{|W|}$), the second because $|W|\ge n/2$.  Summing this bound
over the at
most $\rho n^2$ pairs of $D$, we get
\[
  \Prob[\text{the sample contains an unordered pair from }D] \leq 
  \rho n^2\cdot\frac{4q^2}{n^2}=4q^2\rho\le4m^2\rho.
\]

Finally, compare the two sampling experiments for $p(F,(H,\widehat\theta))$ and
$p(F,(G,\theta))$.  Write $B_1$ for the first bad event (the
sample meets $R$) and $B_2$ for the second (the sample contains an altered pair), so that the
two estimates above read $\Prob[B_1]\le2mk\lambda$ and $\Prob[B_2]\le4m^2\rho$.

Realise both experiments on one probability space, as a coupled pair $(S,T)$ of samples---$S$
for the $H$-experiment, $T$ for the $G$-experiment---by the following recipe.  Draw $S$ a
uniformly random $q$-subset of $W$.  If $S\subseteq U$ (equivalently, on $B_1^c$), set
$T:=S$; otherwise (on $B_1$) draw $T$ a uniformly random $q$-subset of $U$, independent of
$S$.  Then $S$ is a uniformly random $q$-subset of $W$ by construction.  Moreover $T$ is a
uniformly random $q$-subset of $U$. Writing $A_H$ for the event
that $S$, together with the roots $\operatorname{im}\widehat\theta$, induces $F$ in $H$, and
$A_G$ for the event that $T$, together with the roots $\operatorname{im}\theta$, induces $F$
in $G$, we thus have $p(F,(H,\widehat\theta))=\Prob[A_H]$ and $p(F,(G,\theta))=\Prob[A_G]$,
while $S=T$ on $B_1^c$.

On $B_1^c\cap B_2^c$ the two samples are equal, lie in $U$, and span no altered pair, so $H$
and $G$ agree on them; together with conditions (i) and (ii) of
\cref{def:sparse-root-repair}---which make $\operatorname{im}\widehat\theta$ induce $\sigma$
and match its adjacencies to $U$ with those of the original roots
$\operatorname{im}\theta$---the $\sigma$-flags induced by $(\operatorname{im}\widehat\theta,S)$
in $H$ and by $(\operatorname{im}\theta,T)$ in $G$ coincide, so $A_H$ and $A_G$ occur
together.  Thus $A_H$ and $A_G$ differ only within $B_1\cup B_2$, and
\[
\begin{aligned}
  \bigl|p(F,(H,\widehat\theta))-p(F,(G,\theta))\bigr|
  &=\bigl|\Prob[A_H]-\Prob[A_G]\bigr|
  \le\Prob[B_1\cup B_2]\\
  &\le\Prob[B_1]+\Prob[B_2]
  \le 2mk\lambda+4m^2\rho .
\end{aligned}
\]
Since $C_m=4m^2$ bounds both $2mk$ (as $k\le m$) and $4m^2$, the right-hand side is at most
$C_m(\lambda+\rho)$, which by the choice of $\lambda$ and $\rho$ is $<\eps$.  This is (iii).

This is the finite planting property at $\sigma$.
\end{proof}

\begin{remark}
\Cref{subsec:c5-planting} below applies this to the $C_5$-free class, where cloning a root forces
deleting only the old edges inside its neighbourhood---a sparse set by \cref{lem:c5-nbhd}.
\end{remark}

\subsection{Root-plantable \texorpdfstring{$C_5$}{C5}-free types}
\label{subsec:c5-planting}

The $C_5$-free graphs are dense and not blow-up-closed, so they are the first natural place to
test root-plantability beyond the global closure theorem.  Indeed, in $K_4$, blowing up one
vertex $v$ to two vertices $x,y$ creates a $5$-cycle $x-a-b-y-c-x$ through the three neighbours
$a,b,c$ of $v$, regardless of the edge between $x$ and $y$, so a naive blow-up leaves the class.
The repair is local, and \cref{thm:sparse-repair-planting} applies: root-plantability holds at the
one-vertex type and at the two-root non-edge type.  (The two-root \emph{edge} type, by contrast,
obstructs; see \cref{sec:c5-edge}.)

The ambient convention is unchanged here: throughout the paper, a hereditary class is ordered by
induced subgraphs.  The local convention is that $C_5$-free has its usual extremal-graph-theory
meaning, namely containing no $C_5$ as an ordinary subgraph.  Equivalently, in the
induced-hereditary language this forbids all five-vertex induced graphs that contain a $C_5$.
Balanced complete bipartite graphs are $C_5$-free and attain edge density $1/2$; conversely, the
Erd\H{o}s--Stone--Simonovits theorem
\cite{ErdosStone,ErdosSimonovits} gives $\mathrm{ex}(n,C_5)=(1/2+o(1))\binom n2$, and Simonovits's
stability/exact theorem for colour-critical graphs \cite{Simonovits1968} identifies the unique
dense extremal limit as the balanced complete bipartite graphon.

The key fact is a strong local constraint.

\begin{lemma}[Neighbourhoods in $C_5$-free graphs]\label{lem:c5-nbhd}
If $G$ is $C_5$-free and $v\in V(G)$, then $G[N(v)]$ contains no path on four
vertices as a subgraph.  Consequently every connected component of $G[N(v)]$ is
a star or a triangle, and in particular
\[
  e(G[N(v)])\le |N(v)|.
\]
\end{lemma}

\begin{proof}
A path $a$-$x$-$y$-$b$ inside $N(v)$ together with $v$ forms the $5$-cycle
$v$-$a$-$x$-$y$-$b$-$v$.  Thus $G[N(v)]$ is $P_4$-subgraph-free.

Now let $C$ be a connected component of a graph with no $P_4$ subgraph.  If $C$
is acyclic, then its diameter is at most $2$, hence it is a star.  If $C$
contains a cycle, no cycle has length at least $4$, and a triangle cannot have
any further vertex in the same component: a shortest path from such a vertex to
the triangle would contain a $P_4$.  Thus $C$ is a triangle.  Each component has
at most as many edges as vertices, giving the displayed bound.
\end{proof}

\subsubsection{The one-root type}

The one-root case is the first example where root-plantability holds without a
global blow-up-closure hypothesis.  Ordinary false-twin duplication of a vertex may create $C_5$s:
if the root sits in a triangle, cloning the root and one neighbour creates a
five-cycle.  The repair is local.  Clone the root into a small positive-density
independent set, but first delete all old edges inside its neighbourhood.  The
deleted set is only linear in the number of vertices by \cref{lem:c5-nbhd}, so
it is invisible to the densities of fixed labelled flags.

\begin{definition}[One-root planting]\label{def:c5-one-root-planting}
Let $(G,r)$ be a rooted graph and let $L\ge1$.  The planted graph
$P_L(G,r)$ is defined as follows.  Remove $r$, add an independent set $R$ of
$L$ new vertices, and keep the old vertex set
\[
  U:=V(G)\setminus\{r\}.
\]
Inside $U$, keep every edge of $G-r$ except the edges with both endpoints in
$N_G(r)$.  Finally, join every vertex of $R$ to every vertex of $N_G(r)$, and
to no other vertex of $U$.
\end{definition}

\begin{lemma}[The planted graph is still $C_5$-free]\label{lem:c5-planting-free}
If $G$ is $C_5$-free, then $P_L(G,r)$ is $C_5$-free for every $L\ge1$.
\end{lemma}

\begin{proof}
Put $H=P_L(G,r)$ and let $R$ be the planted root cluster.  Suppose, for a
contradiction, that $H$ contains a $5$-cycle $C$.

If $C$ uses no vertex of $R$, then all its edges are old edges of $G-r$, so it
is a $C_5$ in $G$.

If $C$ uses exactly one vertex $x\in R$, write the cycle as
\[
  x-a-b-c-d-x.
\]
The edges incident with $x$ force $a,d\in N_G(r)$, and the three middle edges
are old edges of $G$.  Hence
\[
  r-a-b-c-d-r
\]
is a $C_5$ in $G$, again impossible.

Finally suppose $C$ uses at least two vertices of $R$.  Since $R$ is
independent and the independence number of $C_5$ is $2$, it uses exactly two,
say $x,y\in R$.  On one of the two arcs between $x$ and $y$ in the cycle there
are two consecutive old vertices, say $b,c$.  Both are adjacent in $H$ to a
vertex of $R$, so both lie in $N_G(r)$; but the edge $bc$ would then be an edge
inside $N_G(r)$, and all such edges were deleted in the construction of $H$.
This contradiction proves the lemma.
\end{proof}

\begin{lemma}[One-root sparse repairs]\label{lem:c5-one-root-sparse-repair}
The $C_5$-free graph class has sparse root-blow-up repairs at the one-vertex
type.
\end{lemma}

\begin{proof}
Let $0<\lambda\le1$ and $\rho>0$.  Choose
$n_0\ge \max\{2/\lambda,1/\rho\}$.  Let $(G,r)$ be a one-rooted $C_5$-free graph
with $n:=|G|\ge n_0$, and put $L=\lfloor\lambda n\rfloor$.  Then
$\lambda n/2\le L\le\lambda n$.  Let $H=P_L(G,r)$ and let $R$ be the planted
root cluster.  By \cref{lem:c5-planting-free}, $H$ is $C_5$-free.

The sparse-repair conditions are immediate from the construction.  There are no
root--root adjacencies to check for the one-vertex type.  Each vertex of $R$ has
exactly the same neighbours in $U:=V(G)\setminus\{r\}$ as $r$ had in $G$.  Finally,
$H[U]$ differs from $G[U]$ only on the edges inside $N_G(r)$, and by
\cref{lem:c5-nbhd} there are at most $|N_G(r)|\le n\le\rho n^2$ such edges.
\end{proof}

\begin{theorem}[One-root $C_5$-free root-plantability]\label{thm:c5-one-root}
The $C_5$-free graph class is root-plantable at the one-vertex type:
\[
  S_{\vtype}=Q_{\vtype}.
\]
\end{theorem}

\begin{proof}
The one-vertex type is non-degenerate because the $C_5$-free class contains
arbitrarily large graphs.  Hence
\cref{lem:c5-one-root-sparse-repair,thm:sparse-repair-planting} apply.
\end{proof}

\begin{remark}
\Cref{thm:c5-one-root} is the promised local replacement for the global
blow-up-closure theorem in one rooted variable.  It uses two facts special to one root:
the only dangerous new $C_5$s created by cloning the root pass through an old
edge inside $N_G(r)$, and those old neighbourhood edges have density
$O(1/|G|)$ in sampling for fixed labelled flags.  For larger types one would have to control
edges inside and between several root-neighbourhood regions.  The next result
shows that this can still be done for two non-adjacent roots; \cref{sec:c5-edge}
shows that it cannot work in general at the edge type.
\end{remark}

\subsubsection{The two-root non-edge type}

Let $\eta$ be the two-vertex non-edge type.  The same idea works for $\eta$:
clone both roots into two independent clusters, leave no edges between the
clusters, and delete the old edges inside each root-neighbourhood.  Mixed
five-cycles using one clone of each root project back to genuine $C_5$s in the
original graph; cycles using two clones from the same root are killed by the
neighbourhood-edge deletion.

\begin{definition}[Two-root non-edge planting]\label{def:c5-nonedge-planting}
Let $(G,r,s)$ be a rooted graph with $rs\notin E(G)$, and let $L\ge1$.  The
planted graph $P_L(G,r,s)$ is defined as follows.  Remove $r,s$, add two
disjoint independent sets $R$ and $S$ of size $L$, and keep the old vertex set
\[
  U:=V(G)\setminus\{r,s\}.
\]
There are no edges inside $R$, inside $S$, or between $R$ and $S$.  Inside
$U$, keep every edge of $G-\{r,s\}$ except the edges with both endpoints in
$N_G(r)$ or both endpoints in $N_G(s)$.  Finally, join every vertex of $R$ to
every vertex of $N_G(r)$, and every vertex of $S$ to every vertex of $N_G(s)$,
with no other new edges.
\end{definition}

\begin{lemma}[The non-edge planted graph is $C_5$-free]\label{lem:c5-nonedge-planting-free}
If $G$ is $C_5$-free and $rs\notin E(G)$, then $P_L(G,r,s)$ is $C_5$-free for
every $L\ge1$.
\end{lemma}

\begin{proof}
Put $H=P_L(G,r,s)$ and let $R,S$ be the planted clusters.  The set
$R\cup S$ is independent, and the independence number of $C_5$ is $2$, so any
$5$-cycle in $H$ uses at most two planted vertices.

If a $5$-cycle uses no planted vertex, all its edges are old edges of
$G-\{r,s\}$, hence it is a $C_5$ in $G$.

If it uses exactly one planted vertex, say $x\in R$, write the cycle as
\[
  x-a-b-c-d-x.
\]
Then $a,d\in N_G(r)$ and the middle edges are old edges of $G$, so
$r-a-b-c-d-r$ is a $C_5$ in $G$.  The case $x\in S$ is identical.

If it uses two planted vertices from the same cluster, say $x,y\in R$, then on
the length-three arc between $x$ and $y$ there are two consecutive old vertices
$b,c$.  Both are adjacent to a vertex of $R$, hence both lie in $N_G(r)$, but
the edge $bc$ would have been deleted.  The same argument applies to two
vertices of $S$.

It remains to consider one planted vertex $x\in R$ and one planted vertex
$y\in S$.  Since $x$ and $y$ are non-adjacent, replacing $x$ by $r$ and $y$ by
$s$ in the cycle preserves all five cycle edges: neighbours of $x$ lie in
$N_G(r)$, neighbours of $y$ lie in $N_G(s)$, and all edges among the three old
vertices are old edges of $G$.  The replacement therefore gives a $C_5$ in
$G$, contradicting the hypothesis.
\end{proof}

\begin{lemma}[Non-edge sparse repairs]\label{lem:c5-nonedge-sparse-repair}
The $C_5$-free graph class has sparse root-blow-up repairs at the two-root
non-edge type $\eta$.
\end{lemma}

\begin{proof}
Let $0<\lambda\le1$ and $\rho>0$.  Choose
$n_0\ge \max\{2/\lambda,2/\rho\}$.  Let $(G,r,s)$ be a $C_5$-free graph rooted at
a non-edge, with $n:=|G|\ge n_0$, and put $L=\lfloor\lambda n\rfloor$.  Then
$\lambda n/2\le L\le\lambda n$.  Let $H=P_L(G,r,s)$ with planted clusters
$R,S$.  By \cref{lem:c5-nonedge-planting-free}, $H$ is $C_5$-free.

There are no edges between $R$ and $S$, as required for the non-edge type, and the
adjacencies from $R$ and $S$ to $U:=V(G)\setminus\{r,s\}$ agree with the adjacencies from
$r$ and $s$ in $G$.  The old part $H[U]$ differs from $G[U]$ only on edges inside
$N_G(r)$ or inside $N_G(s)$.  By \cref{lem:c5-nbhd}, the number of such edges is
at most
\[
  |N_G(r)|+|N_G(s)|\le 2n\le \rho n^2 .
\]
Thus the sparse root-blow-up repair condition holds.
\end{proof}

\begin{theorem}[Non-edge type $C_5$-free root-plantability]\label{thm:c5-nonedge-root}
The $C_5$-free graph class is root-plantable at the two-root non-edge type
$\eta$:
\[
  S_\eta=Q_\eta.
\]
\end{theorem}

\begin{proof}
The non-edge type $\eta$ is non-degenerate because the $C_5$-free class contains
arbitrarily large empty graphs.  Hence
\cref{lem:c5-nonedge-sparse-repair,thm:sparse-repair-planting} apply.
\end{proof}

For comparison, ordinary independent blow-ups detect exactly the triangles.

\begin{lemma}\label{lem:c5-blowup}
Every independent blow-up of a $C_5$-free graph $G$ is again $C_5$-free if and only
if $G$ is triangle-free.
\end{lemma}

\begin{proof}
A $C_5$ in a blow-up projects to a closed walk of length $5$ in $G$, and conversely
any closed walk of length $5$ in $G$ lifts to a $C_5$ in a blow-up with classes of
size $\ge2$.  Since $G$ is $C_5$-free, such a walk---being of odd length $5$, hence
containing an odd cycle of length $3$ or $5$---exists if and only if $G$ has a
triangle.
\end{proof}

\section{Pinning obstructions to root-plantability}

We now turn from planting criteria to obstructions, beginning with the simplest
pinning mechanism: \emph{degeneracy}.  Here pinning is a property of a
particular labelled flag, or more generally of a particular labelled algebra
element.  It means that, for some type
\(\sigma\), element \(a\in\A^\sigma[\T_{\mathrm{Graph}}]\), and constant \(c\),
when any constrained unlabelled limit is rooted at random as a \(\sigma\)-flag,
the resulting homomorphism evaluates \(a\) as \(c\) almost surely,
even though a specially chosen labelled limit can evaluate \(a\) differently.
Degeneracy gives the most basic example.  It means
that graphs in the class have vanishing edge density asymptotically: every
$n$-vertex member has $o(n^2)$ edges, equivalently every dense limit of the
class has edge density zero.  In such a limit a typical vertex sees no
edge from the root to a further vertex, so the density of the one-root edge flag is pinned to \(0\)
under random rooting.  The class may still contain finite graphs with an
\emph{exceptional} high-degree vertex, such as the centre of a star.  The one-root edge flag
separates the typical rooted view from this exceptional rooted view.

For the degeneracy obstruction below, write \(\vtype\) for the one-vertex
type.  The notation \(\one_\sigma\) always denotes the unit of the
\(\sigma\)-typed flag algebra; thus \(\vtype\) is a type, while
\(\one_{\vtype}\) is an algebra element.  Let
$e\in\A^{\vtype}[\T_{\mathrm{Graph}}]$ be the edge flag over the one-vertex type (the
two-vertex flag in which the root is adjacent to the one unlabelled vertex), and
$\rho\in\A^0[\T_{\mathrm{Graph}}]$ is
the unlabelled edge flag; recall $\avg{e}{\vtype}=\rho$ and
$\brak{\vtype}=\one_0$.

\begin{definition}[Edge-degenerate class]\label{def:edge-degenerate}
A hereditary graph class $\K$ is \emph{edge-degenerate} if every constrained
unlabelled limit has edge density zero, that is, $\phi_0(\rho)=0$ for all
$\phi_0\in Q_0$.  Equivalently, the maximum number of edges of an $n$-vertex graph
in $\K$ is $o(n^2)$.
\end{definition}

\begin{theorem}[Degeneracy obstruction]\label{thm:degenerate-obstruction}
Let $\K$ be a hereditary graph class that is edge-degenerate and contains
arbitrarily large stars.  Then $(\T_{\mathrm{Graph}},\T_\K,\sigma{=}\vtype)$ is not
root-plantable.  Concretely, for $f:=-e\in\A^{\vtype}[\T_{\mathrm{Graph}}]$,
\[
  \Prob_{\phi^{\vtype}\sim\Ext_{\vtype}(\phi_0)}[\phi^{\vtype}(f)\ge0]=1
  \qquad\text{for every }\phi_0\in Q_0,
\]
but some $\psi\in Q_{\vtype}$ has $\psi(f)<0$; in fact
\[
  S_{\vtype}\ \subseteq\ \{\chi\in Q_{\vtype}:\chi(e)=0\}\ \subsetneq\ Q_{\vtype}.
\]
\end{theorem}

\begin{proof}
Let $\phi_0\in Q_0$.  Since $\brak{\vtype}=\one_0$ we
have $\phi_0(\brak{\vtype})=1>0$, so $\Ext_{\vtype}(\phi_0)$ is defined.  As $\avg{e}{\vtype}=\rho$,
\eqref{eq:extension-expectation} gives
\[
  \E_{\phi^{\vtype}\sim\Ext_{\vtype}(\phi_0)}[\phi^{\vtype}(e)]
  =\frac{\phi_0(\rho)}{\phi_0(\brak{\vtype})}=\phi_0(\rho)=0,
\]
using edge-degeneracy.  Because $\phi^{\vtype}(e)\ge0$ pointwise, $\phi^{\vtype}(e)=0$ almost
surely, hence $\phi^{\vtype}(-e)=0\ge0$ almost surely.  Moreover,
by \cref{lem:support-as} every support point $\chi$ of every $\Ext_{\vtype}(\phi_0)$
satisfies $\chi(e)=0$, and since $\chi\mapsto\chi(e)$ is continuous this gives the
inclusion $S_{\vtype}\subseteq\{\chi\in Q_{\vtype}:\chi(e)=0\}$.

Let $S_N=K_{1,N-1}$ be the star on $N$
vertices, rooted at its centre; by hypothesis $S_N\in\K$ for arbitrarily large $N$.
Every unlabelled vertex is adjacent to the centre, so $p(e,S_N)=1$ for all $N$.  By
compactness a subsequence converges to some $\psi\in Q_{\vtype}$ with $\psi(e)=1$, whence
$\psi(-e)=-1<0$.  Thus, $\psi\in Q_{\vtype}$ witnesses
$\{\chi\in Q_{\vtype}:\chi(e)=0\}\subsetneq Q_{\vtype}$; in particular
$\psi\in Q_{\vtype}\setminus S_{\vtype}$, so $S_{\vtype}\ne Q_{\vtype}$.
\end{proof}

\begin{remark}
The hypothesis ``contains arbitrarily large stars'' is only used to exhibit, in
$Q_{\vtype}$, a labelled limit of positive one-root edge density; it may be replaced by the
existence of any sequence of one-root graphs in $\K$ whose one-root edge density stays
bounded away from $0$.
\end{remark}

\begin{remark}[The obstruction]
The centre of a star is an exceptional root.  Quotient semantics can see
it, because labelled limits may sit at exceptional vertices.  Ensemble semantics
instead chooses the root in a constrained unlabelled limit at a \emph{typical} vertex; in a
degenerate class the edge density is zero, so a typical vertex has one-root edge
density zero.  The high-degree vertex cannot be planted with positive probability
without leaving the class: duplicating the centre and one neighbour already produces
a $K_{2,2}$, and more generally a denser local configuration than a degenerate class can
sustain.
\end{remark}

\subsection{Sparse degenerate examples}

The hypotheses of \cref{thm:degenerate-obstruction} are met by many familiar
sparse classes.  We first record the original $C_4$-free case, where edge-degeneracy
is the classical K\H{o}v\'ari--S\'os--Tur\'an bound.

Let
\[
  \K_{C_4}:=\{G:\text{$G$ contains no copy of $C_4$ as an ordinary subgraph}\}.
\]
This is hereditary; in the induced-subgraph language of flag algebras it is
axiomatised by forbidding the three four-vertex graphs that contain a $C_4$ as a
subgraph: $C_4$, $K_4$ minus an edge, and $K_4$.  It is not clone-closed: the
independent blow-up of $K_2$ into two clone classes of size $2$ is $K_{2,2}=C_4$.

\begin{lemma}[Dense \texorpdfstring{$C_4$}{C4}-free limits have edge density zero]\label{lem:c4-edge-zero}
$\K_{C_4}$ is edge-degenerate: if $\phi_0\in Q_0$ for the $C_4$-free class, then
$\phi_0(\rho)=0$.
\end{lemma}

\begin{proof}
By the representation theorem in the constrained class, $\phi_0$ is the limit
of a convergent sequence of finite $C_4$-free graphs $G_n$ with $|G_n|\to\infty$.
We show that their edge densities tend to zero.

Let $G$ be a $C_4$-free graph on $N$ vertices, with degrees $d(v)$.  Any two
vertices have at most one common neighbour; otherwise two common neighbours and
the two original vertices form a $C_4$.  Therefore
\[
  \sum_{v\in V(G)} \binom{d(v)}{2}
  \le
  \binom{N}{2}.
\]
By convexity,
\[
  N\binom{\overline d}{2}
  \le
  \binom{N}{2},
  \qquad
  \overline d:=\frac1N\sum_v d(v).
\]
Thus $\overline d(\overline d-1)\le N-1$, so
$\overline d\le 1+\sqrt N$.  The edge density is
$\overline d/(N-1)$, which tends to zero.  Hence
$\phi_0(\rho)=0$.
\end{proof}

\begin{corollary}[\texorpdfstring{$C_4$}{C4}-free counterexample]\label{cor:c4-counterexample}
$\K_{C_4}$ is edge-degenerate (\cref{lem:c4-edge-zero}) and contains all stars,
so by \cref{thm:degenerate-obstruction} it is not root-plantable at
$\sigma=\vtype$: the element $-e$ is non-negative almost surely under every admissible
one-root random extension, but is negative at some point of $Q_{\vtype}$.
\end{corollary}

\begin{corollary}[A family of degenerate obstructions]\label{cor:degenerate-family}
Each of the following hereditary classes is edge-degenerate and contains all stars,
hence is not root-plantable at the one-vertex type:
\begin{enumerate}[label=\textup{(\roman*)}, leftmargin=2.5em]
  \item $K_{s,t}$-subgraph-free graphs for $2\le s\le t$, edge-degenerate by the
        K\H{o}v\'ari--S\'os--Tur\'an bound $\mathrm{ex}(n,K_{s,t})=O(n^{2-1/s})$
        \cite{KovariSosTuran} (a star contains no $K_{2,2}$);
  \item $C_{2k}$-subgraph-free graphs for $k\ge2$, edge-degenerate by the
        Bondy--Simonovits bound $\mathrm{ex}(n,C_{2k})=O(n^{1+1/k})$
        \cite{BondySimonovits};
  \item forests, and more generally any class of bounded average degree containing
        all stars ($O(n)$ edges);
  \item planar graphs ($\le 3n-6$ edges).
\end{enumerate}
\end{corollary}

\subsection{Complementation: a dense obstruction}

One might hope that the failures above are an artefact of \emph{sparsity}, and that
density rescues root-plantability.  It does not.  Root-plantability is preserved by
complementation, which immediately turns each sparse obstruction into a dense one.

\begin{lemma}[Complementation invariance]\label{lem:complementation}
For a hereditary graph class $\K$ let $\overline\K=\{\overline G:G\in\K\}$, which is
again hereditary, and for a type $\sigma$ let $\overline\sigma$ be its complement.
Then $(\T_{\mathrm{Graph}},\T_\K,\sigma)$ is root-plantable if and only if
$(\T_{\mathrm{Graph}},\T_{\overline\K},\overline\sigma)$ is.  For the one-vertex type
$\sigma=\vtype=\overline\sigma$, the class $\K$ is root-plantable at \(\vtype\) if and only if
$\overline\K$ is.
\end{lemma}

\begin{proof}
In this proof write $Q_\sigma(\K)$ and $S_\sigma(\K)$ for the two spaces
defined using the constrained class $\K$.
Let $C_\sigma$ send a finite $\sigma$-flag $F=(G,\theta)$ to the
$\overline\sigma$-flag $\overline F=(\overline G,\theta)$.  This operation is
an involution: applying the corresponding map $C_{\overline\sigma}$ returns
the original flag.  This rule clearly sends isomorphic flags to isomorphic
flags.  We next check that it is compatible with the extension equalities used
to define the flag algebra.  In the flag algebra, a flag is identified with
the weighted average of its one-vertex extensions; therefore, if complementing
a flag is to define a map on algebra elements, it must carry each such equality
to the corresponding equality for complemented flags.  If $F$ is a
$\sigma$-flag and $H$ is a larger $\sigma$-flag, then the same chosen set of
unlabelled vertices, together with the roots, induces $F$ in $H$ if and only if
it induces $\overline F$ in $\overline H$.  Therefore
\[
  p(F,H)=p(\overline F,\overline H).
\]
In particular, if
\[
  F-\sum_{|H|=|F|+1}p(F,H)\,H
\]
is one of the one-vertex extension relations defining
$\A^\sigma[\T_{\mathrm{Graph}}]$, its image under $C_\sigma$ is
\[
  \overline F-\sum_{|H|=|F|+1}p(\overline F,\overline H)\,\overline H,
\]
which is exactly the corresponding extension relation in
$\A^{\overline\sigma}[\T_{\mathrm{Graph}}]$.  Thus two expressions that
represent the same element before complementing still represent the same
element after complementing.  So $C_\sigma$ gives a well-defined linear map
between the two flag algebras.  Since complementation is involutive, this
linear map has inverse $C_{\overline\sigma}$, and hence is a linear isomorphism
\[
  C_\sigma:\A^\sigma[\T_{\mathrm{Graph}}]\cong
  \A^{\overline\sigma}[\T_{\mathrm{Graph}}].
\]

We next check that this linear isomorphism is multiplicative.  Let $F_1$ and
$F_2$ be $\sigma$-flags, and let $H$ be a $\sigma$-flag of the product size
$|F_1|+|F_2|-|\sigma|$.  By definition, the coefficient of $H$ in
$F_1\cdot F_2$ is the probability that a uniformly random split of the unlabelled
vertices of $H$ into two parts of the required sizes exhibits $F_1$ on the
first part together with the roots and $F_2$ on the second part together with
the roots.  For each fixed split this event holds in $H$ if and only if the
same split in $\overline H$ exhibits $\overline F_1$ and $\overline F_2$.
Thus the coefficient of $H$ in $F_1\cdot F_2$ is the coefficient of $\overline H$ in
$\overline F_1\cdot\overline F_2$.  Summing over all product-size flags gives
\[
  C_\sigma(F_1\cdot F_2)=C_\sigma(F_1)\cdot C_\sigma(F_2).
\]
By bilinearity the same identity holds for arbitrary elements of
$\A^\sigma[\T_{\mathrm{Graph}}]$, so $C_\sigma$ is an algebra isomorphism.

Now consider the constrained classes.  A $\sigma$-flag $F$ is admissible for
$\K$ exactly when its underlying unlabelled graph lies in $\K$.  The underlying
graph of $C_\sigma(F)$ is the complement of that graph, so $F$ is admissible
for $\K$ exactly when $C_\sigma(F)$ is admissible for $\overline\K$; likewise,
forbidden flags correspond to forbidden flags.  Hence the algebra isomorphism
identifies the constrained quotient for $\K$ with the constrained quotient for
$\overline\K$.

Let
\[
  \widehat C_\sigma:X_\sigma\longrightarrow X_{\overline\sigma}
\]
be the induced map on positive homomorphisms, defined by
\[
  \widehat C_\sigma(\chi)(a)=\chi(C_{\overline\sigma}a)
  \qquad
  (a\in\A^{\overline\sigma}[\T_{\mathrm{Graph}}]).
\]
This is precomposition with the inverse algebra isomorphism.  We check
continuity directly from the definition of the topology.  For any type $\tau$
and any $a\in\A^\tau[\T_{\mathrm{Graph}}]$, write
\[
  \operatorname{ev}_a:X_\tau\to\mathbb R,
  \qquad
  \operatorname{ev}_a(\eta)=\eta(a),
\]
for the corresponding evaluation map.  The topology on $X_\tau$ is the weakest
topology in which all such maps are continuous.  Thus it suffices to check
that each composite ${\operatorname{ev}_a}\circ{\widehat C_\sigma}$ is
continuous.  For $a\in\A^{\overline\sigma}[\T_{\mathrm{Graph}}]$,
\[
  ({\operatorname{ev}_a} \circ {\widehat C_\sigma})(\chi)
  =
  \widehat C_\sigma(\chi)(a)
  =
  \chi(C_{\overline\sigma}a)
  =
  \operatorname{ev}_{C_{\overline\sigma}a}(\chi),
\]
which is one of the defining continuous evaluation maps on $X_\sigma$.  Thus
$\widehat C_\sigma$ is continuous.  The same argument with $\sigma$ and
$\overline\sigma$ interchanged shows that $\widehat C_{\overline\sigma}$ is
continuous, and the two maps are inverse to each other.  Thus
$\widehat C_\sigma$ is a continuous bijection with continuous inverse.  Because
forbidden flags correspond under complementation, the intrinsic description of
$Q_\sigma$ as the homomorphisms vanishing on forbidden flags shows that
\[
  \widehat C_\sigma(Q_\sigma(\K))=Q_{\overline\sigma}(\overline\K).
\]

It remains to check the ensemble supports.  Let $\phi_0\in Q_0(\K)$ and put
$\overline\phi_0:=\widehat C_\emptyset(\phi_0)\in Q_0(\overline\K)$.  Since
complementation preserves embeddings of the type, it also commutes with
unlabelling:
\[
  C_\emptyset(\avg{g}{\sigma})
  =
  \avg{C_\sigma(g)}{\overline\sigma},
  \qquad
  C_\emptyset(\brak\sigma)=\brak{\overline\sigma}.
\]
Therefore $\phi_0(\brak\sigma)>0$ if and only if
$\overline\phi_0(\brak{\overline\sigma})>0$.  Moreover, for
$g\in\A^\sigma[\T_{\mathrm{Graph}}]$ the expectation formula gives
\[
  \E_{\psi\sim\Ext_\sigma(\phi_0)}[\psi(g)]
  =
  \frac{\phi_0(\avg{g}{\sigma})}{\phi_0(\brak\sigma)}
  =
  \frac{\overline\phi_0(\avg{C_\sigma(g)}{\overline\sigma})}
       {\overline\phi_0(\brak{\overline\sigma})}.
\]
The right-hand side is the expectation of $C_\sigma(g)$ under
$\Ext_{\overline\sigma}(\overline\phi_0)$.  Let
\[
  \mu:=(\widehat C_\sigma)_*\Ext_\sigma(\phi_0),
  \qquad
  \mu':=\Ext_{\overline\sigma}(\overline\phi_0).
\]
We show that $\mu=\mu'$.  By the definition of pushforward and the identity
\(\widehat C_\sigma(\psi)(C_\sigma g)=\psi(g)\), the preceding expectation
calculation gives, for every $g\in\A^\sigma[\T_{\mathrm{Graph}}]$,
\[
  \int_{X_{\overline\sigma}}\eta(C_\sigma g)\,d\mu(\eta)
  =
  \E_{\eta\sim\Ext_{\overline\sigma}(\overline\phi_0)}[\eta(C_\sigma g)]
  =
  \int_{X_{\overline\sigma}}\eta(C_\sigma g)\,d\mu'(\eta).
\]
Since $C_\sigma$ is onto, every element of
$\A^{\overline\sigma}[\T_{\mathrm{Graph}}]$ has the form $C_\sigma g$.  Hence
$\mu$ and $\mu'$ have the same integral against every evaluation function
$\operatorname{ev}_a$ on $X_{\overline\sigma}$.

It remains to explain why agreement on all evaluation functions determines the
measure.  Let
\[
  \mathcal E
  :=
  \{\operatorname{ev}_a:a\in\A^{\overline\sigma}[\T_{\mathrm{Graph}}]\}.
\]
This class already contains all finite linear combinations of its elements:
\[
  \sum_i c_i\,\operatorname{ev}_{a_i}
  =
  \operatorname{ev}_{\sum_i c_i a_i}.
\]
Thus it is closed under sums and scalar multiples.  It contains the constant
function $1$ because $\operatorname{ev}_{\one_{\overline\sigma}}=1$, and it is
closed under products because
\[
  \operatorname{ev}_a(\eta)\operatorname{ev}_b(\eta)
  =
  \eta(a)\eta(b)
  =
  \eta(a\cdot b)
  =
  \operatorname{ev}_{a\cdot b}(\eta).
\]
It also separates points: if two homomorphisms $\eta_1,\eta_2$ are different,
then by definition they differ on some flag-algebra element $a$, so
\(\operatorname{ev}_a(\eta_1)\ne\operatorname{ev}_a(\eta_2)\).  The
Stone--Weierstrass theorem therefore says the following concrete approximation
statement: for every continuous function
\(h:X_{\overline\sigma}\to\mathbb R\) and every \(\eps>0\), there is some
\(\operatorname{ev}_a\in\mathcal E\) such that
\[
  \sup_{\eta\in X_{\overline\sigma}} |h(\eta)-\operatorname{ev}_a(\eta)|<\eps .
\]
This is what ``uniformly dense'' means here.  Since $\mu$ and $\mu'$ agree on
each evaluation function, they agree on this approximating
\(\operatorname{ev}_a\).  Hence
\[
  \left|\int h\,d\mu-\int h\,d\mu'\right|
  \le
  \left|\int (h-\operatorname{ev}_a)\,d\mu\right|
  +
  \left|\int (h-\operatorname{ev}_a)\,d\mu'\right|
  <2\eps,
\]
because both $\mu$ and $\mu'$ are probability measures.  Since $\eps$ is
arbitrary, the two integrals of $h$ are equal.  Thus $\mu$ and $\mu'$ agree on
all continuous functions, and probability measures on a compact space are
determined by their integrals against continuous functions.  Hence
\[
  (\widehat C_\sigma)_*\Ext_\sigma(\phi_0)
  =
  \Ext_{\overline\sigma}(\overline\phi_0).
\]

As $\phi_0$ ranges over $Q_0(\K)$, $\overline\phi_0$ ranges over
$Q_0(\overline\K)$.  It remains only to translate supports.  We use the following
elementary fact.  If $T:X\to Y$ is a continuous bijection with continuous inverse
and $\nu=T_*\mu$, then
\[
  T(\supp\mu)=\supp\nu .
\]
Indeed, if $x\in\supp\mu$ and $V$ is an open neighbourhood of $T(x)$, then
$T^{-1}(V)$ is an open neighbourhood of $x$, so
\(\nu(V)=\mu(T^{-1}(V))>0\).  Thus $T(x)\in\supp\nu$.  Conversely, if
$T(x)\in\supp\nu$ and $U$ is an open neighbourhood of $x$, then $T(U)$ is an
open neighbourhood of $T(x)$; by injectivity,
\(\mu(U)=\nu(T(U))>0\).  Thus $x\in\supp\mu$.  Applying this fact to
$T=\widehat C_\sigma$ and to the equality of extension measures above,
$\widehat C_\sigma$ takes each support appearing in the definition of
$S_\sigma(\K)$ to the corresponding support appearing in the definition of
$S_{\overline\sigma}(\overline\K)$.  After taking the union and closure, we get
\[
  \widehat C_\sigma(S_\sigma(\K))=S_{\overline\sigma}(\overline\K).
\]

Thus the same explicit map $\widehat C_\sigma$ sends the quotient space
$Q_\sigma(\K)$ exactly onto $Q_{\overline\sigma}(\overline\K)$, and sends the
ensemble-support space $S_\sigma(\K)$ exactly onto
$S_{\overline\sigma}(\overline\K)$.  Consequently
$S_\sigma(\K)=Q_\sigma(\K)$ if and only if
$S_{\overline\sigma}(\overline\K)=Q_{\overline\sigma}(\overline\K)$.  This is
the desired equivalence of root-plantability.  The final assertion follows
because the one-vertex type is equal to its complement.
\end{proof}

Call $\K$ \emph{co-edge-degenerate} if $\overline\K$ is edge-degenerate; equivalently
every $\phi_0\in Q_0$ has $\phi_0(\rho)=1$, equivalently every $n$-vertex graph in
$\K$ has at least $\binom n2-o(n^2)$ edges.  Such classes are as dense as possible,
with edge density tending to $1$, yet they fail just as the sparse ones do.

\begin{corollary}[A dense obstruction]\label{cor:codegenerate}
Let $\K$ be hereditary and co-edge-degenerate, and suppose that $\K$ contains
the co-stars
\[
  \overline{K_{1,N-1}}=K_{N-1}\uplus K_1
\]
for arbitrarily large $N$.  Then $\K$ is not root-plantable at
$\sigma=\vtype$.  More concretely, the element
$e-\one_{\vtype}\in\A^{\vtype}[\T_{\mathrm{Graph}}]$ is non-negative almost
surely under every admissible one-root random extension, but is negative at
some point of $Q_{\vtype}(\K)$.

In particular, the graph-complement classes of the examples in
\cref{cor:degenerate-family} are dense hereditary classes that are not
root-plantable at the one-vertex type.
\end{corollary}

\begin{proof}
Since $\K$ is co-edge-degenerate, $\overline\K$ is edge-degenerate.  Since
$\K$ contains arbitrarily large co-stars, $\overline\K$ contains arbitrarily
large stars.  Thus \cref{thm:degenerate-obstruction} shows that
$\overline\K$ is not root-plantable at \(\vtype\).  The one-vertex type is
equal to its complement, so \cref{lem:complementation} transfers this failure
back to $\K$.

Let us also identify the witness directly.  If $\phi_0\in Q_0(\K)$, then
co-edge-degeneracy gives $\phi_0(\rho)=1$.  Since
\(\avg{e}{\vtype}=\rho\) and \(\brak{\vtype}=\one_0\), the extension formula
gives
\[
  \E_{\phi^{\vtype}\sim\Ext_{\vtype}(\phi_0)}
  [\phi^{\vtype}(e)]
  =
  \phi_0(\rho)
  =
  1.
\]
The random variable \(\phi^{\vtype}(e)\) is a flag density, so it lies in
\([0,1]\).  Hence \(\phi^{\vtype}(e)=1\) almost surely, and therefore
\[
  \phi^{\vtype}(e-\one_{\vtype})=0
  \qquad\text{almost surely.}
\]
Thus \(e-\one_{\vtype}\) is ensemble-non-negative.

On the other hand, choose co-stars
\(\overline{K_{1,N_j-1}}=K_{N_j-1}\uplus K_1\in\K\) with \(N_j\to\infty\),
and root each one at its isolated vertex.  The one-root edge density is then
always \(0\).  By compactness, a subsequence converges to a point
\(\psi\in Q_{\vtype}(\K)\) with \(\psi(e)=0\), and hence
\[
  \psi(e-\one_{\vtype})=-1<0.
\]
So quotient semantics fails for the same element.

For the final claim, take the complements of the classes listed in
\cref{cor:degenerate-family}.  The original classes are edge-degenerate and
contain arbitrarily large stars; their complement classes are therefore
co-edge-degenerate and contain arbitrarily large co-stars, so the first part
applies.
\end{proof}

\begin{remark}[The dichotomy: density is not the dividing line]
\cref{thm:degenerate-obstruction} and \cref{cor:codegenerate}
together show that root-plantability is governed neither by heredity nor by density:
edge-degenerate classes (one-root edge density pinned to $0$) and co-edge-degenerate
classes (pinned to $1$) both fail, for the same reason.  What separates a typical
vertex from the exceptional one is that the one-root edge density is \emph{pinned} to a
boundary value of $[0,1]$ across all constrained unlabelled limits, while the
constrained class still contains a labelled limit attaining the opposite boundary.  The
governing invariant is this boundary pinning, not sparsity or density.

This refines, rather than closes, the problem of characterising root-plantability.
The $C_5$-free graphs show why the question must be asked at a fixed type.  At the
one-vertex type the edge flag is \emph{not} pinned: empty graphs give one-root edge
density $0$, while balanced complete bipartite graphs give one-root edge density
$\tfrac12$.  Thus no boundary-pinning obstruction comes from the one-root edge flag, and
a local root-blow-up repair proves root-plantability there.  The same repair also handles
the two-root non-edge type by cloning the two non-adjacent roots simultaneously.  At the
two-root edge type, by contrast, a pinning obstruction does arise, so the full all-types
$C_5$-free conjecture is false (\cref{sec:c5-edge}), even though the positive one-root and
non-edge cases remain useful models for planting beyond clone-closure.
\end{remark}

\subsection{The general obstruction}

Both obstructions are instances of one mechanism: the density of a flag pinned almost
surely to a single value across all constrained limits, while the quotient escapes
it.

\begin{theorem}[Pinning obstruction]\label{thm:pinning}
Let $\sigma$ be a non-degenerate type, let $F$ be a $\sigma$-flag, and let
$c\in[0,1]$.  Suppose the density of $F$ is almost surely pinned to $c$ across
all constrained unlabelled limits: for every $\phi_0\in Q_0$ with
$\phi_0(\brak\sigma)>0$, one has $\phi^\sigma(F)=c$ almost surely under
$\phi^\sigma \sim \Ext_\sigma(\phi_0)$.  If the constrained class contains a labelled limit attaining a different
value---some $\psi\in Q_\sigma$ with $\psi(F)\ne c$---then $(\T_0,\T_1,\sigma)$ is not
root-plantable.
\end{theorem}

\begin{proof}
Let
\[
  Z:=\{\chi\in X_\sigma:\chi(F)=c\}.
\]
This set is closed.  By the pinning hypothesis, $Z$ has full measure under every
admissible $\Ext_\sigma(\phi_0)$, so the support of every such extension is contained in
$Z$.  Since $S_\sigma$ is the closure of the union of these supports, and $Z$ is closed,
we have $S_\sigma\subseteq Z$.  But the assumed point $\psi\in Q_\sigma$ satisfies
$\psi(F)\ne c$, hence $\psi\notin S_\sigma$.  Therefore $S_\sigma\ne Q_\sigma$, so the
triple is not root-plantable.
\end{proof}

\begin{remark}[Endpoints are automatic and checkable]
At the endpoints $c\in\{0,1\}$ the almost-sure pinning follows from a mere
expectation condition, since a $[0,1]$-valued random variable with mean $0$
(respectively $1$) is almost surely $0$ (respectively $1$).  Concretely, by
\eqref{eq:extension-expectation}, $\phi_0(\avg{F}{\sigma})=0$ for all admissible
$\phi_0$ gives the case $c=0$; if there is $\psi\in Q_\sigma$ with $\psi(F)>0$, then
the signed element $-F$ is non-negative on $S_\sigma$ and negative at $\psi$.  Similarly,
$\phi_0(\avg{F}{\sigma})=\phi_0(\brak\sigma)$ gives the case $c=1$; if there is
$\psi\in Q_\sigma$ with $\psi(F)<1$, then $F-\one_\sigma$ is non-negative on
$S_\sigma$ and negative at $\psi$.  The degeneracy obstruction
(\cref{thm:degenerate-obstruction}) and its dense
counterpart (\cref{cor:codegenerate}) are these two endpoint cases with $\sigma=\vtype$
and $F=e$.  An interior value $c\in(0,1)$ would instead require the almost-sure pinning as a
genuine hypothesis---but for an edge-deletion-closed hereditary class no interior value is
ever pinned, as \cref{subsec:boundary} shows.
\end{remark}

\subsection{Boundary pinning for edge-deletion-closed classes}\label{subsec:boundary}

The obstructions exhibited above are all at the boundary, $c\in\{0,1\}$, and the remark
after \cref{thm:pinning} observed that an interior value would require the almost-sure
constancy as a genuine hypothesis.  We now show this boundary phenomenon for hereditary
classes that are also closed under deleting edges.  This is not a change of ambient order:
the flag algebra still uses induced subgraphs.  The extra edge-deletion closure is exactly
what is needed for the thinning argument below, and it is automatic for the usual
ordinary-subgraph-free classes.  This matters for the next section, where the final objects
of interest are \emph{density bounds}: unlabelled extremal inequalities in $\A^0[\T_1]$,
such as Mantel- or Tur\'an-type inequalities among graph densities.  Boundary pinning can
expose a real failure of root-plantability at a non-empty type, but after unlabelling it
gives no stronger asymptotic bound at the empty type.

\begin{theorem}[No interior pinning]\label{thm:no-interior}
Let $\K$ be a hereditary graph class that is closed under deleting edges, and let $\sigma$
be any non-degenerate type.  If a $\sigma$-flag $F$ is pinned to $c$ on $S_\sigma$, then
$c\in\{0,1\}$.
\end{theorem}

\begin{proof}
It is enough to produce a positive homomorphism $\psi_\sigma\in S_\sigma$ that is
$\{0,1\}$-valued on every flag: pinning forces $F\equiv c$ on $S_\sigma$, so then
$c=\psi_\sigma(F)\in\{0,1\}$.

As $\sigma$ is non-degenerate, some $\phi_0\in Q_0$ has positive $\sigma$-density.  Choose
finite graphs $G_n\in\K$ converging to $\phi_0$.  For $\lambda\in(0,1]$, form a random
spanning subgraph $G_n^\lambda$ by keeping each edge of $G_n$ independently with
probability $\lambda$.  Since $\K$ is closed under deleting edges, every realisation of
$G_n^\lambda$ still lies in $\K$.

For each fixed finite graph $H$, the expected induced density of $H$ in $G_n^\lambda$ is
determined by the induced densities of graphs on $|H|$ vertices in $G_n$.  Indeed, for
two graphs $H$ and $M$ on the same number of vertices, let $q_{H,M}(\lambda)$ be the
probability that the random graph obtained from $M$ by keeping each edge independently
with probability $\lambda$ is isomorphic to $H$.  Then
\[
  \E[p(H,G_n^\lambda)]
  =
  \sum_{|M|=|H|} q_{H,M}(\lambda)\,p(M,G_n).
\]
The sum is finite, and $p(M,G_n)$ converges for every $M$ because $G_n\to\phi_0$; hence
the expectation on the left has a limit as $n\to\infty$.  Call this limit
$\alpha_H^\lambda$.

We now choose actual thinned graphs whose finite densities follow these limiting
expectations.  Fix $r$ and $\eps>0$.  For all large $n$, the expectations
$\E[p(H,G_n^\lambda)]$ are within $\eps/2$ of $\alpha_H^\lambda$ for every graph $H$
with at most $r$ vertices.

We claim that, with probability tending to $1$, all the actual densities
$p(H,G_n^\lambda)$ are also close to their expectations for this finite list of graphs.
Here is the estimate.  Write $N=|V(G_n)|$ and fix a graph $H$ on $h\le r$ vertices.  The
random thinning is determined by independent coin flips, one for each edge of $G_n$.  If
we change the outcome of just one coin, say the coin for the edge $xy$, then the induced
$H$-density can change only through samples that contain both $x$ and $y$.  Among all
$h$-vertex samples, the fraction containing $x$ and $y$ is
\[
  \frac{\binom{N-2}{h-2}}{\binom Nh}=O_h(N^{-2})
\]
(and it is $0$ if $h<2$).  Thus changing one coin changes $p(H,G_n^\lambda)$ by at most
$C_hN^{-2}$ for a constant $C_h$ depending only on $h$.  We use the following elementary
concentration estimate.  Let $Y_1,\ldots,Y_m$ be independent random variables and let
$X=f(Y_1,\ldots,Y_m)$.  Suppose that for each $i$, changing only the value of $Y_i$ while
leaving all other $Y_j$ fixed can change $f$ by at most $c_i$.  Then
\[
  \Prob[\,|X-\E[X]|>t\,]\le
  2\exp\left(-\frac{2t^2}{\sum_i c_i^2}\right).
\]
Here $Y_1,\ldots,Y_m$ are the edge-retention coins of the thinning process and
$X=p(H,G_n^\lambda)$.  There are at most $\binom N2$ such coins, and we may take every
$c_i$ to be $C_hN^{-2}$.  Hence
\[
  \sum_i c_i^2 \le \binom N2 C_h^2N^{-4}=O_h(N^{-2}),
\]
so, with $t=\eps/2$, the right-hand side tends to $0$ as $N\to\infty$.  Therefore
\[
  \Prob\bigl[\,|p(H,G_n^\lambda)-\E[p(H,G_n^\lambda)]|>\eps/2\,\bigr]\to0 .
\]
There are only finitely many graphs $H$ with at most $r$ vertices, so a union bound shows
that, with probability tending to $1$, the same estimate holds for all of them
simultaneously.  Hence, for all large $n$, there is at least one realisation of
$G_n^\lambda$ for which
\[
  |p(H,G_n^\lambda)-\alpha_H^\lambda|<\eps
  \qquad\text{for every } |H|\le r .
\]
Taking $r=1,2,\ldots$ and $\eps=1/r$, and choosing $n$ increasing at each step, gives a
deterministic sequence of thinned graphs whose density of every fixed finite graph $H$
converges to $\alpha_H^\lambda$.  This sequence lies in $\K$, so its limit is a
constrained unlabelled positive homomorphism; denote it by $\phi_0^\lambda\in Q_0$.

We now check that this limit can still be rooted by $\sigma$.  Consider an embedding
$\theta:\sigma\to G_n$.  After thinning, the same map $\theta$ is still an embedding of
$\sigma$ if all edges of $\sigma$ are kept; the non-edges of $\sigma$ stay non-edges
because thinning only deletes edges.  Therefore the expected $\sigma$-density of
$G_n^\lambda$ is at least
$\lambda^{e(\sigma)}$ times the $\sigma$-density of $G_n$.  Passing to the deterministic
thinned limit gives
\[
  \phi_0^\lambda(\brak\sigma)\ge
  \lambda^{e(\sigma)}\phi_0(\brak\sigma)>0.
\]
Thus $\Ext_\sigma(\phi_0^\lambda)$ is defined for every $\lambda>0$.

Now let $F'$ be a $\sigma$-flag, and let $m(F')$ be the number of edges of $F'$ not already
belonging to the root type $\sigma$.  The unlabelled graph underlying $F'$ has
at least $e(\sigma)+m(F')$ edges, and the thinning formula above gives
\[
  \phi_0^\lambda(\avg{F'}{\sigma})=O(\lambda^{e(\sigma)+m(F')})
  \qquad(\lambda\to0).
\]
Dividing by the lower bound for $\phi_0^\lambda(\brak\sigma)$ in
\eqref{eq:extension-expectation}, we get
\[
  \E_{\phi^\sigma\sim\Ext_\sigma(\phi_0^\lambda)}[\phi^\sigma(F')]
  =O(\lambda^{m(F')}).
\]
Hence every flag with $m(F')>0$ has density tending to $0$ in probability under
$\Ext_\sigma(\phi_0^\lambda)$.  For each size there is a unique remaining flag, namely
the extension of $\sigma$ by vertices that are pairwise non-adjacent and non-adjacent to
the roots; since the flags of that size sum to $\one_\sigma$, its density tends to $1$ in
probability.

By compactness, take any weak subsequential limit of the measures
$\Ext_\sigma(\phi_0^\lambda)$ as $\lambda\to0$.  The preceding coordinate convergence
forces this limiting measure to be the Dirac measure at a positive homomorphism
$\psi_\sigma\in X_\sigma$ characterised as follows: for each number of vertices,
$\psi_\sigma$ assigns value $1$ to the unique flag extending $\sigma$ with no new edges,
and value $0$ to every other flag of that size.  Thus $\psi_\sigma$ lies in $S_\sigma$ and
is $\{0,1\}$-valued on every flag.
\end{proof}

\begin{remark}[Scope]
\Cref{thm:no-interior} is uniform in $\sigma$, with no restriction on the type, and
subsumes the one-vertex, edge, and clique cases at once.  The edge-thinned limits
$\phi_0^\lambda$ have positive $\sigma$-density for every fixed $\lambda>0$.  Thus each
$\phi_0^\lambda$ is an admissible unlabelled limit for the extension theorem, and it has
an extension measure $\Ext_\sigma(\phi_0^\lambda)$ obtained by choosing an embedding of
$\sigma$ at random and recording the resulting labelled limit.  The estimates in the proof
say that, as $\lambda\to0$, these measures converge in flag-density coordinates to the
probability measure concentrated at the positive homomorphism $\psi_\sigma$: all flags
with a new edge have value $0$, and the unique edgeless extension of each size has value
$1$.
Here is the support-closure argument explicitly.  Take any open neighbourhood $U$ of
$\psi_\sigma$.  Since the topology on $X_\sigma$ is generated by the flag-density
coordinates, $U$ contains a smaller neighbourhood $V$ of $\psi_\sigma$ described by
finitely many inequalities involving finitely many flag densities.  The coordinate
convergence above implies that
\[
  \Ext_\sigma(\phi_0^\lambda)(V)>0
\]
for all sufficiently small $\lambda>0$, and hence also
$\Ext_\sigma(\phi_0^\lambda)(U)>0$.  An open set of positive measure must meet the support
of that measure, so $U$ meets $\supp(\Ext_\sigma(\phi_0^\lambda))$ for some admissible
$\lambda$.  Since this holds for every neighbourhood $U$ of $\psi_\sigma$,
$\psi_\sigma$ belongs to the closure of the union of all admissible supports, which is
exactly $S_\sigma$.
The role of the parameter $\lambda$ is only to approach this limiting positive
homomorphism through admissible extension measures.  It is harmless if
$\phi_0^\lambda(\brak\sigma)$ tends to $0$ as $\lambda\to0$: for every fixed positive
$\lambda$ the value is still nonzero, so $\Ext_\sigma(\phi_0^\lambda)$ is defined, and the
closure in the definition of $S_\sigma$ records the limit as $\lambda$ goes to zero.
Finally, by applying complementation (\cref{lem:complementation}), the same
no-interior-pinning conclusion applies to every class whose image under graph
complementation is edge-deletion-closed, including the dense co-degenerate ones.  The
consequences for density bounds are drawn in \cref{sec:empty-type}.
\end{remark}

\subsection{The \texorpdfstring{$C_5$}{C5}-free edge-type obstruction}\label{sec:c5-edge}

The $C_5$-free class is root-plantable at the one-vertex type (\cref{thm:c5-one-root}) and at the
two-root non-edge type (\cref{thm:c5-nonedge-root}), shown in \cref{subsec:c5-planting} by a local
sparse repair.  The two-root \emph{edge} type behaves oppositely: a book-edge construction gives a
pinning obstruction, so root-plantability fails there.  We first record that no pinning obstruction
arises at the one-vertex type, and then build the edge-type obstruction.

\begin{corollary}[No pinning obstruction for $C_5$-free]\label{cor:c5-no-pin}
The $C_5$-free class has no pinning obstruction at $\sigma=\vtype$ from the edge or the
triangle flag over the one-vertex type.
\end{corollary}

\begin{proof}
By \cref{lem:c5-nbhd}, $G[N(v)]$ is a disjoint union of stars and triangles, so it has
at most $\deg(v)$ edges; hence the number of triangles through $v$ is at most $\deg(v)\le n$,
and the one-root triangle density at every vertex is $O(1/n)$.  Hence the
triangle flag over the one-vertex type has density $0$ on \emph{all} of $Q_{\vtype}$: it is pinned at $0$, but
the quotient attains only $0$, so no obstruction arises.  The edge flag is not
pinned at all: empty graphs give one-root edge density $0$, while balanced complete
bipartite graphs are $C_5$-free and give one-root edge density $\tfrac12$ in the limit.
Thus, after choosing the root uniformly at random, the edge flag has limiting value $0$ in
one $C_5$-free sequence and limiting value $\tfrac12$ in another; it is therefore not
pinned to a single value.
\end{proof}

The same class is not root-plantable at the two-root edge type.  Let $\tau$ be
the type consisting of two roots joined by an edge, and let
$F_\triangle\in\A^\tau[\T_{\mathrm{Graph}}]$ be the three-vertex $\tau$-flag in which
the one unlabelled vertex is adjacent to both roots; equivalently, $F_\triangle$ is the
triangle $\tau$-flag over the edge type.

\begin{lemma}[Few triangles]\label{lem:c5-few-triangles}
If $G$ is $C_5$-free, then the number $T(G)$ of triangles in $G$ satisfies
\[
  T(G)\le\frac23 e(G).
\]
In particular, every unlabelled $C_5$-free limit has triangle density zero.
\end{lemma}

\begin{proof}
For every $v$, \cref{lem:c5-nbhd} gives $e(G[N(v)])\le d(v)$.  Summing over
vertices,
\[
  3T(G)=\sum_v e(G[N(v)])\le \sum_v d(v)=2e(G),
\]
since every triangle is counted once at each of its three vertices.  As
$e(G)=O(|V(G)|^2)$, the normalised triangle density is $O(1/|V(G)|)$ along any
sequence with $|V(G)|\to\infty$.
\end{proof}

\begin{corollary}[The triangle edge-type flag is pinned under random edge-rooting]
\label{cor:c5-edge-pinned}
Let $\phi_0\in Q_0$ be a $C_5$-free unlabelled limit with
$\phi_0(\brak{\tau})>0$, and let $\phi^\tau\sim\Ext_\tau(\phi_0)$.  Then
\[
  \phi^\tau(F_\triangle)=0
  \qquad\text{almost surely.}
\]
\end{corollary}

\begin{proof}
The downward average $\avg{F_\triangle}{\tau}$ is a positive scalar multiple of the
unlabelled triangle.  By \cref{lem:c5-few-triangles},
$\phi_0(\avg{F_\triangle}{\tau})=0$.  By \eqref{eq:extension-expectation},
\[
  \E[\phi^\tau(F_\triangle)]
  =
  \frac{\phi_0(\avg{F_\triangle}{\tau})}{\phi_0(\brak{\tau})}
  =0.
\]
Since $\phi^\tau(F_\triangle)\ge0$ pointwise, it is almost surely zero.
\end{proof}

\begin{definition}[Book graphs]\label{def:c5-book}
Let $B_n$ be the graph with two distinguished vertices $a,b$, the edge $ab$,
and $n$ further vertices each adjacent to both $a$ and $b$, with no edges among
the further vertices.
\end{definition}

\begin{lemma}[The book edge is $C_5$-free]\label{lem:c5-book}
Each $B_n$ is $C_5$-free.  With the two roots placed at $a$ and $b$ in that order, the sequence
$(B_n,a,b)$ has a convergent subsequence whose limit is a point $\psi\in Q_\tau$
with
\[
  \psi(F_\triangle)=1.
\]
\end{lemma}

\begin{proof}
Every vertex of $B_n$ other than $a$ and $b$ is adjacent only to $a$ and $b$.
Suppose such a further vertex $x$ lies on a cycle with no repeated vertices.  Then the two
neighbours of $x$ along the cycle must be $a$ and $b$, because $x$ has no other
neighbours in $B_n$.  Thus a cycle using one further vertex is just a triangle through
$a$ and $b$, and a cycle using two further vertices can only be the $4$-cycle
$a,x,b,y,a$ up to reversing the order.  A cycle cannot use three further vertices, since
each of them would need both $a$ and $b$ as its two neighbours on the cycle, forcing $a$ or
$b$ to appear more than once.  Hence $B_n$ has no cycle of length $5$, and so no $C_5$.

For the $\tau$-flag density, every unlabelled vertex of $B_n$ is adjacent to both
roots.  Therefore $p(F_\triangle,(B_n,a,b))=1$ for all $n$, and any convergent subsequence
has the asserted limit.
\end{proof}

\begin{theorem}[The edge type is not root-plantable]\label{thm:c5-edge-not-root-plantable}
The $C_5$-free graph class is not root-plantable at the two-root edge type
$\tau$.  Concretely,
\[
  \Prob_{\phi^\tau\sim\Ext_\tau(\phi_0)}[\phi^\tau(-F_\triangle)\ge0]=1
  \qquad\text{for every }\phi_0\in Q_0\text{ with }\phi_0(\brak{\tau})>0,
\]
but $-F_\triangle$ is negative at some $\psi\in Q_\tau$.
\end{theorem}

\begin{proof}
First fix any constrained unlabelled limit $\phi_0\in Q_0$ with
$\phi_0(\brak\tau)>0$.  This positivity means that the random edge-rooting measure
$\Ext_\tau(\phi_0)$ is defined.  By \cref{cor:c5-edge-pinned}, if
$\phi^\tau$ is sampled from this measure, then
$\phi^\tau(F_\triangle)=0$ almost surely.  Hence
\[
  \phi^\tau(-F_\triangle)=0
\]
almost surely, and in particular it is almost surely non-negative.  This proves the
displayed probability-one assertion.

Next we show that the same signed flag is negative at a labelled quotient homomorphism.
It is useful first to record where the random edge-rooting measures can live.  Let
\[
  Z:=\{\chi\in X_\tau:\chi(F_\triangle)=0\}.
\]
The set $Z$ is closed, since evaluation at $F_\triangle$ is continuous.  The first
paragraph says that, for every $\phi_0\in Q_0$ with $\phi_0(\brak\tau)>0$, the measure
$\Ext_\tau(\phi_0)$ gives probability one to $Z$.  A probability measure whose mass is
one on a closed set has support contained in that closed set: outside $Z$ there is the
open set $X_\tau\setminus Z$, and it has measure zero.  Therefore every support appearing
in the definition of $S_\tau$ is contained in $Z$.  Since $S_\tau$ is the closure of the
union of those supports, and $Z$ is closed, we get
\[
  S_\tau\subseteq Z .
\]

On the other hand, \cref{lem:c5-book} gives a labelled quotient homomorphism
$\psi\in Q_\tau$ with $\psi(F_\triangle)=1$.  Thus $\psi\notin Z$, and hence
$\psi\notin S_\tau$.  This implies that
\[
  \psi(-F_\triangle)=-1<0
\]
and $S_\tau\ne Q_\tau$, i.e., the $C_5$-free class is not root-plantable at the type $\tau$.
\end{proof}

\begin{remark}
The $C_5$-free class therefore gives both a positive and a negative lesson.  At the
one-vertex type, a randomly chosen root can see a genuine positive fraction of the graph:
for example, in a balanced complete bipartite graph its one-root edge density is
$\tfrac12$.  The local planting theorem for this class, \cref{thm:c5-one-root}, says that
every one-root quotient view can be approximated by views obtained from randomly chosen
roots in $C_5$-free graphs.

The two-root edge type behaves differently.  If the two roots are chosen by first taking a
positive edge-density $C_5$-free limit and then choosing a random edge, the
common-neighbour triangle flag has value $0$ almost surely by
\cref{cor:c5-edge-pinned}.  But the book graphs contain a special edge $ab$ whose
common-neighbour density is $1$.  This special edge is not seen by random edge-rooting in
any positive edge-density limit, but it still produces a legitimate quotient homomorphism
in $Q_\tau$.  Thus the edge type has a pinning obstruction even though the one-vertex and
non-edge types are root-plantable; the obstruction analysis has to be done type by type.
\end{remark}

\section{The gap is invisible to density bounds}\label{sec:empty-type}

The obstructions above all occur at non-empty types.  A \emph{pinning}
obstruction means that some labelled flag \(F\) has a fixed value \(c\) in
almost every random rooted view of every constrained unlabelled limit, even
though a specially chosen quotient limit at this type can give \(F\) a different
value.  Algebraically, the corresponding witness is \(F-c\,\one_\sigma\) or
\(c\,\one_\sigma-F\): it vanishes on the ensemble support but not on all of
\(Q_\sigma\).  In the one-vertex examples, this type is \(\sigma=\vtype\), and
the witnesses are \(-e\) and \(e-\one_{\vtype}\), where \(\vtype\) is the
one-vertex type and \(\one_{\vtype}\) is the unit of its flag algebra.  A
flag-algebra \emph{application}, by contrast, outputs a density inequality in
the empty-type algebra $\A^0[\T_1]$, namely, a statement $f\ge0$ on $Q_0$, such as a
Mantel or Tur\'an bound.  One should therefore ask whether the incompleteness
just exhibited can ever affect such a bound.  It cannot, for two independent
reasons: the two semantics coincide at the empty type, and the witnesses that
expose the gap unlabel to zero.

\begin{proposition}[Empty-type collapse]\label{prop:empty-type}
For the empty type, $\Ext_\emptyset(\phi_0)=\delta_{\phi_0}$ for every $\phi_0\in Q_0$;
hence $S_\emptyset=Q_\emptyset=Q_0$.  Thus $(\T_0,\T_1,\emptyset)$ is always
root-plantable, and quotient and ensemble semantics coincide for every $f\in\A^0[\T_0]$.
\end{proposition}

\begin{proof}
The empty type admits a unique embedding into any graph, so the random rooting carries
no information; we read this off \eqref{eq:extension-expectation}.  For
$\sigma=\emptyset$ the unlabelling operator is the identity and $\brak\emptyset=\one_0$, so
$\phi_0(\brak\emptyset)=1$ and $\mu:=\Ext_\emptyset(\phi_0)$ satisfies
$\E_{\psi\sim\mu}[\psi(g)]=\phi_0(g)$ for all $g\in\A^0[\T_0]$.  Applying this to $g$ and to
$g^2$ and using that $\phi_0$ is multiplicative,
\[
  \E_{\psi\sim\mu}\bigl[(\psi(g)-\phi_0(g))^2\bigr]=\phi_0(g^2)-\phi_0(g)^2=0,
\]
so $\psi(g)=\phi_0(g)$ for $\mu$-almost every $\psi$, for each fixed $g\in\A^0[\T_0]$.
Apply this to the countable family of unlabelled flags and intersect the resulting
probability-one events.  Since those flag densities separate the points of $X_0$, this
forces $\mu=\delta_{\phi_0}$, whence
$S_\emptyset=\overline{\bigcup_{\phi_0}\{\phi_0\}}=\overline{Q_0}=Q_0$.  The coincidence
of the two semantics now follows from \cref{thm:support-criterion}.
\end{proof}

\begin{corollary}\label{cor:confined}
No density bound in $\A^0[\T_1]$ is ensemble-true but quotient-false; every failure of
the quotient/ensemble correspondence occurs at a non-empty type $\sigma\ne\emptyset$.  
\end{corollary}

\Cref{cor:confined} settles the question of \emph{truth}.  A flag-algebra proof, though,
reaches its empty-type conclusion through labelled non-negative terms: elements
$s_i\in\A^{\sigma_i}[\T_1]$ whose unlabelled averages $\avg{s_i}{\sigma_i}$ appear in
the final empty-type identity.  Thus one must still ask
whether the labelled gap can affect \emph{provability}: whether an empty-type inequality
could be certified by an ensemble-relaxed proof system but not by the quotient
flag-algebra proof system.  We measure empty-type contributions in the usual constrained seminorm
\[
  \|u\|_{Q_0}:=\sup_{\phi_0\in Q_0}|\phi_0(u)|
  \qquad (u\in\A^0[\T_1]).
\]
A standard flag-algebra certificate writes
\[
  f-\sum_i\avg{s_i}{\sigma_i}=(\text{a manifestly non-negative combination of flags}),
\]
each $s_i$ a sum of squares in $\A^{\sigma_i}[\T_1]$.  Because $\avg{s_i}{\sigma_i}$
integrates $s_i$ against $\Ext_{\sigma_i}(\phi_0)$, a measure supported on
$S_{\sigma_i}$, non-negativity of $s_i$ on $S_{\sigma_i}$ already suffices---weaker than
the sum-of-squares condition, which is non-negativity on all of $Q_{\sigma_i}$.
Equivalently, at a fixed type $\sigma$, an \emph{ensemble-relaxed} flag-algebra proof
system could enlarge the allowed cone of labelled terms to all elements non-negative on
$S_\sigma$.  Write, as cones inside $\A^0[\T_1]$,
\[
  \C^{\mathrm{quot}}_\sigma=\{\avg{s}{\sigma}:s\text{ a sum of squares}\},
  \qquad
  \C^{\mathrm{ens}}_\sigma=\{\avg{s}{\sigma}:s\ge0\text{ on }S_\sigma\},
\]
for the empty-type contributions the type $\sigma$ can make to the two methods, so that
$\C^{\mathrm{quot}}_\sigma\subseteq\C^{\mathrm{ens}}_\sigma$ always.  The next result
shows that this inclusion becomes equality after closure.  Here closure is taken in the
$Q_0$-seminorm because an empty-type expression is used in a density problem only through
its values on constrained limits: $\|u-v\|_{Q_0}<\eps$ means that every admissible limit
$\phi_0\in Q_0$ evaluates $u$ and $v$ within $\eps$.

Concretely, a labelled term $s\in\A^\sigma[\T_1]$ enters an empty-type certificate only
through its averaged term $\avg{s}{\sigma}$.  The theorem says that if
$s$ is merely non-negative on $S_\sigma$, then for every $\eps>0$ there is an ordinary
sum-of-squares term $q\in\A^\sigma[\T_1]$ such that
\[
  \sup_{\phi_0\in Q_0}
  \left|
    \phi_0\!\left(\avg{s}{\sigma}\right)
    -
    \phi_0\!\left(\avg{q}{\sigma}\right)
  \right|
  <\eps .
\]
Thus any density certificate using the ensemble-allowed term $\avg{s}{\sigma}$ gives,
after allowing an arbitrarily small additive slack in the final bound, a quotient
certificate using $\avg{q}{\sigma}$.  This is the natural notion for asymptotic density
bounds, which are statements about the limiting values on $Q_0$ rather than about a
particular finite algebraic identity.  Thus the labelled gap cannot improve an
asymptotic density bound, although exact finite certificates at non-empty types may still
differ.

\begin{theorem}[No closed certificate gap]\label{thm:no-closed-certificate-gap}
For every hereditary class $\T_1$ and every non-degenerate type $\sigma$,
\[
  \overline{\C^{\mathrm{quot}}_\sigma}^{\,\|\cdot\|_{Q_0}}
  =
  \overline{\C^{\mathrm{ens}}_\sigma}^{\,\|\cdot\|_{Q_0}} .
\]
\end{theorem}

Thus allowing labelled terms \(s\in\A^\sigma[\T_1]\) that are non-negative only on
$S_\sigma$ does not enlarge the $Q_0$-closed cone of empty-type certificate
contributions.  It may still give a finite exact certificate for a boundary point that
quotient certificates reach only as a $Q_0$-limit, equivalently only with arbitrarily
small additive slack; any such exact-certifiability gap must occur at non-empty types.

\begin{proof}
It suffices to prove
$\C^{\mathrm{ens}}_\sigma\subseteq
\overline{\C^{\mathrm{quot}}_\sigma}^{\,\|\cdot\|_{Q_0}}$.
Let $s\in\A^\sigma[\T_1]$ be non-negative on $S_\sigma$.  As a function on the
compact set $S_\sigma$, $s$ is continuous and non-negative: each
$\psi\in S_\sigma\subseteq Q_\sigma$ factors through the quotient
$\A^\sigma[\T_0]\twoheadrightarrow\A^\sigma[\T_1]$, say
$\psi=\bar\psi\circ\qmap_\sigma$, so an element
$a\in\A^\sigma[\T_1]$ defines the continuous function
\[
  \widehat a:S_\sigma\to\R,\qquad \widehat a(\psi)=\bar\psi(a).
\]
We use the shorthand $\psi(a)$ for this value.  These functions form a unital real
subalgebra of $C(S_\sigma)$, the algebra of continuous real-valued functions on
$S_\sigma$.  They separate points:
if $\psi_1,\psi_2\in S_\sigma$ are distinct labelled limits, then they differ on some
$\sigma$-flag density, equivalently on some element of $\A^\sigma[\T_1]$.  Applying
Stone--Weierstrass to the continuous function $\sqrt{s}$, and choosing the approximation
sufficiently close, for every
$\delta>0$ there is $h\in\A^\sigma[\T_1]$ such that
\[
  \sup_{\psi\in S_\sigma}\left|s(\psi)-h(\psi)^2\right|<\delta .
\]
The element $h^2$ is a square, so $\avg{h^2}{\sigma}\in\C^{\mathrm{quot}}_\sigma$.

Let $\phi_0\in Q_0$.  If $\phi_0(\brak\sigma)>0$, then
$\Ext_\sigma(\phi_0)$ is supported on $S_\sigma$, and
\eqref{eq:extension-expectation} gives
\begin{align*}
  \left|\phi_0\!\left(\avg{s}{\sigma}-\avg{h^2}{\sigma}\right)\right| 
  =
  \left|\phi_0\!\left(\avg{s-h^2}{\sigma}\right)\right|
  & =
  \phi_0(\brak\sigma)\,
  \left|\E_{\psi\sim\Ext_\sigma(\phi_0)}[\psi(s-h^2)]\right|
  \\
  & = 
  \phi_0(\brak\sigma)\,
  \left|\E_{\psi\sim\Ext_\sigma(\phi_0)}[s(\psi)-h(\psi)^2]\right|
  \\
  & 
  \le \phi_0(\brak\sigma)\,\delta
  \le \delta .
\end{align*}
If $\phi_0(\brak\sigma)=0$, then every unlabelled graph appearing in
$\avg{s-h^2}{\sigma}$ contains an embedding of $\sigma$, so its $\phi_0$-density is zero
by monotonicity; hence the same displayed quantity is $0$.  Therefore
\[
  \left\|\avg{s}{\sigma}-\avg{h^2}{\sigma}\right\|_{Q_0}\le\delta .
\]
Since $\delta$ was arbitrary, $\avg{s}{\sigma}$ lies in the
$Q_0$-closure of $\C^{\mathrm{quot}}_\sigma$.
\end{proof}

The following exact observations complement the closed-cone theorem.  They show the
specific witnesses gain nothing even before closure, and that for the first
degenerate obstructions the two cones coincide outright.

\begin{proposition}[The vanishing ideal unlabels to zero]\label{prop:ideal-zero}
Let $\mathcal I_\sigma=\{a\in\A^\sigma[\T_1]:\psi(a)=0\text{ for all }\psi\in S_\sigma\}$,
where each $\psi\in S_\sigma\subseteq Q_\sigma$ is evaluated on $\A^\sigma[\T_1]$ through
the unique homomorphism $\bar\psi\in\Homp(\A^\sigma[\T_1],\R)$ it factors as,
$\psi=\bar\psi\circ\qmap_\sigma$.
Then $\avg{a}{\sigma}=0$ in $\A^0[\T_1]$ for every $a\in\mathcal I_\sigma$.  In
particular, suppose a flag $F$ is pinned to the value $c$ on $S_\sigma$, so that
$F-c\,\one_\sigma$ vanishes on $S_\sigma$.  Then this pinning witness, and every
flag-multiple $(F-c\,\one_\sigma)\,h$, has zero image under the unlabelling operator
$\avg{\cdot}{\sigma}$; equivalently, its average over random embeddings of the type
$\sigma$ is $0$ in the empty-type algebra.  More generally, if
two labelled elements $s,s'\in\A^\sigma[\T_1]$ have the same evaluation function on
$S_\sigma$,
then $\avg{s}{\sigma}=\avg{s'}{\sigma}$.  Hence a strict exact inclusion
$\C^{\mathrm{ens}}_\sigma\supsetneq\C^{\mathrm{quot}}_\sigma$ can occur only if there is
an element $s\in\A^\sigma[\T_1]$ such that the function
$\psi\mapsto\psi(s)$ is non-negative on $S_\sigma$, but is not equal on
$S_\sigma$ to the function $\psi\mapsto\psi(q)$ for any sum of squares
$q\in\A^\sigma[\T_1]$.
\end{proposition}

\begin{proof}
For $\phi_0\in Q_0$ with $\phi_0(\brak\sigma)>0$,
$\phi_0(\avg{a}{\sigma})=\phi_0(\brak\sigma)\,\E_{\psi\sim\Ext_\sigma(\phi_0)}[\psi(a)]$
by \eqref{eq:extension-expectation}, and the expectation vanishes because the measure is
supported on $S_\sigma$ and $a=0$ there.  If $\phi_0(\brak\sigma)=0$, then
$\avg{a}{\sigma}$ is a combination of unlabelled graphs each containing an embedding of
$\sigma$, so each has $\phi_0$-density $0$ by monotonicity and again
$\phi_0(\avg{a}{\sigma})=0$.  Thus $\avg{a}{\sigma}$ vanishes on $Q_0$ and so is the zero
element.  A pinning witness lies in $\mathcal I_\sigma$ by definition, and
$\mathcal I_\sigma$ is an ideal; the final clause holds because
$s-s'\in\mathcal I_\sigma$ whenever $s,s'$ agree on $S_\sigma$.
\end{proof}

\begin{proposition}[Degenerate roots collapse to a point]\label{prop:single-point}
If $\K$ is edge-degenerate (all constrained limits have edge density $0$), or
co-edge-degenerate (all constrained limits have edge density $1$), then $S_{\vtype}$ is a single point
and $\C^{\mathrm{ens}}_{\vtype}=\C^{\mathrm{quot}}_{\vtype}=\R_{\ge0}\,\one_0$, the ray of
non-negative multiples of the empty-type unit: at the one-vertex type the ensemble
relaxation certifies exactly the density bounds the
quotient method does.
\end{proposition}
Note that because of this proposition, the degeneracy and co-degeneracy counterexamples
(\cref{cor:c4-counterexample,cor:degenerate-family,cor:codegenerate}) therefore cost
nothing for density bounds.
\begin{proof}
Take $\K$ edge-degenerate; the co-degenerate case follows by complementation
(\cref{lem:complementation}).  Let $F$ be any \(\vtype\)-flag whose underlying graph $M$ has an
edge.  Put $k:=|M|$, and choose a $k$-vertex subset $U$ uniformly at random from an
$n$-vertex graph $G$.  Write $\rho(G):=e(G)/\binom n2$ for the edge density of $G$.
Let $X=e(G[U])$ be the number of edges induced by $U$.  Each of
the $\binom{k}{2}$ vertex-pairs inside $U$ is, marginally, a uniformly random pair of
vertices of $G$, so it is an edge with probability $\rho(G)$.  By linearity of
expectation,
\[
  \E X=\binom{k}{2}\rho(G).
\]
If $G[U]$ is isomorphic to $M$, then $X\ge1$ because $M$ has an edge.  Hence
\[
  p(M,G)\le \Prob[X\ge1]\le \E X=\binom{k}{2}\rho(G),
\]
where the middle inequality is Markov's inequality.  Passing to the limit gives
$\phi_0(M)\le\binom{|M|}2\phi_0(\rho)=0$ for every $\phi_0\in Q_0$.  As $\avg{F}{\vtype}$ is a
non-negative multiple of $M$, we get
$\E_{\phi^{\vtype}\sim\Ext_{\vtype}(\phi_0)}[\phi^{\vtype}(F)]=\phi_0(\avg{F}{\vtype})=0$, and since
$\phi^{\vtype}(F)\ge0$, $\phi^{\vtype}(F)=0$ almost surely.  Thus, every support point of every
$\Ext_{\vtype}(\phi_0)$ assigns density $0$ to every edge-containing \(\vtype\)-flag; exactly one
homomorphism $o\in X_{\vtype}$ does so (the labelled empty-graph limit, the others being
determined by the chain relations), so $S_{\vtype}=\{o\}$.

Now let $s\in\A^{\vtype}[\T_1]$ be non-negative on $S_{\vtype}$.  Since
$S_{\vtype}$ has only the point $o$, the only value of $s$ seen by the ensemble support
is the number $s(o)\ge0$.  Hence $s$ agrees on $S_{\vtype}$ with the constant labelled
element $s(o)\,\one_{\vtype}$.  By \cref{prop:ideal-zero}, agreeing on the support gives
the same unlabelled average, so
\[
  \avg{s}{\vtype}
  =
  \avg{s(o)\,\one_{\vtype}}{\vtype}
  =
  s(o)\,\brak{\vtype}
  =
  s(o)\,\one_0 .
\]
The constant labelled element $s(o)\,\one_{\vtype}$ is a sum of squares, namely
$(\sqrt{s(o)}\,\one_{\vtype})^2$.  Thus every ensemble-allowed contribution is already a
quotient sum-of-squares contribution, and both cones are exactly
$\R_{\ge0}\,\one_0$.
\end{proof}

\begin{remark}[What a practically relevant exact gap would require]
\Cref{prop:ideal-zero,prop:single-point} say that the counterexamples above disappear
after applying the downward operator.  The pinned witnesses themselves, and all their
flag-multiples, unlabel to zero; in the degenerate one-root cases, even the whole
ensemble cone collapses under $\avg{\cdot}{\vtype}$ to the single ray
$\R_{\ge0}\one_0$.  Thus, these failures of root-plantability do not create a new
empty-type density certificate.

A gap that survives at the exact certificate level would therefore need something more:
a \emph{Positivstellensatz gap on the support variety}.  In concrete terms, one would need
a non-root-plantable type whose $S_\sigma$ supports a non-negative evaluation function
coming from some $s\in\A^\sigma[\T_1]$, but no sum of squares has the same restriction to
$S_\sigma$.  By \cref{thm:no-closed-certificate-gap}, any such gap is necessarily
non-closed: after unlabelling it can be approximated arbitrarily well by ordinary
quotient sums of squares.  Boundary pinning to a single point supplies no exact gap, and
by \cref{thm:no-interior} pinning is \emph{always} to a boundary value.  The $C_5$-free
class illustrates both possibilities at once: the one-vertex and non-edge-type supports
are rich but plantable, while the edge-type support has a boundary pinning obstruction
coming from common neighbours of an exceptional book edge.
\end{remark}

As a final, concrete instance, take the $C_5$-free class at the edge type $\tau$ of
\cref{sec:c5-edge}, which is not root-plantable (\cref{thm:c5-edge-not-root-plantable});
its gap, too, is closed-cone inert.

\begin{corollary}[The edge-type gap is closed-cone inert]\label{cor:c5-edge-closed-inert}
For the $C_5$-free class and the edge type $\tau$,
\[
  \overline{\C^{\mathrm{quot}}_\tau}^{\,\|\cdot\|_{Q_0}}
  =
  \overline{\C^{\mathrm{ens}}_\tau}^{\,\|\cdot\|_{Q_0}} .
\]
Thus, the ensemble-relaxed edge-type labelled terms do not improve any asymptotic
$C_5$-free density bound, even though $\tau$ is not root-plantable.
\end{corollary}

\begin{proof}
This is the special case of \cref{thm:no-closed-certificate-gap} for the
non-degenerate edge type $\tau$.  The pinned witness $F_\triangle$ (the triangle $\tau$-flag) is even
more inert: by \cref{prop:ideal-zero}, $F_\triangle$ and all its flag-multiples unlabel to zero, since
$F_\triangle=0$ on $S_\tau$ (\cref{cor:c5-edge-pinned}).
\end{proof}

\section{Strengthening by a further constraint}

The framework is not tied to a single forbidden-subgraph constraint.  It applies
unchanged to any hereditary class, in particular to the intersection $\K\cap\K'$ of the
constrained class with a further graph property $\K'$, provided the intersection is again
hereditary.  When the further constraint is \emph{not} hereditary, there is no quotient
algebra to retreat to, but ensemble semantics still makes sense relative to an arbitrary
set of admissible limits.  This section develops that relative theory in its natural
generality---soundness, a support-closure criterion, a complementary-slackness
principle in exact, quantitative, and matrix-kernel forms, and completeness theorems
(a relative Positivstellensatz and a relative closed-cone collapse) showing the
resulting slice method misses nothing---and then instantiates it twice: on the
Tur\'an equality slices, of which the Mantel slice is the first case, and on the equality
slices of verified certificates, where the slackness principle turns a density bound into
labelled local equations, moment identities, rigidity theorems for the extremal graphon,
and stability estimates.

Four classical results enter this section as unproved external inputs: the
Erd\H{o}s--Simonovits stability theorem (uniqueness of the balanced complete $r$-partite
graphon as the extremiser of Tur\'an's theorem at exact equality, \cref{thm:turan-slice},
\cref{thm:relative-mantel}); Zykov's $K_4$-density \emph{bound} for $r\ge4$ (the plain
inequality, which underlies the parametric slice equations of
\cref{thm:parametric-p4-equality-slice} and everything downstream of them; at $r=3$ it is
derived internally from $K_4$-freeness); the \emph{equality case} of the same theorem
(uniqueness of the balanced complete $r$-partite graphon as its extremiser,
\cref{cor:parametric-p4-turan-recovery}, \cref{thm:parametric-quant-stability}); and the
existence half of the Lov\'asz--Szegedy representation theorem (every unlabelled limit is
represented by some graphon, \cref{thm:k4free-p4-tripartite},
\cref{cor:top-endpoint-recovery}).  Where a result of this section depends on one of
these, the dependence is recorded as an explicit hypothesis of the result, in the same
style already used for the Zykov equality case throughout \S11.7--\S11.8; no classical
theorem is left to enter silently through a proof-only citation.

\subsection{When the further constraint is hereditary}

If $\K'$ is itself a hereditary graph property---closed under induced subgraphs---then the
intersection $\T_1:=\K\cap\K'$ is hereditary, the quotient flag algebra $\A^\sigma[\T_1]$
is well-defined, and the support-closure criterion (\cref{thm:support-criterion}) applies
verbatim.  The equivalence of the two semantics is thus governed by root-plantability and
by nothing weaker: it does \emph{not} hold for all hereditary classes, as the degeneracy
obstruction (\cref{thm:degenerate-obstruction})---already in the $C_4$-free
instance---refutes equivalence under heredity alone.  The substantive task is therefore
to find checkable sufficient conditions for root-plantability.  Clone-closure
(\cref{thm:clone-root-plantable}), true-clone-closure (\cref{thm:true-clone-root-plantable}),
and substitution-closure (\cref{thm:substitution-root-plantable}) are three such
conditions, all subsumed by the single \emph{blow-up-closure} criterion
(\cref{thm:blowup-root-plantable})---of which substitution-closure is in fact the strictest,
not the broadest (\cref{rem:strictness}).  The remaining natural candidates involve
structural hypotheses beyond sparse root-blow-up repairs
(\cref{thm:sparse-repair-planting}), such as repairs incident to root clusters,
multi-scale repairs, and amalgamation properties for classes that are not globally
closed under substitution.

\subsection{Relative ensemble enhancements}

If the further constraint is not hereditary---for instance connectivity, a parity
condition, or a prescribed density value---then $\K\cap\K'$ is generally not closed under
induced subgraphs.  This breaks the basic starting point of Razborov's quotient
construction, which assumes that every induced subgraph of a graph in the class is again
in the class.  In this case there may be no canonical quotient flag algebra for
$\K\cap\K'$ of the usual kind.

Ensemble semantics has more room.  Instead of forming a new quotient algebra, fix a
nonempty set
\[
  Y\subseteq Q_0
\]
of admissible unlabelled limits inside an existing hereditary class.  The set $Y$ may
be cut out by arbitrary continuous density constraints, for example
$\phi_0(g_i)=0$, $\phi_0(h_j)\ge0$, or $\phi_0(\rho)=c$---but no structure whatever is
assumed: $Y$ need not be closed, and need not be describable by density constraints at
all.  The relative labelled support at type $\sigma$ is
\[
  S_\sigma(Y):=
  \overline{
  \bigcup_{\substack{\phi_0\in Y\\ \phi_0(\brak\sigma)>0}}
  \supp(\Ext_\sigma(\phi_0))
  }.
\]
Relative ensemble positivity asks only for non-negativity on $S_\sigma(Y)$.  This is a
genuine enhancement of the class: it can impose limit-level conditions that are meaningful for
the intended extremal problem but are invisible to hereditary quotients.  Taking $Y=Q_0$
recovers the absolute theory: $S_\sigma(Q_0)$ is exactly the root-planting set $S_\sigma$
of \cref{def:root-planting}, so every statement below about $S_\sigma(Y)$ generalises a
statement about $S_\sigma$.

Before using the relative supports we record two structural facts.  The first says that
nothing is lost or gained by closing the constraint set; closedness hypotheses can
therefore be dropped throughout.

\begin{lemma}[Closing the constraint set does not change the support]\label{lem:relative-closure}
Let $Y\subseteq Q_0$ be nonempty, and let $\overline Y$ be its closure in
$X_0:=\Homp(\A^0[\T_0],\R)$ (equivalently in $Q_0$, which is closed in $X_0$).  Then, for
every type $\sigma$,
\[
  S_\sigma(\overline Y)=S_\sigma(Y).
\]
\end{lemma}

\begin{proof}
Since $Y\subseteq\overline Y$, the inclusion $S_\sigma(Y)\subseteq S_\sigma(\overline Y)$
is immediate; both sides are closed, so it suffices to show
$\supp(\Ext_\sigma(\phi_0))\subseteq S_\sigma(Y)$ for every $\phi_0\in\overline Y$ with
$\phi_0(\brak\sigma)>0$.

First, the assignment $\phi\mapsto\Ext_\sigma(\phi)$ is weakly continuous on the set
$\{\phi\in X_0:\phi(\brak\sigma)>0\}$.  The evaluation functions
$\chi\mapsto\chi(g)$, $g\in\A^\sigma[\T_0]$, form a unital \emph{subalgebra} of
$C(X_\sigma)$: linear combinations of evaluations are evaluations by the linearity of
each $\chi$, and products are evaluations because each $\chi$ is multiplicative,
$\chi(g_1)\chi(g_2)=\chi(g_1g_2)$.  This subalgebra separates the points of $X_\sigma$,
so it is dense in $C(X_\sigma)$ by Stone--Weierstrass.  By
\eqref{eq:extension-expectation} the integral of the evaluation of $g$ against
$\Ext_\sigma(\phi)$ equals $\phi(\avg{g}{\sigma})/\phi(\brak\sigma)$, a continuous
function of $\phi$ wherever the denominator is positive.  For probability measures on a
compact space, convergence of integrals on a dense subspace of $C(X_\sigma)$ upgrades to
convergence of the integrals of every continuous function, that is, to weak convergence.

Second, supports are lower semicontinuous along weak convergence: if $P_n\to P$ weakly on
the compact metrizable space $X_\sigma$, then
$\supp(P)\subseteq\overline{\bigcup_n\supp(P_n)}$.  Indeed, let $\chi\in\supp(P)$ and let
$U\ni\chi$ be open; the portmanteau theorem gives $\liminf_nP_n(U)\ge P(U)>0$, so $U$
meets $\supp(P_n)$ for all large $n$, and letting $U$ shrink through a countable
neighbourhood basis of $\chi$ produces points of $\bigcup_n\supp(P_n)$ converging to
$\chi$.

Now let $\phi_0\in\overline Y$ with $\phi_0(\brak\sigma)>0$.  By metrizability choose
$\phi_n\in Y$ with $\phi_n\to\phi_0$.  Evaluation at $\brak\sigma$ is continuous, so
$\phi_n(\brak\sigma)>0$ for all large $n$; discard the finitely many others.  By the
first step $\Ext_\sigma(\phi_n)\to\Ext_\sigma(\phi_0)$ weakly, and by the second step
\[
  \supp(\Ext_\sigma(\phi_0))
  \subseteq
  \overline{\bigcup_n\supp(\Ext_\sigma(\phi_n))}
  \subseteq
  S_\sigma(Y),
\]
the last inclusion because each $\phi_n$ lies in $Y$ and $S_\sigma(Y)$ is closed.
\end{proof}

The second fact is soundness, now in its general form: the constraint set is arbitrary,
and the tested element may be taken in the ambient labelled algebra.

\begin{proposition}[Relative ensemble soundness]\label{prop:relative-soundness}
Let $Y\subseteq Q_0$ be nonempty, let $\sigma$ be a type, and let
$f\in\A^\sigma[\T_0]$ satisfy $\chi(f)\ge0$ for every $\chi\in S_\sigma(Y)$.  Then
\[
  \phi_0(\avg{f}{\sigma})\ge0
  \qquad(\phi_0\in Y).
\]
\end{proposition}

\begin{proof}
Let $\phi_0\in Y$.  If $\phi_0(\brak\sigma)>0$, then $\supp(\Ext_\sigma(\phi_0))$ is one
of the sets whose union defines $S_\sigma(Y)$, so $\psi(f)\ge0$ for
$\Ext_\sigma(\phi_0)$-almost every $\psi$ by \cref{lem:support-as}, and
\eqref{eq:extension-expectation} gives
\[
  \phi_0(\avg{f}{\sigma})
  =
  \phi_0(\brak\sigma)\,
  \E_{\psi\sim\Ext_\sigma(\phi_0)}[\psi(f)]
  \ge0.
\]
If $\phi_0(\brak\sigma)=0$, then every unlabelled graph appearing in
$\avg{f}{\sigma}$ contains an embedding of $\sigma$, hence has $\phi_0$-density $0$ by
monotonicity.  Thus $\phi_0(\avg{f}{\sigma})=0$.
\end{proof}

\begin{remark}[Elements of the quotient algebra]\label{rem:quotient-elements}
The statement for $s\in\A^\sigma[\T_1]$, which is how a certificate produced by a
quotient computation presents itself, is the same statement in disguise: the quotient map
$\qmap_\sigma$ is surjective, every point of $S_\sigma(Y)\subseteq Q_\sigma$ factors
through $\qmap_\sigma$, and unlabelling commutes with the quotient maps (the unlabelling
of a forbidden flag is a non-negative multiple of its forbidden underlying graph, which
the unlabelled quotient sends to $0$).  So one may replace $s$ by any preimage
$f\in\qmap_\sigma^{-1}(s)$ without changing either the hypothesis or the conclusion.  We
use the two formulations interchangeably below.
\end{remark}

The relative theory also has a support-closure criterion, and---in sharp contrast with
the absolute criterion of \cref{thm:support-criterion}---it holds unconditionally.

\begin{proposition}[Relative support-closure criterion]\label{prop:relative-criterion}
Let $Y\subseteq Q_0$ be nonempty, let $\sigma$ be a type, and let $f\in\A^\sigma[\T_0]$.
The following are equivalent.
\begin{enumerate}[label=\textup{(\alph*)}, leftmargin=2.8em]
  \item\label{rel-crit-support} $\chi(f)\ge0$ for every $\chi\in S_\sigma(Y)$.
  \item\label{rel-crit-ensemble} For every $\phi_0\in Y$ with $\phi_0(\brak\sigma)>0$,
  \[
    \Prob_{\psi\sim\Ext_\sigma(\phi_0)}[\psi(f)\ge0]=1.
  \]
\end{enumerate}
\end{proposition}

\begin{proof}
\ref{rel-crit-support}$\Rightarrow$\ref{rel-crit-ensemble}: each support
$\supp(\Ext_\sigma(\phi_0))$ with $\phi_0\in Y$ is contained in $S_\sigma(Y)$, so
\cref{lem:support-as} applies.
\ref{rel-crit-ensemble}$\Rightarrow$\ref{rel-crit-support}: by \cref{lem:support-as},
$f\ge0$ on each such support, hence on their union; the evaluation
$\chi\mapsto\chi(f)$ is continuous, so $f\ge0$ on the closure of the union, which is
$S_\sigma(Y)$.
\end{proof}

\begin{remark}[Relative semantics is complete by construction]
In the absolute setting, ensemble semantics is compared against an \emph{external}
benchmark, the quotient set $Q_\sigma$, and \cref{thm:support-criterion} says the two
agree exactly under root-plantability.  Relative to $Y$ there is no quotient to compare
against---the support set $S_\sigma(Y)$ \emph{is} the semantics---so the analogue of
root-plantability is automatic and the relative theory never has a completeness gap: an
inequality holds almost surely under random rooting from every limit in $Y$ if and only
if it holds everywhere on $S_\sigma(Y)$.  With $Y=Q_0$,
\cref{prop:relative-criterion} says that ensemble semantics coincides with
non-negativity on $S_\sigma$; the root-planting question of this paper is precisely
whether that internal benchmark matches the external one $Q_\sigma$.
\end{remark}

\subsection{Complementary slackness on extremal slices}

\Cref{prop:relative-soundness} says which terms may appear in a relative certificate:
anything non-negative on the relevant relative support.  The next theorem is the general
workhorse of this section.  It packages, over an arbitrary constraint set, the three uses
one makes of such a certificate: soundness; the quantitative control that near-equality
exerts on every term; and the exact vanishing forced on the equality slice.  Every
concrete result in the remainder of the section is an instance.

\begin{theorem}[Relative complementary slackness]\label{thm:relative-slackness}
Let $Y\subseteq Q_0$ be nonempty.  Let $h,n\in\A^0[\T_0]$, let $c\in\R$, and for finitely
many indices $1\le i\le m$ let $\sigma_i$ be a type, $\lambda_i>0$ a real, and
$f_i\in\A^{\sigma_i}[\T_0]$.  Assume:
\begin{enumerate}[label=\textup{(\roman*)}, leftmargin=2.8em]
  \item\label{cs-hyp-supports} $\chi(f_i)\ge0$ for every $\chi\in S_{\sigma_i}(Y)$ and
  every $i$;
  \item\label{cs-hyp-slack} $\phi_0(n)\ge0$ for every $\phi_0\in Y$;
  \item\label{cs-hyp-certificate} the certificate inequality
  \[
    \phi_0(h)+\sum_{i=1}^m\lambda_i\,\phi_0(\avg{f_i}{\sigma_i})+\phi_0(n)\le c
    \qquad(\phi_0\in Y)
  \]
  holds.
\end{enumerate}
Then the following hold.
\begin{enumerate}[label=\textup{(\arabic*)}, leftmargin=2.8em]
  \item\label{cs-soundness} \textup{(Soundness)}  $\phi_0(h)\le c$ for every
  $\phi_0\in Y$.
  \item\label{cs-approx} \textup{(Approximate slackness)}  For every $\phi_0\in Y$,
  writing $\Delta:=c-\phi_0(h)\ge0$,
  \[
    \sum_{i=1}^m\lambda_i\,\phi_0(\avg{f_i}{\sigma_i})+\phi_0(n)\le\Delta;
  \]
  in particular $\phi_0(\avg{f_i}{\sigma_i})\le\Delta/\lambda_i$ for every $i$, and
  $\phi_0(n)\le\Delta$.
  \item\label{cs-exact} \textup{(Exact slackness)}  If $\phi_0\in Y$ attains
  $\phi_0(h)=c$, then $\phi_0(n)=0$ and, for every $i$,
  $\phi_0(\avg{f_i}{\sigma_i})=0$; moreover, when $\phi_0(\brak{\sigma_i})>0$,
  \[
    \psi(f_i)=0
    \qquad\text{for }\Ext_{\sigma_i}(\phi_0)\text{-almost every }\psi.
  \]
  \item\label{cs-global} \textup{(Global vanishing)}  If every $\phi_0\in Y$ attains
  $\phi_0(h)=c$, then $f_i=0$ identically on $S_{\sigma_i}(Y)$, for every $i$.
\end{enumerate}
\end{theorem}

\begin{proof}
By \cref{prop:relative-soundness} and hypothesis \ref{cs-hyp-supports}, each
$\phi_0(\avg{f_i}{\sigma_i})$ with $\phi_0\in Y$ is non-negative.  Together with
\ref{cs-hyp-slack}, every term on the left of \ref{cs-hyp-certificate} other than
$\phi_0(h)$ is non-negative; dropping those terms gives \ref{cs-soundness}, and moving
$\phi_0(h)$ to the right-hand side gives the displayed inequality of \ref{cs-approx},
whose stated consequences follow because each summand is non-negative and each
$\lambda_i$ is positive.

For \ref{cs-exact}, apply \ref{cs-approx} with $\Delta=0$: a sum of non-negative terms is
bounded by $0$, so every term vanishes.  If moreover $\phi_0(\brak{\sigma_i})>0$, then by
\eqref{eq:extension-expectation}
\[
  0=\phi_0(\avg{f_i}{\sigma_i})
  =\phi_0(\brak{\sigma_i})\,\E_{\psi\sim\Ext_{\sigma_i}(\phi_0)}[\psi(f_i)],
\]
so $\psi(f_i)$ is a random variable with zero mean.  It is almost surely non-negative,
because the support of $\Ext_{\sigma_i}(\phi_0)$ is contained in $S_{\sigma_i}(Y)$, where
\ref{cs-hyp-supports} holds; a non-negative random variable with zero mean vanishes
almost surely.

For \ref{cs-global}, fix $i$ and let $\phi_0\in Y$ with $\phi_0(\brak{\sigma_i})>0$.  By
\ref{cs-exact}, both $\psi(f_i)\ge0$ and $\psi(-f_i)\ge0$ hold almost surely, so
\cref{lem:support-as}, applied to $f_i$ and to $-f_i$, makes $f_i$ vanish on
$\supp(\Ext_{\sigma_i}(\phi_0))$.  Thus $f_i$ vanishes on the union of these supports
over $\phi_0\in Y$, and by continuity of $\chi\mapsto\chi(f_i)$ also on the closure of
that union, which is $S_{\sigma_i}(Y)$.
\end{proof}

\begin{remark}[The shape of the hypotheses]\label{rem:cs-shape}
Three comments on how \cref{thm:relative-slackness} is meant to be used.  First, the
standard instance is $f_i=\ell_i^2$ with $\ell_i\in\A^{\sigma_i}[\T_0]$ arbitrary: then
hypothesis \ref{cs-hyp-supports} is automatic at \emph{every} homomorphism, since
$\chi(\ell_i^2)=\chi(\ell_i)^2\ge0$, and conclusions
\ref{cs-exact}--\ref{cs-global} read $\psi(\ell_i)=0$ almost surely, respectively
$\ell_i=0$ identically on $S_{\sigma_i}(Y)$.  But squares are not required: any element
non-negative merely on the relative support qualifies, which is strictly weaker than
being a sum of squares.  Second, the unlabelled slack term $n$ absorbs ingredients that
are non-negative for classical reasons rather than by virtue of being averaged
squares---the Zykov $K_4$-density term in \cref{thm:parametric-p4-equality-slice} below
is of this kind---and the theorem returns the vanishing $\phi_0(n)=0$ on the slice,
which is itself usable information.  Third, hypothesis \ref{cs-hyp-certificate} is
required only \emph{on} $Y$.  A certificate proved in the quotient holds on all of
$Q_0\supseteq Y$ and so qualifies, but genuinely relative certificates, with terms
non-negative only on the relative supports, qualify as well.
\end{remark}

Near-extremal limits admit a quantitative form of the vanishing conclusion.  The bridge
is the Cauchy--Schwarz inequality, which converts the $O(\Delta)$ control on averaged
squares into $O(\sqrt\Delta)$ control on every first moment.

\begin{lemma}[Cauchy--Schwarz for unlabelled averages]\label{lem:relative-cauchy-schwarz}
For every $\phi_0\in X_0$, every type $\sigma$, and all $\ell,g\in\A^\sigma[\T_0]$,
\[
  \phi_0(\avg{\ell g}{\sigma})^2
  \le
  \phi_0(\avg{\ell^2}{\sigma})\,
  \phi_0(\avg{g^2}{\sigma}).
\]
\end{lemma}

\begin{proof}
If $\phi_0(\brak\sigma)=0$, then all three averages vanish at $\phi_0$---each is a
combination of unlabelled graphs containing an embedding of $\sigma$, as in the proof of
\cref{prop:relative-soundness}---and the inequality reads $0\le0$.  Otherwise, by
\eqref{eq:extension-expectation} and the multiplicativity of each $\psi$,
\[
  \phi_0(\avg{\ell g}{\sigma})
  =\phi_0(\brak\sigma)\,\E[\psi(\ell)\psi(g)],
  \qquad
  \phi_0(\avg{\ell^2}{\sigma})
  =\phi_0(\brak\sigma)\,\E[\psi(\ell)^2],
  \qquad
  \phi_0(\avg{g^2}{\sigma})
  =\phi_0(\brak\sigma)\,\E[\psi(g)^2],
\]
the expectations being over $\psi\sim\Ext_\sigma(\phi_0)$.  The claim is the
Cauchy--Schwarz inequality for these expectations, multiplied through by
$\phi_0(\brak\sigma)^2$.  (This is Razborov's Cauchy--Schwarz inequality
\cite{Razborov2007}, re-proved through the extension measure.)
\end{proof}

\begin{corollary}[Certificates control first moments at rate $\sqrt\Delta$]\label{cor:sos-first-moments}
In the setting of \cref{thm:relative-slackness}, suppose $f_i=\ell_i^2$.  Then for every
$\phi_0\in Y$, every $i$, and every $g\in\A^{\sigma_i}[\T_0]$,
\[
  \bigl|\phi_0(\avg{\ell_i\,g}{\sigma_i})\bigr|
  \le
  \sqrt{\tfrac{\Delta}{\lambda_i}}\,
  \sqrt{\phi_0(\avg{g^2}{\sigma_i})},
  \qquad
  \Delta:=c-\phi_0(h).
\]
In particular, taking $g=\one_{\sigma_i}$,
\[
  \bigl|\phi_0(\avg{\ell_i}{\sigma_i})\bigr|
  \le
  \sqrt{\tfrac{\Delta}{\lambda_i}\,\phi_0(\brak{\sigma_i})}
  \le
  \sqrt{\tfrac{\Delta}{\lambda_i}}\,.
\]
\end{corollary}

\begin{proof}
Combine \cref{lem:relative-cauchy-schwarz} with the bound
$\phi_0(\avg{\ell_i^2}{\sigma_i})\le\Delta/\lambda_i$ of
\cref{thm:relative-slackness}\ref{cs-approx}.  For the final display,
$\avg{\one_{\sigma_i}^2}{\sigma_i}=\avg{\one_{\sigma_i}}{\sigma_i}=\brak{\sigma_i}$, and
$\phi_0(\brak{\sigma_i})\le1$.
\end{proof}

In practice a certificate's labelled terms do not arrive as an explicit list of squares.
An SDP solver returns, for each type, a positive semidefinite matrix $Q$ over a vector
$v$ of flags, and after rounding $Q$ is rational; extracting squares means
eigendecomposing $Q$, which in general destroys rationality.  The next form of the
slackness theorem consumes $Q$ directly.  Its equality conclusion is that the labelled
moment vector falls into $\ker Q$, so a certificate yields, mechanically and in exact
arithmetic, precisely $\operatorname{rank}Q$ independent labelled equations per type.

\begin{theorem}[Kernel form of complementary slackness]\label{thm:kernel-slackness}
Let $Y\subseteq Q_0$ be nonempty.  Let $h,n\in\A^0[\T_0]$, let $c\in\R$, and for finitely
many indices $1\le t\le m$ let $\sigma_t$ be a type, $v^{(t)}=(v^{(t)}_1,\dots,v^{(t)}_{k_t})$
a vector of elements of $\A^{\sigma_t}[\T_0]$, and $Q^{(t)}\in\R^{k_t\times k_t}$ a
symmetric positive semidefinite matrix; write
\[
  \bigl\langle Q^{(t)}v^{(t)},v^{(t)}\bigr\rangle
  :=
  \sum_{a,b=1}^{k_t}Q^{(t)}_{ab}\,v^{(t)}_a v^{(t)}_b
  \;\in\;\A^{\sigma_t}[\T_0].
\]
Assume $\phi_0(n)\ge0$ for every $\phi_0\in Y$, and the certificate inequality
\[
  \phi_0(h)
  +\sum_{t=1}^m\phi_0\Bigl(\avg{\langle Q^{(t)}v^{(t)},v^{(t)}\rangle}{\sigma_t}\Bigr)
  +\phi_0(n)\le c
  \qquad(\phi_0\in Y).
\]
Then, writing $\Delta:=c-\phi_0(h)$ for $\phi_0\in Y$:
\begin{enumerate}[label=\textup{(\arabic*)}, leftmargin=2.8em]
  \item\label{ks-soundness} $\phi_0(h)\le c$ for every $\phi_0\in Y$, so $\Delta\ge0$;
  \item\label{ks-approx} for every $\phi_0\in Y$, every $t$, and every $w\in\R^{k_t}$,
  \[
    \phi_0\Bigl(\avg{\bigl(w^{\mathsf T}Q^{(t)}v^{(t)}\bigr)^2}{\sigma_t}\Bigr)
    \le
    \bigl\langle Q^{(t)}w,w\bigr\rangle\,\Delta ;
  \]
  \item\label{ks-exact} if $\phi_0\in Y$ attains $\phi_0(h)=c$, then $\phi_0(n)=0$ and,
  for every $t$ with $\phi_0(\brak{\sigma_t})>0$, the moment vector
  $\psi(v^{(t)}):=(\psi(v^{(t)}_1),\dots,\psi(v^{(t)}_{k_t}))$ satisfies
  \[
    Q^{(t)}\,\psi(v^{(t)})=0
    \qquad\text{for }\Ext_{\sigma_t}(\phi_0)\text{-almost every }\psi ;
  \]
  \item\label{ks-global} if every $\phi_0\in Y$ attains $\phi_0(h)=c$, then
  $\chi(v^{(t)})\in\ker Q^{(t)}$ for every $\chi\in S_{\sigma_t}(Y)$ and every $t$;
  equivalently, each row $w^{\mathsf T}Q^{(t)}$ gives the labelled equation
  $w^{\mathsf T}Q^{(t)}v^{(t)}=0$ on $S_{\sigma_t}(Y)$, and a basis of the row space of
  $Q^{(t)}$ gives $\operatorname{rank}Q^{(t)}$ linearly independent such equations.
\end{enumerate}
\end{theorem}

\begin{proof}
Set $f_t:=\langle Q^{(t)}v^{(t)},v^{(t)}\rangle$.  Every homomorphism $\chi$ is linear
and multiplicative, so
$\chi(f_t)=\langle Q^{(t)}\chi(v^{(t)}),\chi(v^{(t)})\rangle\ge0$ by positive
semidefiniteness; thus hypothesis \ref{cs-hyp-supports} of
\cref{thm:relative-slackness} holds automatically---at every homomorphism, not merely on
the relative supports---and that theorem applies with $\lambda_t=1$.  Conclusion
\ref{ks-soundness} is \cref{thm:relative-slackness}\ref{cs-soundness}.

For \ref{ks-approx}, fix $w$ and put $\ell_w:=w^{\mathsf T}Q^{(t)}v^{(t)}
\in\A^{\sigma_t}[\T_0]$.  For every homomorphism $\psi$, the Cauchy--Schwarz inequality
for the positive semidefinite form $Q^{(t)}$ gives
\[
  \psi(\ell_w)^2
  =\bigl\langle Q^{(t)}\psi(v^{(t)}),w\bigr\rangle^2
  \le
  \bigl\langle Q^{(t)}w,w\bigr\rangle\,
  \bigl\langle Q^{(t)}\psi(v^{(t)}),\psi(v^{(t)})\bigr\rangle
  =
  \bigl\langle Q^{(t)}w,w\bigr\rangle\,\psi(f_t).
\]
If $\phi_0(\brak{\sigma_t})=0$, both sides of \ref{ks-approx} vanish by the monotonicity
argument of \cref{prop:relative-soundness}.  Otherwise, taking expectations over
$\psi\sim\Ext_{\sigma_t}(\phi_0)$ and multiplying by $\phi_0(\brak{\sigma_t})$,
\[
  \phi_0(\avg{\ell_w^2}{\sigma_t})
  \le
  \bigl\langle Q^{(t)}w,w\bigr\rangle\,\phi_0(\avg{f_t}{\sigma_t})
  \le
  \bigl\langle Q^{(t)}w,w\bigr\rangle\,\Delta,
\]
the last step by \cref{thm:relative-slackness}\ref{cs-approx}.

For \ref{ks-exact}, \cref{thm:relative-slackness}\ref{cs-exact} gives $\phi_0(n)=0$ and
$\psi(f_t)=\langle Q^{(t)}\psi(v^{(t)}),\psi(v^{(t)})\rangle=0$ almost surely; for a
positive semidefinite matrix, $\langle Qx,x\rangle=0$ forces $Qx=0$.  For
\ref{ks-global}, \ref{ks-exact} makes each coordinate function
$w^{\mathsf T}Q^{(t)}v^{(t)}$ vanish almost surely under every admissible extension
measure, hence---by \cref{lem:support-as} applied to $\pm\,w^{\mathsf T}Q^{(t)}v^{(t)}$---%
on every support, hence by continuity on the closure $S_{\sigma_t}(Y)$; linear
independence of the equations is linear independence of the rows.
\end{proof}

\begin{remark}[Certificates in practice]\label{rem:kernel-practice}
\Cref{thm:kernel-slackness} is the form of slackness adapted to certificates as they are
actually produced and verified.  A rounded certificate ships rational matrices
$Q^{(t)}$; a rational basis of $\ker Q^{(t)}$, computed by row reduction in exact
arithmetic, determines via \ref{ks-global} the complete list of labelled slice
equations the certificate yields at the type $\sigma_t$---no eigendecomposition, and no
irrational square roots, ever enter.  The scalar form (\cref{prop:equality-slice-vanishing})
is recovered by choosing any decomposition $Q^{(t)}=\sum_i\lambda_iu_iu_i^{\mathsf T}$
with $\lambda_i>0$: the equations $u_i^{\mathsf T}v^{(t)}=0$ span the same space as the
rows of $Q^{(t)}$.  The count is worth noting: the number of independent local equations
mined per type equals $\operatorname{rank}Q^{(t)}$, so a certificate with small kernels
pins the extremal ensemble tightly, while a certificate whose matrices have large
kernels leaves the slice correspondingly unconstrained.  The parametric certificate of
\cref{thm:parametric-p4-equality-slice} corresponds to rank-one and rank-two blocks
whose row spaces are spanned by the displayed equations.
\end{remark}

Finally, a slice that consists of a single limit upgrades automatically to a stability
statement, by nothing more than compactness.

\begin{proposition}[Unique slices give qualitative stability]\label{prop:unique-slice-stability}
Let $Z\subseteq X_0$ be closed---for instance $Z=Q_0$ for any hereditary class---let
$h_j\in\A^0[\T_0]$ and $c_j\in\R$ for countably many indices $j$, and set
\[
  Y:=\{\phi\in Z:\phi(h_j)=c_j\ \text{for all }j\}.
\]
If $Y=\{\phi^\ast\}$ is a single point, then every sequence $(\phi_k)$ in $Z$ with
$\phi_k(h_j)\to c_j$ for every $j$ converges to $\phi^\ast$.
\end{proposition}

\begin{proof}
$Z$ is a closed subspace of the compact metrizable space $X_0$, hence itself compact
metrizable.  If $\phi$ is the limit of a convergent subsequence of $(\phi_k)$, then
$\phi\in Z$, and continuity of each evaluation $\chi\mapsto\chi(h_j)$ gives
$\phi(h_j)=c_j$; so $\phi\in Y$, that is, $\phi=\phi^\ast$.  A sequence in a compact
metric space all of whose convergent subsequences have the same limit converges to that
limit.
\end{proof}

\begin{remark}
Under the standard identification of unlabelled limits with graphons up to
measure-preserving relabelling, in which the density topology coincides with the
cut-metric topology on the compact graphon space,
\cref{prop:unique-slice-stability} is a cut-distance stability statement;
\cref{cor:k4free-p4-qualitative-stability} below is the instance for the $P_4$ slice.
\end{remark}

\subsection{Completeness of the slice method}\label{subsec:slice-completeness}

Before turning to examples we settle how much the slice method can see.  Three
questions arise.  Does the root-planting question of this paper have a relative
analogue, and how does it behave on slices?  Do labelled averaged terms
\(\avg{s}{\sigma}\), with \(s\) non-negative merely on the relative support, prove more
density consequences over the slice than averaged sums of squares do?  And is every
density inequality valid on a slice
reachable by a certificate at all?  The first question we answer by example: slices
generically \emph{break} root-plantability, and that breakage is precisely the mining
signal.  The other two we answer in the affirmative, in the same ``up to closure''
sense as \cref{sec:empty-type}: over a slice, sums of squares remain a complete
source of labelled averaged terms, and one squared penalty in the constraints is the only
genuinely new ingredient a slice certificate ever needs.

\begin{definition}[Relative planted set and relative root-plantability]\label{def:relative-plantability}
Let $Y\subseteq Q_0$ be nonempty and let $\sigma$ be a type.  A point $\chi\in X_\sigma$
is a \emph{$Y$-planted view} if there are finite $\sigma$-flags $(G_n,\theta_n)$ with
every $G_n\in\T_1$ such that $p(F,(G_n,\theta_n))\to\chi(F)$ for every $\sigma$-flag
$F$, while the underlying unlabelled graphs $G_n$ converge to a limit belonging to
$\overline Y$.  Write $Q_\sigma(Y)\subseteq X_\sigma$ for the set of $Y$-planted views.
The pair $(Y,\sigma)$ is \emph{relatively root-plantable} if
$S_\sigma(Y)=Q_\sigma(Y)$.
\end{definition}

Thus $Q_\sigma(Y)$ collects the views obtained by planting the roots \emph{anywhere}
---including at exceptional vertices---in large constrained graphs whose global
statistics converge into $Y$, while $S_\sigma(Y)$ keeps only the views a random rooting
sees.  Taking $Y=Q_0$ recovers the absolute theory.

\begin{proposition}[Structure of the planted set]\label{prop:relative-plantability}
Let $Y\subseteq Q_0$ be nonempty and $\sigma$ a type.
\begin{enumerate}[label=\textup{(\roman*)}, leftmargin=2.5em]
  \item\label{rp-basic} $Q_\sigma(Y)$ is a closed subset of $X_\sigma$ with
  $S_\sigma(Y)\subseteq Q_\sigma(Y)\subseteq Q_\sigma$, and $Q_\sigma(Q_0)=Q_\sigma$;
  in particular, for $Y=Q_0$ relative root-plantability is root-plantability.
  \item\label{rp-criterion} For every $f\in\A^\sigma[\T_0]$, non-negativity of $f$ on
  $Q_\sigma(Y)$ implies the relative ensemble semantics of
  \cref{prop:relative-criterion}\ref{rel-crit-ensemble}; the two are equivalent for
  every $f$ if and only if $(Y,\sigma)$ is relatively root-plantable.
\end{enumerate}
\end{proposition}

\begin{proof}
\ref{rp-basic}.  A $Y$-planted view is a limit of finite flag density vectors, hence a
positive homomorphism \cite[Theorem 3.3]{Razborov2007}, and it vanishes on every
forbidden flag because forbidden flags have density $0$ in every flag with underlying
graph in $\T_1$; by the intrinsic description of $Q_\sigma$ this gives
$Q_\sigma(Y)\subseteq Q_\sigma$.  Closedness is a diagonal argument: if
$\chi^{(k)}\to\chi$ with witnessing sequences for each $\chi^{(k)}$, choose from the
$k$-th sequence a flag whose labelled densities are within $1/k$ of $\chi^{(k)}$ on the
first $k$ flags and whose underlying graph is within $1/k$ of its base limit
$\phi^{(k)}\in\overline Y$ on the first $k$ graphs; since $\overline Y$ is compact, a
subsequence of the $\phi^{(k)}$ converges in $\overline Y$, and along it the chosen
flags witness $\chi\in Q_\sigma(Y)$.  For $S_\sigma(Y)\subseteq Q_\sigma(Y)$, let
$\psi\in\supp(\Ext_\sigma(\phi_0))$ with $\phi_0\in Y$, $\phi_0(\brak\sigma)>0$, and
let finite graphs $G_n\in\T_1$ converge to $\phi_0$.  The finite random-rooting
distributions of the $G_n$ converge weakly to $\Ext_\sigma(\phi_0)$
\cite[Theorems 3.5 and 3.12]{Razborov2007}, so every neighbourhood of $\psi$ receives
positive rooting probability for all large $n$; choosing for each $k$ an $n_k$ and a
rooting $\theta_k$ whose view lies within $1/k$ of $\psi$ produces a witnessing
sequence with base limit $\phi_0\in\overline Y$.  Hence
$\psi\in Q_\sigma(Y)$, and $S_\sigma(Y)\subseteq Q_\sigma(Y)$ follows since the latter
is closed.  Finally, for $Y=Q_0$: every $\chi\in Q_\sigma$ is a limit of finite
$\T_1$-flags \cite[Theorem 3.3]{Razborov2007}, and after passing to a subsequence the
underlying graphs converge, necessarily to a point of $Q_0=\overline{Q_0}$; so
$Q_\sigma\subseteq Q_\sigma(Q_0)$, and the reverse inclusion was shown above.

\ref{rp-criterion}.  By \cref{prop:relative-criterion}, relative ensemble semantics for
$f$ is non-negativity of $f$ on $S_\sigma(Y)$, so the implication follows from
$S_\sigma(Y)\subseteq Q_\sigma(Y)$, and equivalence for every $f$ compares
non-negativity on two closed subsets of $X_\sigma$, one contained in the other.  If the
sets are equal the two conditions coincide; if not, the separation argument in the
proof of \cref{thm:support-criterion}---Urysohn on the compact metrizable
$Q_\sigma(Y)$, Tietze extension to $X_\sigma$, and Stone--Weierstrass
approximation by an element of $\A^\sigma[\T_0]$---produces an $f$ strictly positive
on $S_\sigma(Y)$ and negative somewhere on $Q_\sigma(Y)$.
\end{proof}

\begin{proposition}[Slices break root-plantability: the Mantel slice]\label{prop:mantel-not-plantable}
Let $\T_1$ be the triangle-free class and $Y_{\mathrm{Mantel}}$ the Mantel slice of
\cref{thm:relative-mantel}, and assume the Erd\H{o}s--Simonovits hypothesis of
\cref{thm:relative-mantel}.  Then $(Y_{\mathrm{Mantel}},\vtype)$ is not relatively
root-plantable:
\[
  S_{\vtype}(Y_{\mathrm{Mantel}})\subsetneq Q_{\vtype}(Y_{\mathrm{Mantel}}),
\]
although the triangle-free class itself is root-plantable at every type
(\cref{cor:clique-free}).
\end{proposition}

\begin{proof}
Let $G_n$ be $K_{n,n}$ together with one additional isolated vertex $w_n$, and let
$\theta_n$ root the one-vertex type at $w_n$.  Each $G_n$ is triangle-free, and the
single extra vertex changes every induced density by $O(1/n)$, so the underlying graphs
converge to the balanced complete bipartite limit, which is the point of
$Y_{\mathrm{Mantel}}$.  The rooted views converge as well: a flag containing an edge at
the root has density $0$ for every $n$, and a flag whose root is isolated has, up to
$O(1/n)$, the density of its unrooted part in $K_{n,n}$.  The limit $\chi$ is therefore
a $Y_{\mathrm{Mantel}}$-planted view with one-root edge density $\chi(e)=0$.  By
\cref{thm:relative-mantel}, $S_{\vtype}(Y_{\mathrm{Mantel}})$ is a single point with
$e=\frac12\one_{\vtype}$, so $\chi\notin S_{\vtype}(Y_{\mathrm{Mantel}})$.
\end{proof}

\begin{remark}[Relative pinning is the mining signal]\label{rem:relative-pinning}
\Cref{prop:mantel-not-plantable} is the relative avatar of the pinning phenomenon of
this paper, with its sign reversed.  On the Mantel slice the element
$e-\frac12\one_{\vtype}$ vanishes identically on $S_{\vtype}(Y_{\mathrm{Mantel}})$ yet is
non-zero on $Q_{\vtype}(Y_{\mathrm{Mantel}})$: a \emph{relative pinning witness}.  In
the absolute theory such a witness certified an incompleteness of the quotient
calculus; relative to $Y$ there is no quotient benchmark
(\cref{prop:relative-criterion}), and the witness \emph{is} the extracted information
---the local law of extremisers that mining reports.  The analogue of
\cref{prop:ideal-zero} still holds, with ``zero element'' weakened to ``$Y$-null'': if
$a\in\A^\sigma[\T_1]$ vanishes on $S_\sigma(Y)$, then $\phi_0(\avg{a}{\sigma})=0$ for
every $\phi_0\in Y$, by the argument of \cref{prop:ideal-zero} verbatim.  The Mantel
witness illustrates the sharpened conclusion: $e-\frac12\one_{\vtype}$ unlabels to
$\rho-\frac12\one_0$, which is not the zero element of $\A^0[\T_1]$ but is exactly the
defining constraint of the slice---a $Y$-null element, as it must be.
\end{remark}

We now turn to certificates.  Following \cref{sec:empty-type}, measure empty-type
contributions relative to the slice by the seminorm
\[
  \|u\|_{Y}:=\sup_{\phi_0\in Y}|\phi_0(u)|
  \qquad(u\in\A^0[\T_1]),
\]
so that $\|\cdot\|_{Q_0}$ is the case $Y=Q_0$, and write
\[
  \C^{\mathrm{ens}}_\sigma(Y):=\{\avg{s}{\sigma}:s\in\A^\sigma[\T_1],\
  s\ge0\text{ on }S_\sigma(Y)\},
\]
the empty-type contributions available to an ensemble-relaxed method over the slice;
$\C^{\mathrm{quot}}_\sigma\subseteq\C^{\mathrm{ens}}_\sigma(Y)$ always, since squares
are non-negative everywhere.  The absolute closed-cone collapse of
\cref{thm:no-closed-certificate-gap} survives relativisation unchanged.

\begin{theorem}[No closed certificate gap over a slice]\label{thm:relative-certificate-gap}
For every nonempty $Y\subseteq Q_0$ and every type $\sigma$,
\[
  \overline{\C^{\mathrm{quot}}_\sigma}^{\,\|\cdot\|_{Y}}
  =
  \overline{\C^{\mathrm{ens}}_\sigma(Y)}^{\,\|\cdot\|_{Y}} .
\]
Thus, even over a slice, labelled terms non-negative merely on the relative support prove
no density consequence beyond averaged sums of squares, except possibly at the exact
finite-certificate level.
\end{theorem}

\begin{proof}
If $S_\sigma(Y)=\emptyset$, then $\phi_0(\brak\sigma)=0$ for every $\phi_0\in Y$, and
by the monotonicity argument of \cref{prop:relative-soundness} every unlabelled average
$\avg{s}{\sigma}$ is $Y$-null; both closures then coincide trivially.  Otherwise let
$s\in\A^\sigma[\T_1]$ be non-negative on the compact set $S_\sigma(Y)$ and let
$\delta>0$.  Exactly as in the proof of \cref{thm:no-closed-certificate-gap},
Stone--Weierstrass applied to $\sqrt{s}$ on $S_\sigma(Y)$ produces
$h\in\A^\sigma[\T_1]$ with $\sup_{\psi\in S_\sigma(Y)}|s(\psi)-h(\psi)^2|<\delta$.  For
$\phi_0\in Y$ with $\phi_0(\brak\sigma)>0$ the measure $\Ext_\sigma(\phi_0)$ is
supported on $S_\sigma(Y)$ by the definition of the relative support, so
\eqref{eq:extension-expectation} gives
$|\phi_0(\avg{s-h^2}{\sigma})|\le\delta\,\phi_0(\brak\sigma)\le\delta$; for
$\phi_0(\brak\sigma)=0$ the same quantity is $0$.  Hence
$\|\avg{s}{\sigma}-\avg{h^2}{\sigma}\|_{Y}\le\delta$ with
$\avg{h^2}{\sigma}\in\C^{\mathrm{quot}}_\sigma$.
\end{proof}

The remaining question is whether the certificates reach every inequality that is
\emph{true} on the slice.  They do, up to an arbitrarily small constant, and the only
new ingredient needed is a penalty term: a single non-negative multiple of the sum of
squared constraints.

\begin{theorem}[Relative Positivstellensatz]\label{thm:relative-positivstellensatz}
Let $J$ be a countable index set, let $g_j\in\A^0[\T_1]$ for $j\in J$, and let
\[
  Y:=\{\phi_0\in Q_0:\phi_0(g_j)=0\ \text{for all }j\in J\}
\]
(a slice $\phi_0(h)=c$ enters with $g:=h-c\,\one_0$).  For $f\in\A^0[\T_1]$ the
following are equivalent.
\begin{enumerate}[label=\textup{(\alph*)}, leftmargin=2.8em]
  \item\label{psatz-semantic} $\phi_0(f)\ge0$ for every $\phi_0\in Y$.
  \item\label{psatz-penalty} For every $\eps>0$ there are finitely many indices
  $j_1,\dots,j_N\in J$ and a constant $M\ge0$ such that
  \[
    \phi_0\Bigl(f+\eps\,\one_0+M\sum_{i=1}^{N}g_{j_i}^2\Bigr)\ge0
    \qquad\text{for every }\phi_0\in Q_0 .
  \]
\end{enumerate}
Equivalently, writing $\C_{Q_0}:=\{u\in\A^0[\T_1]:\phi_0(u)\ge0\ \forall\phi_0\in Q_0\}$
for the semantic cone of the class,
\[
  \{f:\phi_0(f)\ge0\ \forall\phi_0\in Y\}
  =
  \overline{\C_{Q_0}+\operatorname{span}\{g_j^2:j\in J\}}^{\,\|\cdot\|_{Q_0}} .
\]
The statement holds for $Y=\emptyset$ as well, both sides then being all of
$\A^0[\T_1]$.
\end{theorem}

\begin{proof}
\ref{psatz-penalty}$\Rightarrow$\ref{psatz-semantic}: evaluate at $\phi_0\in Y$; by
multiplicativity $\phi_0(g_j^2)=\phi_0(g_j)^2=0$, so $\phi_0(f)\ge-\eps$ for every
$\eps>0$.

\ref{psatz-semantic}$\Rightarrow$\ref{psatz-penalty}: fix $\eps>0$ and enumerate
$J=\{j_1,j_2,\dots\}$.  For $n\ge1$ let
\[
  K_n:=\{\phi_0\in Q_0:\phi_0(g_{j_i})^2\le1/n\ \text{for }1\le i\le n\},
\]
a decreasing sequence of closed subsets of the compact space $Q_0$ with
$\bigcap_nK_n=Y$.  The closed sets $K_n\cap\{\phi_0:\phi_0(f)\le-\eps/2\}$ decrease and
have empty intersection (the intersection would lie in $Y$, where $f\ge0$), so by
compactness some $K_N\cap\{\phi_0(f)\le-\eps/2\}=\emptyset$; that is, $\phi_0(f)>-\eps/2$
on $K_N$.  Put $u:=\sum_{i=1}^{N}g_{j_i}^2$, let
$B:=\sup_{\phi_0\in Q_0}|\phi_0(f)|<\infty$ (a continuous function on a compact space),
and set $M:=BN$.  For $\phi_0\in Q_0$: if $\phi_0\in K_N$, then
$\phi_0(f)+\eps+M\phi_0(u)\ge-\eps/2+\eps+0=\eps/2$, since
$\phi_0(u)=\sum_i\phi_0(g_{j_i})^2\ge0$; if $\phi_0\notin K_N$, then
$\phi_0(g_{j_i})^2>1/N$ for some $i\le N$, so $\phi_0(u)>1/N$ and
$\phi_0(f)+\eps+M\phi_0(u)\ge-B+\eps+B=\eps$.  In either case the displayed inequality
of \ref{psatz-penalty} holds (strictly).  For $Y=\emptyset$ the same compactness
argument makes some $K_N$ empty, and the second branch alone applies to any $f$.

For the cone identity: each element of $\C_{Q_0}$ is non-negative on $Y\subseteq Q_0$,
each $g_j^2$ is $Y$-null by multiplicativity, and $\|\cdot\|_{Q_0}$-limits preserve
non-negativity on $Y$ since $\|\cdot\|_Y\le\|\cdot\|_{Q_0}$; this proves
``$\supseteq$''.  Conversely, if $f$ is non-negative on $Y$, then by
\ref{psatz-penalty} $f+\eps\,\one_0\in\C_{Q_0}-M\sum_ig_{j_i}^2\subseteq
\C_{Q_0}+\operatorname{span}\{g_j^2\}$ for every $\eps$, and
$\|f-(f+\eps\one_0)\|_{Q_0}=\eps\to0$.
\end{proof}

\begin{remark}[The slice method is complete, modulo the absolute question]\label{rem:slice-complete}
\Cref{thm:relative-positivstellensatz} reduces non-negativity over a slice to
non-negativity over the whole class, at the price of one squared-constraint penalty and
an arbitrarily small shift of the constant term; and by
\cref{thm:relative-certificate-gap} the labelled-term side over the slice collapses,
after closure, to ordinary sums of squares.  Together: every density
inequality valid on an equality slice is provable, up to $\eps$ and up to closure, by a
certificate of the standard quotient shape augmented by the single term
$M\sum_ig_{j_i}^2$.  In this precise sense the mining procedure of this section is
semantically complete---nothing true on the slice is invisible to it.  What the
reduction does \emph{not} settle is the absolute question it reduces to: whether
$\C_{Q_0}$ itself is the closure of the finite sum-of-squares certificates of the
class.  For the unconstrained class this is the weak Positivstellensatz of
Lov\'asz and Szegedy \cite{LovaszSzegedy2012}; for hereditary subclasses it is open in
general, and it is independent of everything relative in this section.  Inequality
constraints $\phi_0(h)\ge0$, by contrast, genuinely require Positivstellensatz-style
terms attached to the inequality constraints; equality constraints cover every slice used
in this paper.
\end{remark}

\subsection{Example: Tur\'an equality slices}

The simplest slices pin a single density at the extreme value of a classical theorem
whose equality case is known.  The slice is then one limit, and the relative supports
read off its local statistics---which the hereditary quotient deliberately forgets.  We
record the whole Tur\'an family at once; the Mantel slice is the case $r=2$.

\begin{theorem}[Tur\'an equality slices]\label{thm:turan-slice}
Fix $r\ge2$.  Let $\T_r$ be the $K_{r+1}$-free graph class, let $Q_0^{(r)}$ be its space
of constrained unlabelled limits, let $\rho$ be the unlabelled edge density, and set
\[
  Y_{\mathrm{Tur}}^{(r)}
  :=
  \Bigl\{\phi_0\in Q_0^{(r)}:\phi_0(\rho)=\tfrac{r-1}{r}\Bigr\}.
\]
The set $Y_{\mathrm{Tur}}^{(r)}$ contains the balanced complete $r$-partite limit,
represented by the graphon $T_r$; this existence half needs no further input.  Assume, in
addition, the following instance of the Erd\H{o}s--Simonovits stability theorem: $T_r$ is
the \emph{only} point of $Y_{\mathrm{Tur}}^{(r)}$, i.e. $Y_{\mathrm{Tur}}^{(r)}=\{T_r\}$
(equivalently, every $K_{r+1}$-free graphon of edge density $\frac{r-1}r$ is, up to
measure-preserving relabelling, $T_r$).  Then the relative supports at the one-vertex type
$\vtype$, the ordered edge type $\tau$, and the ordered non-edge type $\eta$ are single
points, and the following identities hold.
\begin{enumerate}[label=\textup{(\roman*)}, leftmargin=2.5em]
  \item If $e\in\A^{\vtype}[\T_r]$ is the one-root edge flag, then on
        $S_{\vtype}(Y_{\mathrm{Tur}}^{(r)})$,
        \[
          e=\frac{r-1}{r}\,\one_{\vtype} .
        \]
  \item At $\tau$, let $a_\tau$ and $b_\tau$ be the two three-vertex flags in which the
        new vertex is adjacent to exactly the first, respectively exactly the second,
        root; let $z_\tau$ be the flag in which it is adjacent to neither root; and let
        $g_\tau$ be the common-neighbour flag, in which it is adjacent to both.  On
        $S_\tau(Y_{\mathrm{Tur}}^{(r)})$,
        \[
          a_\tau=b_\tau=\frac1r\,\one_\tau,
          \qquad
          g_\tau=\frac{r-2}{r}\,\one_\tau,
          \qquad
          z_\tau=0 .
        \]
  \item At $\eta$, let $a_\eta,b_\eta,z_\eta,g_\eta$ denote the analogous flags.  On
        $S_\eta(Y_{\mathrm{Tur}}^{(r)})$,
        \[
          z_\eta=\frac1r\,\one_\eta,
          \qquad
          g_\eta=\frac{r-1}{r}\,\one_\eta,
          \qquad
          a_\eta=b_\eta=0 .
        \]
\end{enumerate}
\end{theorem}

\begin{proof}
Existence needs only Tur\'an's theorem: the balanced complete $r$-partite graph attains
edge density $\frac{r-1}r$ exactly, so $T_r\in Y_{\mathrm{Tur}}^{(r)}$.  Uniqueness is
exactly the assumed hypothesis; it is the content of the Erd\H{o}s--Simonovits stability
theorem \cite{ErdosSimonovits,Simonovits1968}, which says that a $K_{r+1}$-free sequence
whose edge density tends to $\frac{r-1}{r}$ converges, after changing $o(n^2)$ edges, to
balanced complete $r$-partite graphs---equivalently, that $T_r$ is the only $K_{r+1}$-free
graphon of edge density $\frac{r-1}r$.

The identities are direct readings of random labelled extensions of $T_r$.  The parts of
$T_r$ are interchangeable and its vertices within a part indistinguishable, so the views
from almost every vertex, from almost every ordered edge, and from almost every ordered
non-edge coincide; each extension measure is therefore a Dirac measure and each relative
support a single point.  A random vertex is adjacent to all parts except its own, giving
(i).  A random ordered edge has its endpoints in two distinct parts; a further random
vertex lies in the first endpoint's part, the second endpoint's part, or one of the
remaining $r-2$ parts with probabilities $\frac1r$, $\frac1r$, and $\frac{r-2}r$, and is
then adjacent to exactly the second root, exactly the first root, or both, respectively;
no position leaves it adjacent to neither root.  This gives (ii).  A random ordered
non-edge has both endpoints in one part; a further random vertex lies in that part with
probability $\frac1r$, adjacent to neither root, and in one of the other $r-1$ parts
with probability $\frac{r-1}r$, adjacent to both.  This gives (iii).
\end{proof}

At $r=2$ this is the Mantel slice of the introduction.  Let $\T_1$ be the triangle-free
graph class (that is, $\T_2$ above), let $\rho\in\A^0[\T_1]$ be the unlabelled edge
density, and set
\[
  Y_{\mathrm{Mantel}}
  :=
  Y_{\mathrm{Tur}}^{(2)}
  =
  \{\phi_0\in Q_0:\phi_0(\rho)=1/2\}.
\]
This is a closed, non-hereditary enhancement: it describes the equality case of Mantel's
theorem at the level of limits, not a hereditary graph class.

\begin{theorem}[Relative Mantel equality support]\label{thm:relative-mantel}
The set $Y_{\mathrm{Mantel}}$ contains the balanced complete bipartite limit; this
existence half needs no further input.  Assume, in addition, the $r=2$ instance of the
Erd\H{o}s--Simonovits hypothesis of \cref{thm:turan-slice}: the balanced complete
bipartite limit is the \emph{only} point of $Y_{\mathrm{Mantel}}$.  Then the relative
supports at the one-vertex type, the edge type $\tau$, and the non-edge type $\eta$ are
all single points.

More explicitly:
\begin{enumerate}[label=\textup{(\roman*)}, leftmargin=2.5em]
  \item If $e\in\A^{\vtype}[\T_1]$ is the one-root edge flag, then on
        $S_{\vtype}(Y_{\mathrm{Mantel}})$ one has
        \[
          e=\frac12\,\one_{\vtype} .
        \]
  \item At the ordered edge type $\tau$, let $a_\tau$ and $b_\tau$ be the two
        three-vertex flags in which the new vertex is adjacent to exactly the first,
        respectively exactly the second, root endpoint; let $z_\tau$ be the flag in
        which it is adjacent to neither endpoint; and let $g_\tau$ denote the quotient
        image of the common-neighbour triangle flag, already zero in the triangle-free
        class.  On $S_\tau(Y_{\mathrm{Mantel}})$,
        \[
          a_\tau=b_\tau=\frac12\,\one_\tau,
          \qquad
          z_\tau=g_\tau=0 .
        \]
  \item At the ordered non-edge type $\eta$, let $a_\eta,b_\eta,z_\eta,g_\eta$ denote the
        analogous flags, where $g_\eta$ is now the flag in which the new vertex is
        adjacent to both non-adjacent roots.  On $S_\eta(Y_{\mathrm{Mantel}})$,
        \[
          g_\eta=z_\eta=\frac12\,\one_\eta,
          \qquad
          a_\eta=b_\eta=0 .
        \]
\end{enumerate}
\end{theorem}

\begin{proof}
This is the case $r=2$ of \cref{thm:turan-slice}, whose Erd\H{o}s--Simonovits
hypothesis specialises at $r=2$ to the one assumed here: the identities there read
$e=\frac12\one_{\vtype}$; $a_\tau=b_\tau=\frac12\one_\tau$ with
$g_\tau=\frac{r-2}{r}\one_\tau=0$ and $z_\tau=0$ (at $r=2$ the common-neighbour flag
$g_\tau$ is the triangle flag, already zero in the triangle-free quotient); and
$z_\eta=g_\eta=\frac12\one_\eta$ with $a_\eta=b_\eta=0$.
\end{proof}

\begin{remark}[Why this is not quotient semantics]
The equality constraint $\phi_0(\rho)=1/2$ is not hereditary, so it is not naturally a
quotient flag algebra.  The relative ensemble support nevertheless records strong
labelled consequences of being an extremal Mantel limit.  For instance,
$e=\frac12\one_{\vtype}$ on $S_{\vtype}(Y_{\mathrm{Mantel}})$, while this equality is false in the
ordinary triangle-free quotient: stars rooted at their centres give quotient points
with $e=1$, and empty graphs give quotient points with $e=0$.  Thus relative ensemble semantics can express
local information about extremisers that the hereditary quotient deliberately forgets.
The same contrast holds for every $r$: on $S_{\vtype}(Y_{\mathrm{Tur}}^{(r)})$ the
one-root edge density is pinned to $\frac{r-1}r$, a statement no hereditary quotient of
the $K_{r+1}$-free class can make.
\end{remark}

\subsection{Example: equality slices from a certificate}

The Mantel slice is a useful sanity check, but the graphon is already known.  A more
exploratory use of relative ensemble semantics is to take a flag-algebra certificate and
study the equality slice it defines.  The certificate then becomes a source of labelled
equations for extremal limits, even when the extremal graphons themselves have not been
classified.

\begin{proposition}[Equality slices force certificate terms to vanish]\label{prop:equality-slice-vanishing}
Let $h\in\A^0[\T_1]$, let $c\in\R$, and let
\[
  Y=\{\phi_0\in Q_0:\phi_0(h)=c\}.
\]
Suppose a certificate proves, on $Q_0$,
\[
  h+\sum_{i=1}^m \lambda_i\,\avg{\ell_i^2}{\sigma_i}\le c\,\one_0,
  \qquad
  \lambda_i>0 .
\]
Then each linear certificate term vanishes on the corresponding relative support:
\[
  \psi(\ell_i)=0
  \qquad
  (\psi\in S_{\sigma_i}(Y),\ 1\le i\le m).
\]
\end{proposition}

\begin{proof}
This is \cref{thm:relative-slackness} with $f_i:=\ell_i^2$ and $n:=0$.  Hypothesis
\ref{cs-hyp-supports} holds at every homomorphism because
$\chi(\ell_i^2)=\chi(\ell_i)^2\ge0$; hypothesis \ref{cs-hyp-certificate} holds on
$Q_0\supseteq Y$ because the certificate is proved on $Q_0$; and every $\phi_0\in Y$
attains $\phi_0(h)=c$ by the definition of $Y$.  Conclusion \ref{cs-global} gives
$\chi(\ell_i)^2=\chi(\ell_i^2)=0$, hence $\chi(\ell_i)=0$, for every
$\chi\in S_{\sigma_i}(Y)$.
\end{proof}

\begin{theorem}[The $K_4$-free $P_4$ equality slice]\label{thm:k4free-p4-equality-slice}
Let $\T_1$ be the $K_4$-free graph class, and let $\pi_{P_4}\in\A^0[\T_1]$ denote
the $P_4$ density expression used in the verified certificate
\texttt{API.K4freeP4}.  Define the extremal slice
\[
  Y_{P_4}:=\{\phi_0\in Q_0:\phi_0(\pi_{P_4})=32/9\}.
\]
Let $\eta$ be the ordered non-edge type and $\tau$ the ordered edge type.

For $\eta$, let $z_\eta$ be the three-vertex flag in which the new vertex is adjacent
to neither root, and let $g_\eta$ be the flag in which it is adjacent to both roots.
Then on $S_\eta(Y_{P_4})$,
\[
  2z_\eta=g_\eta .
\]

For $\tau$, let $a_\tau$ and $b_\tau$ be the three-vertex flags in which the new
vertex is adjacent to exactly the first, respectively exactly the second, root, and let
$g_\tau$ be the common-neighbour triangle flag.  Then on $S_\tau(Y_{P_4})$,
\[
  a_\tau=b_\tau,
  \qquad
  a_\tau+b_\tau=2g_\tau .
\]
\end{theorem}

\begin{proof}
The verified $K_4$-free $P_4$ certificate has the form
\[
  \pi_{P_4}
  +\frac89\,\avg{(2z_\eta-g_\eta)^2}{\eta}
  +5\,\avg{(a_\tau-b_\tau)^2}{\tau}
  +\frac{35}{9}\,\avg{(a_\tau+b_\tau-2g_\tau)^2}{\tau}
  \le \frac{32}{9}\,\one_0 .
\]
This is exactly the certificate encoded by
\texttt{K4freeP4.f\_1}, \texttt{K4freeP4.f\_2}, and \texttt{K4freeP4.f\_3}:
the generated flags
\[
  \texttt{FlagAlgebra\_3\_2\_0\_0},\quad
  \texttt{FlagAlgebra\_3\_2\_0\_3}
\]
are $z_\eta$ and $g_\eta$, while
\[
  \texttt{FlagAlgebra\_3\_2\_1\_1},\quad
  \texttt{FlagAlgebra\_3\_2\_1\_2},\quad
  \texttt{FlagAlgebra\_3\_2\_1\_3}
\]
are $a_\tau$, $b_\tau$, and $g_\tau$.  Applying
\cref{prop:equality-slice-vanishing} to the equality slice $Y_{P_4}$ gives the three
displayed labelled equations.
\end{proof}

\begin{remark}
\Cref{thm:k4free-p4-equality-slice} is the kind of example the relative method is meant
to produce.  It does not start from a known extremal graphon and read off its local
densities.  Instead, it takes a bound already proved by a flag-algebra certificate and
extracts necessary local laws for every extremal limit.  These laws are invisible to the
ordinary hereditary quotient: the quotient sees all $K_4$-free labelled limits, while the
relative support sees only roots that occur with positive probability inside
$P_4$-extremal limits.
\end{remark}

The same mechanism applies uniformly to the whole parametric family of certificates, and
there it also processes terms that are not averaged squares.

\begin{theorem}[Parametric $K_{r+1}$-free $P_4$ equality equations]
\label{thm:parametric-p4-equality-slice}
Fix an integer $r\ge3$; for $r\ge4$, assume Zykov's $K_4$-density \emph{bound}
\cite{Zykov1949}: every $K_{r+1}$-free graphon has induced-$K_4$ density at most
$\frac{(r-1)(r-2)(r-3)}{r^3}$ (at $r=3$ this is the vacuous bound $\phi_0(K_4)\le0$,
derived internally from $K_4$-freeness).  Let $\T_r$ be the $K_{r+1}$-free graph class, and
let $\pi_{P_4}^{(r)}\in\A^0[\T_r]$ be the $P_4$ density expression used in
\texttt{API.CompleteGraphFreeP4}.  Define the extremal slice
\[
  Y_r:=\left\{\phi_0\in Q_0^{(r)}:
    \phi_0(\pi_{P_4}^{(r)})=12\left(\frac{r-1}{r}\right)^3\right\}.
\]
Let $\eta$ be the ordered non-edge type and $\tau$ the ordered edge type.

For $\eta$, let $z_\eta$ be the three-vertex flag in which the new vertex is adjacent
to neither root, and let $g_\eta$ be the flag in which it is adjacent to both roots.
Then on $S_\eta(Y_r)$,
\[
  (r-1)z_\eta=g_\eta .
\]

For $\tau$, let $a_\tau$ and $b_\tau$ be the three-vertex flags in which the new
vertex is adjacent to exactly the first, respectively exactly the second, root, and let
$g_\tau$ be the common-neighbour flag.  Then on $S_\tau(Y_r)$,
\[
  a_\tau=b_\tau,
  \qquad
  (r-2)(a_\tau+b_\tau)=2g_\tau .
\]

Moreover, every $\phi_0\in Y_r$ has extremal $K_4$ density:
\[
  \phi_0(K_4)=\frac{(r-1)(r-2)(r-3)}{r^3},
\]
where $K_4$ denotes the unlabelled four-clique density element.
\end{theorem}

\begin{proof}
The verified certificate in \texttt{API.CompleteGraphFreeP4} has the exact form
\[
\begin{aligned}
  \pi_{P_4}^{(r)}
  &+p_1(r)\avg{((r-1)z_\eta-g_\eta)^2}{\eta}
   +p_2(r)\avg{(a_\tau-b_\tau)^2}{\tau} \\
  &+p_3(r)\avg{((r-2)a_\tau+(r-2)b_\tau-2g_\tau)^2}{\tau}
   +p_0(r)\kappa_4^{(r)}
   +L_r
   =
   12\left(\frac{r-1}{r}\right)^3\one_0 .
\end{aligned}
\]
Here $p_i(r)>0$ for $i=0,1,2,3$, the term
\[
  \kappa_4^{(r)}
  :=
  \frac{(r-1)(r-2)(r-3)}{r^3}\one_0-K_4
\]
is non-negative on $K_{r+1}$-free limits by the Zykov $K_4$-density bound
\cite{Zykov1949}, and
$L_r$ is the non-negative leftover term in the Lean file.

This is an instance of \cref{thm:relative-slackness}, with $Y:=Y_r$,
$h:=\pi_{P_4}^{(r)}$, the three averaged squares as the terms
$\lambda_i\avg{f_i}{\sigma_i}$ (so hypothesis \ref{cs-hyp-supports} is automatic), and
\[
  n:=p_0(r)\,\kappa_4^{(r)}+L_r,
\]
which satisfies hypothesis \ref{cs-hyp-slack} because both parts are non-negative on
$Q_0^{(r)}\supseteq Y_r$.  The displayed identity gives hypothesis
\ref{cs-hyp-certificate} on all of $Q_0^{(r)}$, and every $\phi_0\in Y_r$ attains the
bound by the definition of the slice.  Conclusion \ref{cs-global} makes the three squares
vanish identically on the corresponding relative supports, which---taking square roots,
as in \cref{prop:equality-slice-vanishing}---is exactly the displayed labelled equation
for $\eta$ and the two for $\tau$.  Conclusion \ref{cs-exact} gives $\phi_0(n)=0$; both
parts of $n$ are non-negative, so $\phi_0(\kappa_4^{(r)})=0$, which is precisely the
stated extremal $K_4$ density.
\end{proof}

\begin{remark}[The parametric slice is non-empty]\label{rem:parametric-nonempty}
For every $r\ge3$ the slice $Y_r$ is non-empty: the balanced complete $r$-partite limit
attains $\pi_{P_4}^{(r)}=12\bigl(\frac{r-1}r\bigr)^3$.  This achievability direction is
formalised, for every $r\ge3$, as
\texttt{Kr\_plus\_1\_free\_P4\_density\_achievable} in
\texttt{API.CompleteGraphFreeP4}; \cref{thm:k4free-p4-tripartite} below uses its $r=3$
instance.
\end{remark}

\subsection{Moment consequences and the extremal graphon}

The labelled slice equations integrate against a representing graphon to exact moment
identities, and we work these out for the whole parametric family at once.  For a
graphon $W$ write
\[
  d(x):=\int W(x,u)\,du,\qquad
  c(x,y):=\int W(x,u)W(y,u)\,du,
\]
\[
  p:=\int d(x)\,dx,\qquad
  D:=\int d(x)^2\,dx,\qquad
  T:=\int W(x,y)c(x,y)\,dx\,dy,
\]
so that $p$ is the edge density and $T$ the triangle density in the usual graphon
normalisation, and $D-p^2=\int(d(x)-p)^2\,dx\ge0$ is the variance of the degree
function.  The translation between three-vertex rooted flags and these kernels is: at an
ordered non-edge root $(x,y)$,
\[
  z_\eta=1-d(x)-d(y)+c(x,y),\qquad g_\eta=c(x,y),
\]
while at an ordered edge root $(x,y)$,
\[
  a_\tau=d(x)-c(x,y),\qquad b_\tau=d(y)-c(x,y),\qquad g_\tau=c(x,y).
\]
An identity between rooted flags valid on a relative support translates, for a graphon
$W$ representing a point of the slice, into the corresponding kernel identity holding
almost everywhere with respect to the natural rooting measure: the extension measure at
$\tau$, respectively $\eta$, is the law of the view from a $W$-random ordered edge,
respectively ordered non-edge, so the identity holds for $W(x,y)$-almost every,
respectively $(1-W(x,y))$-almost every, ordered pair $(x,y)$.

\begin{theorem}[Moment identities from the local equations]\label{thm:parametric-moments}
Fix $r\ge3$ and set
\[
  \alpha_r^-:=\frac{r-1}{2r-3},
  \qquad
  \alpha_r^+:=\frac{r-1}{r},
\]
so that $\alpha_r^-\le\alpha_r^+$ for all $r\ge3$, with equality exactly at $r=3$, where
both equal $\frac23$.  Let $W$ be \emph{any} graphon whose kernels satisfy the two local
equations
\[
  (r-2)\bigl(d(x)+d(y)-2c(x,y)\bigr)=2c(x,y)
  \qquad\text{for $W$-almost every }(x,y),
\]
\[
  (r-1)\bigl(1-d(x)-d(y)+c(x,y)\bigr)=c(x,y)
  \qquad\text{for $(1-W)$-almost every }(x,y).
\]
By \cref{thm:parametric-p4-equality-slice} and the translation above, every graphon
representing a point of $Y_r$ is of this kind---and in addition satisfies
$d(x)=d(y)$ for $W$-almost every pair, an equation the moment identities below do not
need; at $r=3$, by \cref{thm:k4free-p4-equality-slice}, so is every graphon
representing a point of $Y_{P_4}$.  Then:
\begin{enumerate}[label=\textup{(\roman*)}, leftmargin=2.5em]
  \item\label{pm:T} $(r-1)\,T=(r-2)\,D$;
  \item\label{pm:D} $r(2r-3)\,D=(r-1)^2(3p-1)$;
  \item\label{pm:var} the degree variance is an explicit function of the edge density:
  \[
    \int(d(x)-p)^2\,dx
    =
    D-p^2
    =
    (\alpha_r^+-p)(p-\alpha_r^-);
  \]
  \item\label{pm:interval} consequently $\alpha_r^-\le p\le\alpha_r^+$, and $W$ is
  degree-regular ($d=p$ almost everywhere) precisely when
  $p\in\{\alpha_r^-,\alpha_r^+\}$.
\end{enumerate}
\end{theorem}

\begin{proof}
Integrating the edge equation against $W(x,y)$ and the non-edge equation against
$1-W(x,y)$, using
\[
  \int W(x,y)\,d(x)\,dx\,dy=\int d(x)^2\,dx=D,
  \qquad
  \int c(x,y)\,dx\,dy=\int d(u)^2\,du=D,
\]
gives
\[
  (r-2)(2D-2T)=2T,
  \qquad
  (r-1)\bigl((1-p)-(2p-2D)+(D-T)\bigr)=D-T .
\]
The first display is \ref{pm:T}.  Substituting $T=\frac{r-2}{r-1}D$ into the second and
collecting terms yields
\[
  (r-1)(1-3p)+(2r-1)D=\frac{D}{r-1},
  \qquad\text{i.e.}\qquad
  \bigl[(2r-1)(r-1)-1\bigr]D=(r-1)^2(3p-1),
\]
and $(2r-1)(r-1)-1=2r^2-3r=r(2r-3)$, which is \ref{pm:D}.

For \ref{pm:var}, \ref{pm:D} gives
\[
  D-p^2
  =
  \frac{(r-1)^2(3p-1)-r(2r-3)\,p^2}{r(2r-3)} .
\]
The numerator is $-r(2r-3)(p-\alpha_r^+)(p-\alpha_r^-)$: the monic quadratic
$p^2-\frac{3(r-1)^2}{r(2r-3)}p+\frac{(r-1)^2}{r(2r-3)}$ has root product
$\frac{(r-1)^2}{r(2r-3)}=\alpha_r^+\alpha_r^-$ and root sum
$\frac{3(r-1)^2}{r(2r-3)}=\alpha_r^++\alpha_r^-$, as one checks directly.  Hence
$D-p^2=(\alpha_r^+-p)(p-\alpha_r^-)$, and the left-hand side is the degree variance
$\int(d-p)^2$.

For \ref{pm:interval}: the variance is non-negative, and the product
$(\alpha_r^+-p)(p-\alpha_r^-)$ is non-negative exactly for
$\alpha_r^-\le p\le\alpha_r^+$; it vanishes exactly at the two endpoints, and the
variance vanishes exactly when $d=p$ almost everywhere.
\end{proof}

Note that \cref{thm:parametric-moments} uses neither $K_{r+1}$-freeness nor the
certificate: it is a statement about any graphon obeying the two local equations.
Remarkably, identifying the extremiser needs no more.  At the regular endpoint the two
equations are \emph{rigid} by themselves---no forbidden subgraph, no certificate, and no
classical equality theorem enters.

\begin{theorem}[Rigidity at the regular endpoint]\label{thm:slice-rigidity}
Fix $r\ge3$ and let $W$ be any graphon satisfying the two local equations of
\cref{thm:parametric-moments}.  If the edge density of $W$ is
$p=\alpha_r^+=\frac{r-1}r$, then $W$ is, up to a measure-preserving relabelling, the
balanced complete $r$-partite graphon $T_r$.
\end{theorem}

\begin{proof}
\emph{Step 1: the local profile.}
By \cref{thm:parametric-moments}\ref{pm:var}, $p=\alpha_r^+$ forces
$\int(d-p)^2=0$, so $d(x)=\frac{r-1}r$ for almost every $x$.  Substituting into the
non-edge equation: for $(1-W)$-almost every $(x,y)$,
\[
  (r-1)\Bigl(1-\tfrac{2(r-1)}r+c(x,y)\Bigr)=c(x,y),
  \qquad\text{i.e.}\qquad
  (r-2)\,c(x,y)=\frac{(r-1)(r-2)}{r},
\]
so $c(x,y)=\frac{r-1}r$ for $(1-W)$-almost every $(x,y)$ (here $r-2\ge1$).

\emph{Step 2: the equality case of $c\le d$.}
For such a pair,
\[
  0=d(y)-c(x,y)=\int W(y,u)\bigl(1-W(x,u)\bigr)\,du,
\]
so $W(x,u)=1$ for almost every $u$ with $W(y,u)>0$, and symmetrically with $x$ and $y$
exchanged.  Writing $\{W(x,\cdot)>0\}$ for the support of the section, this gives the
chain of inclusions, all modulo null sets,
\[
  \{W(y,\cdot)>0\}\subseteq\{W(x,\cdot)=1\}\subseteq\{W(x,\cdot)>0\}
  \subseteq\{W(y,\cdot)=1\}\subseteq\{W(y,\cdot)>0\},
\]
so all five sets coincide up to null sets.  In particular $W(x,\cdot)$ is almost
everywhere $\{0,1\}$-valued, equal to the indicator of a set
$N(x):=\{W(x,\cdot)>0\}$ of measure $d(x)=\frac{r-1}r$, and $N(x)=N(y)$.  Since
$\int(1-W(x,y))\,dy=1-d(x)=\frac1r>0$ for almost every $x$, Fubini's theorem upgrades
this to: for almost every $x$, $W(x,\cdot)=\one_{N(x)}$ almost everywhere with
$|N(x)|=\frac{r-1}r$, and $N(x)=N(y)$ for almost every $y$ with $W(x,y)<1$.

\emph{Step 3: the partition.}
Set $C(x):=[0,1]\setminus N(x)$, so $|C(x)|=\frac1r$ and, up to a null set,
$C(x)=\{y:W(x,y)=0\}$.  Step 2 reads: for almost every $x$ and almost every
$y\in C(x)$, $C(y)=C(x)$.  Two consequences, for $x,x'$ in a suitable full-measure
set.  First, any two of the sets $C(x)$ are equal or disjoint modulo null sets: if
$|C(x)\cap C(x')|>0$, then a generic $y$ in the intersection satisfies both
$C(y)=C(x)$ and $C(y)=C(x')$.  Second, almost every $x$ lies in its own set: choosing
a generic $y\in C(x)$, symmetry of $W$ gives $W(y,x)=W(x,y)=0$, so
$x\in C(y)=C(x)$.  The distinct sets among the $C(x)$ are therefore pairwise disjoint
modulo null sets, each of measure exactly $\frac1r$, and they cover almost every point
of $[0,1]$; hence there are exactly $r$ of them, say $P_1,\dots,P_r$, and for almost
every $x\in P_i$ one has $C(x)=P_i$.

Thus, for almost every pair $(x,y)$: $W(x,y)=0$ when $x$ and $y$ lie in the same part
and $W(x,y)=1$ otherwise.  A measure-preserving relabelling sending each $P_i$ to an
interval of length $\frac1r$ carries $W$ to the balanced complete $r$-partite graphon.
\end{proof}

\begin{corollary}[Unconditional rigidity at $r=3$]\label{cor:r3-rigidity}
Any graphon satisfying the two local equations at $r=3$---that is,
$d(x)+d(y)=4c(x,y)$ for $W$-almost every pair and
$2\bigl(1-d(x)-d(y)+c(x,y)\bigr)=c(x,y)$ for $(1-W)$-almost every pair---is, up to a
measure-preserving relabelling, the balanced complete tripartite graphon $T_3$.
\end{corollary}

\begin{proof}
At $r=3$ the roots coincide, $\alpha_3^-=\alpha_3^+=\frac23$, so
\cref{thm:parametric-moments}\ref{pm:var} reads
$0\le\int(d-p)^2=-(p-\frac23)^2\le0$; hence $p=\frac23=\alpha_3^+$, and
\cref{thm:slice-rigidity} applies.
\end{proof}

\begin{theorem}[A certificate route to the tripartite extremiser]\label{thm:k4free-p4-tripartite}
Let $W$ be a graphon representing a point of the equality slice $Y_{P_4}$, and suppose the
ordered edge type $\tau$ and the ordered non-edge type $\eta$ both occur with positive
probability under random rooting from $W$ (this holds automatically at every point of
$Y_{P_4}$, since the graphon identified below has edge density $2/3$ and non-edge density
$1/3$ --- a routine consequence of the identification, not separately verified here; we
record it as an explicit hypothesis because passing from a point of $Y_{P_4}$ to some
representing graphon $W$ already invokes the existence half of the Lov\'asz--Szegedy
representation theorem).  Then $W$ is, up to a measure-preserving relabelling, the
balanced complete tripartite graphon; in particular every extremiser has edge density
$2/3$ and is degree-regular.  The identification itself uses no further classical equality
theorem: the two labelled slice equations are rigid by themselves
(\cref{cor:r3-rigidity}).  The slice is moreover non-empty---the balanced complete
tripartite graphon attains $\pi_{P_4}=32/9$---so $Y_{P_4}$ consists of exactly that
graphon.
\end{theorem}

\begin{proof}
By \cref{thm:k4free-p4-equality-slice}, the
kernels of $W$ satisfy the two local equations of \cref{thm:parametric-moments} at
$r=3$: the equation $a_\tau+b_\tau=2g_\tau$ gives the edge line and $2z_\eta=g_\eta$
gives the non-edge line (the third slice equation, $a_\tau=b_\tau$, is not needed).
\Cref{cor:r3-rigidity} identifies $W$, up to measure-preserving relabelling, with the
balanced complete tripartite graphon; the edge density $2/3$ and the degree-regularity
are read off from it.  Note that even the $K_4$-freeness of $W$ played no role in the
identification.

Conversely, the balanced complete tripartite graphon is itself $K_4$-free and attains
$\pi_{P_4}=32/9$: this is the matching lower bound for the $P_4$ density, the achievability
direction formalised in \texttt{API.CompleteGraphFreeP4} at $r=3$
(\cref{rem:parametric-nonempty}).  It therefore lies in
$Y_{P_4}$, so the slice is non-empty and equals exactly this graphon.
\end{proof}

\begin{remark}
An earlier version of \cref{thm:k4free-p4-tripartite} pinned the edge density from the
slice equations and then invoked the equality case of Tur\'an's theorem to identify the
extremiser.  \Cref{cor:r3-rigidity} shows the detour is unnecessary: the relative
ensemble enhancement converts the vanishing of labelled SOS terms into labelled local
equations, and \emph{those two equations alone characterise the extremal graphon}.  The
$P_4$ certificate is thus a first instance of a slice that is not controlled by a
pre-existing extremal equality theorem but by its own mined equations---the situation
open problem (3) below asks to produce systematically.
\end{remark}

\begin{corollary}[Qualitative stability]\label{cor:k4free-p4-qualitative-stability}
Let $(W_n)$ be a sequence of $K_4$-free graphons such that
\[
  \phi_{W_n}(\pi_{P_4})\longrightarrow \frac{32}{9}.
\]
Then $W_n$ converges in cut distance, up to measure-preserving relabelling, to the
balanced complete tripartite graphon $T_3$.

Equivalently, if $(G_n)$ is a sequence of finite $K_4$-free graphs with
$\pi_{P_4}(G_n)\to32/9$, then $G_n$ converges to $T_3$ in the graphon sense.
\end{corollary}

\begin{proof}
By \cref{thm:k4free-p4-tripartite}, the slice $Y_{P_4}$ is the single point represented
by $T_3$.  The space $Q_0$ of $K_4$-free unlabelled limits is closed in $X_0$, and the
$P_4$ density expression evaluates continuously, so
\cref{prop:unique-slice-stability}---with $Z:=Q_0$ and the single constraint
$\phi(\pi_{P_4})=32/9$---gives $\phi_{W_n}\to\phi_{T_3}$ in the density topology.  Under
the identification of unlabelled limits with graphons up to measure-preserving
relabelling, in which the density topology coincides with the cut metric, this is the
assertion; the finite-graph statement follows by applying it to the graphons of the
graphs $G_n$.
\end{proof}

\begin{corollary}[Recovery modulo Zykov equality]
\label{cor:parametric-p4-turan-recovery}
Assume $r\ge3$ and the equality case of Zykov's $K_4$-density theorem
\cite{Zykov1949}: among $K_{r+1}$-free graphons, equality in the $K_4$-density bound
occurs only at the balanced complete $r$-partite graphon.  Then every element of $Y_r$
is the balanced complete $r$-partite graphon.  (At $r=3$ the Zykov-equality hypothesis
is not actually needed, by \cref{cor:r3-rigidity}; the statement is recorded for
$r\ge3$ uniformly because it remains true, and the $r=3$ case costs nothing extra.)

Consequently, the relative supports at the one-vertex type, the ordered edge type
$\tau$, and the ordered non-edge type $\eta$ are all single points.  More explicitly, if
$e$ is the one-root edge flag, then on $S_{\vtype}(Y_r)$,
\[
  e=\frac{r-1}{r}\one_{\vtype} .
\]
At the edge type $\tau$,
\[
  a_\tau=b_\tau=\frac1r\one_\tau,
  \qquad
  g_\tau=\frac{r-2}{r}\one_\tau,
  \qquad
  z_\tau=0 .
\]
At the non-edge type $\eta$,
\[
  z_\eta=\frac1r\one_\eta,
  \qquad
  g_\eta=\frac{r-1}{r}\one_\eta,
  \qquad
  a_\eta=b_\eta=0 .
\]
\end{corollary}

\begin{proof}
By \cref{thm:parametric-p4-equality-slice}, every $\phi_0\in Y_r$ has the extremal
$K_4$ density among $K_{r+1}$-free graphons.  The assumed equality case of Zykov's
theorem therefore identifies $\phi_0$ with the balanced complete $r$-partite graphon.
The displayed labelled identities are the readings of random labelled extensions of that
graphon computed in the proof of \cref{thm:turan-slice}.
\end{proof}

The rigidity theorem localises what the Zykov assumption is actually for: not the
structure of the extremiser, only its edge density.

\begin{corollary}[Unconditional recovery at the regular endpoint]\label{cor:top-endpoint-recovery}
Fix $r\ge4$, let $W$ be a graphon representing a point $\phi_0\in Y_r$ with edge density
$\phi_0(\rho)=\frac{r-1}r$, and suppose the ordered edge type $\tau$ and the ordered
non-edge type $\eta$ both occur with positive probability under random rooting from $W$
(automatic at this edge density, exactly as in \cref{thm:k4free-p4-tripartite}; recorded
as an explicit hypothesis for the same reason: identifying a representing $W$ already
invokes Lov\'asz--Szegedy existence).  Then $W$ is, up to a measure-preserving
relabelling, the balanced complete $r$-partite graphon---with no appeal to the Zykov
\emph{equality} case (the Zykov \emph{bound} underlying
\cref{thm:parametric-p4-equality-slice} is still assumed, as throughout).  Consequently,
the hypothesis of \cref{cor:parametric-p4-turan-recovery} can be weakened from the full
equality case of Zykov's theorem to the single scalar statement that every element of
$Y_r$ has edge density $\frac{r-1}r$.
\end{corollary}

\begin{proof}
A graphon representing such a $\phi_0$ satisfies the two local equations by
\cref{thm:parametric-p4-equality-slice}, and its edge density is $\alpha_r^+$, so
\cref{thm:slice-rigidity} identifies it with $T_r$; the labelled identities of
\cref{cor:parametric-p4-turan-recovery} then follow as before.
\end{proof}

\begin{corollary}[Parametric qualitative stability]\label{cor:parametric-qualitative-stability}
Fix $r\ge4$ and assume the equality case of Zykov's theorem, as in
\cref{cor:parametric-p4-turan-recovery}.  If $(W_k)$ is a sequence of $K_{r+1}$-free
graphons with
\[
  \phi_{W_k}(\pi_{P_4}^{(r)})\longrightarrow12\left(\frac{r-1}{r}\right)^3,
\]
then $W_k$ converges in cut distance, up to measure-preserving relabelling, to the
balanced complete $r$-partite graphon $T_r$.
\end{corollary}

\begin{proof}
By \cref{cor:parametric-p4-turan-recovery}, every element of $Y_r$ is the limit
represented by $T_r$, and by \cref{rem:parametric-nonempty} the slice is non-empty; so
$Y_r$ is the single point represented by $T_r$, and
\cref{prop:unique-slice-stability} applies exactly as in
\cref{cor:k4free-p4-qualitative-stability}.
\end{proof}

\begin{remark}[Why $r=3$ is special]\label{rem:r3-special}
For $r\ge4$ the two moment identities alone do not pin the edge density: by
\cref{thm:parametric-moments} they force the exact trade-off between edge density and
degree variance recorded in \ref{pm:var}, confining $p$ to the interval
$[\alpha_r^-,\alpha_r^+]$ whose endpoints are the only degree-regular options, and no
interior value of $p$ is excluded by the moment identities themselves.  The case $r=3$
is exactly the degenerate case $\alpha_r^-=\alpha_r^+$ of a double root, and it is this
coincidence---not any additional strength of the certificate---that pins $p$
unconditionally and hands the argument to the rigidity theorem.  For $r\ge4$ the
structure question is settled by rigidity at the regular endpoint
(\cref{thm:slice-rigidity,cor:top-endpoint-recovery}); what remains classical is the
single scalar step of pinning $p$ to that endpoint, for which the certificate's
$\kappa_4^{(r)}$ term---pinning the $K_4$ density at its Zykov-extremal value
(\cref{thm:parametric-p4-equality-slice})---delegates to the Zykov equality case
(\cref{cor:parametric-p4-turan-recovery,cor:parametric-qualitative-stability}).
Whether that last step can also be internalised---for instance by playing the pinned
$K_4$ density against the moment identities---is open problem (4) below.  The equality
analysis behind \cref{thm:slice-rigidity} genuinely fails at the other endpoint: at
$p=\alpha_r^-$ the non-edge codegree equals $\frac{d}{r-2}<d$, so no equality case of
$c\le d$ is available there.
\end{remark}

\subsection{Quantitative stability}

Near the slice, the exact moment identities survive with explicit errors.  The next
theorem is certificate-free: it is pure Cauchy--Schwarz and holds for \emph{every}
graphon; certificates enter only afterwards, to bound the two square errors by the
certificate deficit.

\begin{theorem}[Approximate moment identities]\label{thm:approximate-moments}
Fix $r\ge3$ and let $W$ be any graphon.  Define the local errors
\[
  \ell_\eta(x,y):=(r-1)\bigl(1-d(x)-d(y)+c(x,y)\bigr)-c(x,y),
  \qquad
  \ell_\tau(x,y):=(r-2)\bigl(d(x)+d(y)-2c(x,y)\bigr)-2c(x,y),
\]
and their square averages
\[
  R_\eta:=\int(1-W(x,y))\,\ell_\eta(x,y)^2\,dx\,dy,
  \qquad
  R_\tau:=\int W(x,y)\,\ell_\tau(x,y)^2\,dx\,dy .
\]
Then, with $\alpha_r^\pm$ as in \cref{thm:parametric-moments},
\[
  \Bigl|
  r(2r-3)\bigl(D-p^2\bigr)
  -
  r(2r-3)\,(\alpha_r^+-p)(p-\alpha_r^-)
  \Bigr|
  \le
  (r-1)\sqrt{R_\eta}+\frac{r-2}{2}\sqrt{R_\tau}.
\]
In particular, since $D-p^2\ge0$,
\[
  r(2r-3)\,(p-\alpha_r^-)(p-\alpha_r^+)
  \le
  (r-1)\sqrt{R_\eta}+\frac{r-2}{2}\sqrt{R_\tau},
\]
and
\[
  r(2r-3)\int(d(x)-p)^2\,dx
  \le
  r(2r-3)\,(\alpha_r^+-p)(p-\alpha_r^-)
  +
  (r-1)\sqrt{R_\eta}+\frac{r-2}{2}\sqrt{R_\tau}.
\]
\end{theorem}

\begin{proof}
Define the integrated local errors
\[
  E_\eta:=\int(1-W(x,y))\,\ell_\eta(x,y)\,dx\,dy,
  \qquad
  E_\tau:=\int W(x,y)\,\ell_\tau(x,y)\,dx\,dy .
\]
Expanding as in the proof of \cref{thm:parametric-moments},
\[
  E_\tau=2(r-2)D-2(r-1)T,
  \qquad
  E_\eta=(r-1)(1-3p)+(3r-4)D-(r-2)T .
\]
Eliminating $T$,
\[
  (r-1)E_\eta-\frac{r-2}{2}E_\tau
  =
  (r-1)^2(1-3p)+\bigl[(r-1)(3r-4)-(r-2)^2\bigr]D
  =
  (r-1)^2(1-3p)+r(2r-3)D,
\]
since $(r-1)(3r-4)-(r-2)^2=2r^2-3r=r(2r-3)$.  Subtracting $r(2r-3)p^2$ from
$r(2r-3)D=(r-1)^2(3p-1)+(r-1)E_\eta-\frac{r-2}{2}E_\tau$ and using the factorisation
$(r-1)^2(3p-1)-r(2r-3)p^2=r(2r-3)(\alpha_r^+-p)(p-\alpha_r^-)$ from the proof of
\cref{thm:parametric-moments},
\[
  r(2r-3)\bigl(D-p^2\bigr)
  =
  r(2r-3)\,(\alpha_r^+-p)(p-\alpha_r^-)
  +(r-1)E_\eta-\frac{r-2}{2}E_\tau .
\]
By the Cauchy--Schwarz inequality,
\[
  |E_\eta|
  \le
  \Bigl(\int(1-W)\Bigr)^{1/2}
  \Bigl(\int(1-W)\,\ell_\eta^2\Bigr)^{1/2}
  \le\sqrt{R_\eta},
  \qquad
  |E_\tau|\le\sqrt{R_\tau},
\]
and the first display follows.  The second display drops the non-negative term
$r(2r-3)(D-p^2)$ from the right-hand side of the identity (after moving the product term
to the left), and the third drops nothing but the sign of the error.
\end{proof}

The certificate now supplies the bounds on $R_\eta$ and $R_\tau$.  We first record the
$r=3$ instance in the notation of the verified certificate \texttt{API.K4freeP4}, then
the parametric family.

\begin{proposition}[Near equality forces approximate local equations]
\label{prop:k4free-p4-certificate-stability}
Let $W$ be a $K_4$-free graphon, and put
\[
  \Delta:=\frac{32}{9}-\phi_W(\pi_{P_4})\ge0 .
\]
As before, write
\[
  d(x):=\int W(x,u)\,du,\qquad
  c(x,y):=\int W(x,u)W(y,u)\,du .
\]
Define the three local square errors
\[
\begin{aligned}
  R_\eta(W)
    &:=\int (1-W(x,y))
      \bigl(2(1-d(x)-d(y)+c(x,y))-c(x,y)\bigr)^2\,dx\,dy,\\
  R_\tau^{-}(W)
    &:=\int W(x,y)\bigl(d(x)-d(y)\bigr)^2\,dx\,dy,\\
  R_\tau^{+}(W)
    &:=\int W(x,y)\bigl(d(x)+d(y)-4c(x,y)\bigr)^2\,dx\,dy .
\end{aligned}
\]
Then
\[
  R_\eta(W)\le\frac98\,\Delta,\qquad
  R_\tau^{-}(W)\le\frac15\,\Delta,\qquad
  R_\tau^{+}(W)\le\frac9{35}\,\Delta .
\]
\end{proposition}

\begin{proof}
The three integrals are the graphon evaluations of the three averaged square terms
\[
  \avg{(2z_\eta-g_\eta)^2}{\eta},\qquad
  \avg{(a_\tau-b_\tau)^2}{\tau},\qquad
  \avg{(a_\tau+b_\tau-2g_\tau)^2}{\tau},
\]
by the translation between rooted three-vertex flags and the kernels $d,c$ recorded
before \cref{thm:parametric-moments}.  Evaluating the verified certificate
\[
  \pi_{P_4}
  +\frac89\,\avg{(2z_\eta-g_\eta)^2}{\eta}
  +5\,\avg{(a_\tau-b_\tau)^2}{\tau}
  +\frac{35}{9}\,\avg{(a_\tau+b_\tau-2g_\tau)^2}{\tau}
  \le \frac{32}{9}\one_0
\]
at $W$ gives
\[
  \frac89 R_\eta(W)+5R_\tau^{-}(W)+\frac{35}{9}R_\tau^{+}(W)\le\Delta ,
\]
which is \cref{thm:relative-slackness}\ref{cs-approx} for this certificate.  Each
summand is non-negative, so the three displayed estimates follow.
\end{proof}

\begin{theorem}[Certificate-to-Tur\'an stability]\label{thm:k4free-p4-quant-stability}
Let $W$ be a $K_4$-free graphon, let
\[
  \Delta:=\frac{32}{9}-\phi_W(\pi_{P_4}),
\]
and let $p:=\int W(x,y)\,dx\,dy$ be its edge density.  Then
\[
  (3p-2)^2
  \le
  \left(\frac{3}{\sqrt2}+\frac{3}{2\sqrt{35}}\right)\sqrt{\Delta},
\]
and hence
\[
  \left|p-\frac23\right|
  \le
  \frac13
  \left(\frac{3}{\sqrt2}+\frac{3}{2\sqrt{35}}\right)^{1/2}
  \Delta^{1/4}.
\]
Moreover the degree function concentrates: with $d(x):=\int W(x,u)\,du$,
\[
  \int\bigl(d(x)-p\bigr)^2\,dx
  \le
  \frac19\left(\frac{3}{\sqrt2}+\frac{3}{2\sqrt{35}}\right)\sqrt{\Delta}.
\]

Finally, let $\omega_{\mathrm{Tur}}$ be any graphon Tur\'an-stability modulus for
$K_4$-free graphons: for every $\gamma>0$, if a $K_4$-free graphon $U$ has edge
density at least $2/3-\omega_{\mathrm{Tur}}(\gamma)$, then
$\delta_\square(U,T_3)<\gamma$.  Such a modulus exists by compactness and the equality
case of Tur\'an's theorem.  Then
\[
  \Delta
  \le
  \left(
    \frac{3\,\omega_{\mathrm{Tur}}(\gamma)}
    {\left(\frac{3}{\sqrt2}+\frac{3}{2\sqrt{35}}\right)^{1/2}}
  \right)^4
  \quad\Longrightarrow\quad
  \delta_\square(W,T_3)<\gamma .
\]
\end{theorem}

\begin{proof}
At $r=3$ the local errors of \cref{thm:approximate-moments} are
$\ell_\eta=2(1-d(x)-d(y)+c)-c$ and $\ell_\tau=d(x)+d(y)-4c$, so that
$R_\eta=R_\eta(W)$ and $R_\tau=R_\tau^{+}(W)$ in the notation of
\cref{prop:k4free-p4-certificate-stability}; moreover $r(2r-3)=9$ and
$\alpha_3^\pm=\frac23$, so $9(\alpha_3^+-p)(p-\alpha_3^-)=-(3p-2)^2$.  The first display
of \cref{thm:approximate-moments} therefore reads
\[
  \Bigl|\,9\bigl(D-p^2\bigr)+(3p-2)^2\,\Bigr|
  \le
  2\sqrt{R_\eta(W)}+\frac12\sqrt{R_\tau^{+}(W)} .
\]
Both $9(D-p^2)$ and $(3p-2)^2$ are non-negative, so each is at most the right-hand side,
which by \cref{prop:k4free-p4-certificate-stability} is at most
\[
  2\left(\frac98\Delta\right)^{1/2}
  +\frac12\left(\frac9{35}\Delta\right)^{1/2}
  =
  \left(\frac{3}{\sqrt2}+\frac{3}{2\sqrt{35}}\right)\sqrt{\Delta}.
\]
This gives the first assertion and, after dividing by $9$, the degree-concentration
estimate; the displayed edge-density estimate follows by taking square roots and
dividing by $3$.

The final implication is just Tur\'an stability applied after this edge-density estimate:
the displayed upper bound on $\Delta$ gives
$p\ge2/3-\omega_{\mathrm{Tur}}(\gamma)$, hence
$\delta_\square(W,T_3)<\gamma$.
\end{proof}

\begin{theorem}[Parametric quantitative stability]\label{thm:parametric-quant-stability}
Fix $r\ge3$, let $W$ be a $K_{r+1}$-free graphon, and let
\[
  \Delta:=12\left(\frac{r-1}{r}\right)^3-\phi_W(\pi_{P_4}^{(r)}),
\]
which is non-negative by \cref{thm:relative-slackness}\ref{cs-soundness} applied to the
parametric certificate of \cref{thm:parametric-p4-equality-slice}.  Let
$p_0(r),\dots,p_3(r)>0$ be the coefficients of that certificate and let
$R_\eta,R_\tau$ be as in \cref{thm:approximate-moments}.  Then:
\begin{enumerate}[label=\textup{(\roman*)}, leftmargin=2.5em]
  \item\label{pq:squares} $R_\eta\le\Delta/p_1(r)$, $R_\tau\le\Delta/p_3(r)$, and
  $\displaystyle\int W(x,y)\bigl(d(x)-d(y)\bigr)^2\,dx\,dy\le\Delta/p_2(r)$;
  \item\label{pq:k4} the $K_4$ density of $W$ is nearly extremal:
  \[
    \phi_W(K_4)\ge\frac{(r-1)(r-2)(r-3)}{r^3}-\frac{\Delta}{p_0(r)}\,;
  \]
  \item\label{pq:quadratic} with
  $C_r:=\dfrac{r-1}{\sqrt{p_1(r)}}+\dfrac{r-2}{2\sqrt{p_3(r)}}$,
  \[
    r(2r-3)\,(p-\alpha_r^-)(p-\alpha_r^+)
    \le
    C_r\sqrt\Delta
    \qquad\text{and}\qquad
    r(2r-3)\Bigl|\,\bigl(D-p^2\bigr)-(\alpha_r^+-p)(p-\alpha_r^-)\Bigr|
    \le
    C_r\sqrt\Delta ;
  \]
  in particular $p$ lies within $C_r\sqrt\Delta/\bigl((r-1)(r-3)\bigr)$ of the interval
  $[\alpha_r^-,\alpha_r^+]$ when $r\ge4$, and
  $|p-\tfrac23|\le\tfrac13C_3^{1/2}\Delta^{1/4}$ when $r=3$;
  \item\label{pq:zykov} if $r\ge4$, assume the equality case of Zykov's theorem as in
  \cref{cor:parametric-p4-turan-recovery}, and let $\omega_{\mathrm{Zyk}}$ be any
  corresponding stability modulus: for every $\gamma>0$, if a $K_{r+1}$-free graphon $U$
  has $\phi_U(K_4)\ge\frac{(r-1)(r-2)(r-3)}{r^3}-\omega_{\mathrm{Zyk}}(\gamma)$, then
  $\delta_\square(U,T_r)<\gamma$.  (Such a modulus exists by compactness and the assumed
  equality case, exactly as for $\omega_{\mathrm{Tur}}$ in
  \cref{thm:k4free-p4-quant-stability}.)  Then
  \[
    \Delta\le p_0(r)\,\omega_{\mathrm{Zyk}}(\gamma)
    \quad\Longrightarrow\quad
    \delta_\square(W,T_r)<\gamma .
  \]
\end{enumerate}
\end{theorem}

\begin{proof}
The parametric certificate is an identity whose left-hand side evaluates at $\phi_W$ to
$\phi_W(\pi_{P_4}^{(r)})$ plus non-negative terms, so
\cref{thm:relative-slackness}\ref{cs-approx}---with $Y=Q_0^{(r)}$, the three squares as
the $f_i$, and $n=p_0(r)\kappa_4^{(r)}+L_r$ as in the proof of
\cref{thm:parametric-p4-equality-slice}---bounds their weighted sum by $\Delta$.  The
three squares evaluate at $W$ to the three integrals of \ref{pq:squares}, by the same
translation as in \cref{prop:k4free-p4-certificate-stability}; this gives
\ref{pq:squares}, and the bound $\phi_W(n)\le\Delta$ gives
$p_0(r)\,\phi_W(\kappa_4^{(r)})\le\Delta$, which is \ref{pq:k4}.

\Cref{thm:approximate-moments} combined with \ref{pq:squares} gives \ref{pq:quadratic}:
both displayed inequalities follow from the identity in that proof once
$(r-1)\sqrt{R_\eta}+\frac{r-2}2\sqrt{R_\tau}\le C_r\sqrt\Delta$.  For $r\ge4$ the
quadratic $q(p):=r(2r-3)(p-\alpha_r^-)(p-\alpha_r^+)$ has simple roots at distance
$\alpha_r^+-\alpha_r^-=\frac{(r-1)(r-3)}{r(2r-3)}>0$; if, say, $p>\alpha_r^+$, then
\[
  q(p)\ge r(2r-3)\,(p-\alpha_r^+)(\alpha_r^+-\alpha_r^-)
  =(r-1)(r-3)\,(p-\alpha_r^+),
\]
so $q(p)\le C_r\sqrt\Delta$ places $p$ within $C_r\sqrt\Delta/\bigl((r-1)(r-3)\bigr)$
of the interval, and symmetrically below $\alpha_r^-$.  At $r=3$ the roots coincide and
the same inequality reads $(3p-2)^2\le C_3\sqrt\Delta$, giving the $\Delta^{1/4}$ rate.

For \ref{pq:zykov}, \ref{pq:k4} gives
$\phi_W(K_4)\ge\frac{(r-1)(r-2)(r-3)}{r^3}-\omega_{\mathrm{Zyk}}(\gamma)$, and the
modulus applies.
\end{proof}

\begin{remark}[On the exponents]
The certificate controls the squared local errors by $O(\Delta)$, hence the integrated
local equations by $O(\sqrt\Delta)$, and the moment argument spends one more square root
exactly when the quadratic in $p$ has a double root.  Thus the direct certificate route
gives the fourth-root modulus of \cref{thm:k4free-p4-quant-stability} at $r=3$, where
$\alpha_3^-=\alpha_3^+$; for $r\ge4$ the same inequality confines $p$ to an
$O_r(\sqrt\Delta)$-neighbourhood of the interval $[\alpha_r^-,\alpha_r^+]$---a
square-root rate, but only interval localisation.  What pins the extremiser for
$r\ge4$ is instead the $\kappa_4^{(r)}$ term, whose control in
\cref{thm:parametric-quant-stability}\ref{pq:k4} is \emph{linear} in $\Delta$; the only
non-explicit ingredient of \cref{thm:parametric-quant-stability}\ref{pq:zykov} is the
modulus of the Zykov equality case.  A sharper modulus at $r=3$ would require additional
information beyond the three square terms.
\end{remark}

\section{Open problems}

Four questions remain.

\begin{enumerate}[label=\textup{(\arabic*)}, leftmargin=2.8em]
  \item \textbf{Beyond sparse root repairs.}
  \Cref{thm:blowup-root-plantable} unifies pure false-clone blow-ups, pure true-clone
  blow-ups, and full substitution-closure
  (\cref{thm:clone-root-plantable,thm:true-clone-root-plantable,thm:substitution-root-plantable})
  under blow-up-closure; \cref{thm:finite-local-planting} gives the abstract finite criterion; and
  \cref{thm:sparse-repair-planting} verifies it for sparse old-edge repairs after
  root blow-up.  The next useful theorem would allow denser but structured repairs,
  repairs incident to the planted root clusters, or amalgamation hypotheses.

  \item \textbf{Exact certificate gaps before closure.}
  \Cref{thm:no-closed-certificate-gap} rules out any improvement of asymptotic density
  bounds after taking the natural closed cone.  What remains is an exact finite-degree
  question: can $\C^{\mathrm{ens}}_\sigma$ be strictly larger than
  $\C^{\mathrm{quot}}_\sigma$ before closure in a way that changes a finite SDP or a
  rational certificate?  The edge type in the $C_5$-free class is the first natural test:
  it is non-root-plantable, its pinned witness and all multiples unlabel to zero, but its
  support is not a single point.

  \item \textbf{Equality-slice mining beyond classical equality theorems.}
  \Cref{thm:k4free-p4-tripartite,thm:parametric-p4-equality-slice} show that a verified
  density certificate can be read backwards as a system of labelled local equations on its
  equality slice, and \cref{cor:r3-rigidity,thm:slice-rigidity} show that in the
  $K_{r+1}$-free $P_4$ family these equations are rigid by themselves at $r=3$, and at
  the regular endpoint for every $r$; the only remaining classical ingredient is the
  endpoint pinning for $r\ge4$ (\cref{cor:top-endpoint-recovery}).  The next test is to
  mine a certificate for a problem whose extremal structure is genuinely unknown, and
  decide whether the kernel equations of \cref{thm:kernel-slackness} are rigid there
  too.

  \item \textbf{Internalising the endpoint, and relative plantability.}
  Two questions raised by the completeness theory of
  \cref{subsec:slice-completeness}.  First, for $r\ge4$, can the edge density of a
  point of $Y_r$ be pinned to $\frac{r-1}r$ without the Zykov equality case---for
  instance by playing the pinned $K_4$ density of
  \cref{thm:parametric-p4-equality-slice} against the moment identities of
  \cref{thm:parametric-moments}, or against further labelled equations mined at larger
  types?  By \cref{cor:top-endpoint-recovery} this scalar step is all that separates
  the parametric family from a fully certificate-internal identification.  Second,
  \cref{prop:mantel-not-plantable} shows a slice can fail relative root-plantability
  while the ambient class satisfies the absolute property; characterise the relatively
  plantable pairs $(Y,\sigma)$, and decide whether the gap
  $Q_\sigma(Y)\setminus S_\sigma(Y)$ can ever carry an \emph{exact} certificate
  difference, in the sense of question (2), that survives over the slice seminorm.
\end{enumerate}

\begin{thebibliography}{9}

\bibitem{Razborov2007}
A.~A.~Razborov.
\newblock Flag algebras.
\newblock \emph{Journal of Symbolic Logic}, 72(4):1239--1282, 2007.

\bibitem{Zykov1949}
A.~A.~Zykov.
\newblock On some properties of linear complexes.
\newblock \emph{Mat. Sb. (N.S.)}, 24(66), no.~2:163--188, 1949.
\newblock \url{https://www.mathnet.ru/eng/sm5974}.

\bibitem{ErdosStone}
P.~Erd\H{o}s and A.~H.~Stone.
\newblock On the structure of linear graphs.
\newblock \emph{Bulletin of the American Mathematical Society}, 52:1087--1091, 1946.

\bibitem{ErdosSimonovits}
P.~Erd\H{o}s and M.~Simonovits.
\newblock A limit theorem in graph theory.
\newblock \emph{Studia Scientiarum Mathematicarum Hungarica}, 1:51--57, 1966.

\bibitem{Simonovits1968}
M.~Simonovits.
\newblock A method for solving extremal problems in graph theory, stability problems.
\newblock In \emph{Theory of Graphs (Proc. Colloq., Tihany, 1966)}, pages 279--319.
\newblock Academic Press, New York, 1968.

\bibitem{KovariSosTuran}
T.~K\H{o}v\'ari, V.~T.~S\'os, and P.~Tur\'an.
\newblock On a problem of K. Zarankiewicz.
\newblock \emph{Colloquium Mathematicae}, 3:50--57, 1954.

\bibitem{BondySimonovits}
J.~A.~Bondy and M.~Simonovits.
\newblock Cycles of even length in graphs.
\newblock \emph{Journal of Combinatorial Theory, Series B}, 16(2):97--105, 1974.

\bibitem{LovaszSzegedy2012}
L.~Lov\'asz and B.~Szegedy.
\newblock Random graphons and a weak Positivstellensatz for graph algebras.
\newblock \emph{Journal of Graph Theory}, 70(2):214--225, 2012.

\end{thebibliography}

\end{document}
